% Établissement de la chronologie des évènements géologiques d'une région — SVT 1ère D — Cours signé OnBuch+
% Programme officiel MINESEC. Compiler avec : tectonic ND13-chronologie-evenements-geologiques.tex
\newcommand{\DOCMATIERE}{SVT}
\newcommand{\DOCNIVEAU}{1ère D}
\newcommand{\DOCMODULE}{Module 3 : L'Éducation à l'Environnement et au Développement Durable}
\newcommand{\DOCLECON}{ND13}
\newcommand{\DOCTITRE}{Établissement de la chronologie des évènements géologiques d'une région}
% OnBuch+ — préambule « cours signés » ENRICHI (compilé avec tectonic / XeLaTeX)
% Système visuel « Pop » d'OnBuch+ : fond crème (jamais blanc), bordures encre
% épaisses, ombres « dures » décalées, titres Archivo Black en capitales.
% Version MATHÉMATIQUES : graphes (pgfplots/tikz), boîtes pédagogiques
% (définition, propriété, méthode, attention, le savais-tu, exercice résolu,
% exercice, TP, bilan), page de garde et signature OnBuch+ ancrée en bas.
%
% Macros attendues dans main.tex AVANT \input{preamble} :
%   \DOCMATIERE  \DOCNIVEAU  \DOCMODULE  \DOCLECON (numéro)  \DOCTITRE
\documentclass[11pt]{article}

\usepackage{fontspec}
\usepackage[french]{babel}
\usepackage[a4paper,margin=2cm,top=3.4cm,bottom=2.8cm,headheight=1.4cm,footskip=1.3cm]{geometry}
\usepackage{amsmath,amssymb,amsfonts}
\usepackage{mathtools} % \xmapsto, \coloneqq... souvent utilisés par les modèles
\usepackage{cancel} % simplifications barrées (bilans de forces)
% \abs{z} (module, valeur absolue) et \norm{u} (norme), utilisés par les modèles
\providecommand{\abs}{}\let\abs\relax\DeclarePairedDelimiter{\abs}{\lvert}{\rvert}
\providecommand{\norm}{}\let\norm\relax\DeclarePairedDelimiter{\norm}{\lVert}{\rVert}
\usepackage[table]{xcolor} % [table] : fournit aussi \rowcolors (alternance de lignes)
\usepackage{pagecolor}
\usepackage{tikz}
\usetikzlibrary{arrows.meta,positioning,calc,shapes.geometric,shapes.misc,shapes.arrows,decorations.pathmorphing,decorations.pathreplacing,decorations.markings,patterns,shadows,mindmap,trees,fit,backgrounds,matrix,intersections,angles,quotes}
\usepackage{pgfplots}
\usepackage[version=4]{mhchem} % équations nucléaires \ce{^{235}_{92}U}
\usepackage[european,siunitx]{circuitikz} % schémas électriques
\pgfplotsset{compat=1.18}
\usepgfplotslibrary{fillbetween,groupplots} % aires entre courbes (calcul intégral)
\usepackage{enumitem}
\usepackage{fancyhdr}
% [most] charge toutes les bibliothèques tcolorbox (listings, minted, raster,
% documentation...) et gonfle inutilement la mémoire TeX sur un document long
% (~300 boîtes) : on ne charge que ce qui est réellement utilisé.
\usepackage[skins,breakable]{tcolorbox}
\usepackage{graphicx}
\usepackage{colortbl}
\usepackage{booktabs}
\usepackage{tabularx}
\usepackage{diagbox} % case d'en-tête diagonale des tableaux à double entrée
% Colonnes extensibles centrée / gauche / droite (C, L, R), souvent utilisées par les modèles
\newcolumntype{C}{>{\centering\arraybackslash}X}
\newcolumntype{L}{>{\raggedright\arraybackslash}X}
\newcolumntype{R}{>{\raggedleft\arraybackslash}X}
\usepackage{array}
\usepackage{multicol}
\usepackage{siunitx}
% parse-numbers=false : accepte aussi \SI{2,5 \times 10^{-3}}{...}
\sisetup{locale=FR,output-decimal-marker={,},group-minimum-digits=4,parse-numbers=false}
\DeclareSIUnit\ohm{\ensuremath{\symup{\Omega}}} % U+2126 absent de la police
\DeclareSIUnit\division{div} % réglages d'oscilloscope (ms/div, V/div)
\DeclareSIUnit\tr{tr} % tours (vitesse de rotation, ex. \SI{1500}{\tr\per\minute})
\DeclareSIUnit\ch{ch} % cheval-vapeur (puissance moteur, unité usuelle hors SI)
\DeclareSIUnit\franccfa{FCFA} % franc CFA (coûts, contexte camerounais)
\DeclareSIUnit\dioptre{\ensuremath{\delta}} % vergence d'une lentille (dioptrie)
\usepackage{qrcode}
% \degree : commande fréquemment utilisée par les modèles pour un angle ou une
% température en dehors de \SI{}{} (non fournie par les paquets chargés).
\providecommand{\degree}{^{\circ}}
% \qed (fin de démonstration), utilisé par les modèles sans amsthm
\providecommand{\qed}{\hfill$\square$}
% \barre : double barre des valeurs interdites dans un tableau de variations (tkz-tab)
\providecommand{\barre}{\ensuremath{\|}}
% environnement proof (amsthm non chargé) utilisé par les modèles
\ifdefined\proof\else\newenvironment{proof}[1][Démonstration]{\par\noindent\textit{#1.}\ }{\qed\par}\fi
\usepackage{lastpage}
\usepackage{hyperref}
\hypersetup{hidelinks}

% ============================================================
% Palette « Pop » OnBuch+ — mêmes hex que la classe P (plus_ui.dart)
% ============================================================
\definecolor{poporange}{HTML}{F59321}
\definecolor{popdark}{HTML}{E07A0C}
\definecolor{popcream}{HTML}{FFF6E8}
\definecolor{popink}{HTML}{1C1714}
\definecolor{poppurple}{HTML}{7A5AE0}
\definecolor{popgreen}{HTML}{2E9E6A}
\definecolor{poppink}{HTML}{E0457B}
\definecolor{popblue}{HTML}{2F7DE1}
\definecolor{popgold}{HTML}{C98A12}
\definecolor{popmuted}{HTML}{8A7F76}
\definecolor{popline}{HTML}{EADFD2}
\definecolor{popgray}{HTML}{8A7F76}
\colorlet{popgrey}{popgray}
% Alias tolérants (le modèle nomme parfois les couleurs différemment)
\colorlet{popwhite}{white}
\colorlet{popblack}{popink}
\colorlet{popred}{poppink}
\colorlet{popyellow}{popgold}
% Teintes claires pour fonds de tableaux / zones
\colorlet{popblueL}{popblue!10!white}
\colorlet{poporangeL}{poporange!14!white}
\colorlet{popgreenL}{popgreen!12!white}
\colorlet{poppurpleL}{poppurple!12!white}
\colorlet{poppinkL}{poppink!12!white}
\colorlet{popgoldL}{popgold!14!white}

\pagecolor{popcream}

% ============================================================
% Typographie
% ============================================================
\defaultfontfeatures{Ligatures=TeX}
\setmainfont{PlusJakartaSans-Medium.ttf}[Path=../../../fonts/,BoldFont=PlusJakartaSans-Bold.ttf]
% Maths en sans-serif (Fira Math) avec chiffres et lettres droites de Plus
% Jakarta, pour que les formules se fondent dans le texte.
\usepackage[math-style=ISO,bold-style=ISO]{unicode-math}
\setmathfont{FiraMath-Regular.otf}[Path=../../../fonts/]
\setmathfont{PlusJakartaSans-Medium.ttf}[Path=../../../fonts/,range={up/num,up/{Latin,latin},\mathup}]
\usepackage{icomma} % virgule décimale sans espace en mode math
% Caractères mathématiques Unicode absents des polices de texte (Plus Jakarta,
% Archivo Black) : rendus par leur équivalent mathématique, en texte comme en
% formule — sinon ils s'affichent en carré vide (ex. « Arithmétique dans ℤ »).
\usepackage{newunicodechar}
\newunicodechar{ℝ}{\ensuremath{\mathbb{R}}}
\newunicodechar{ℤ}{\ensuremath{\mathbb{Z}}}
\newunicodechar{ℂ}{\ensuremath{\mathbb{C}}}
\newunicodechar{ℕ}{\ensuremath{\mathbb{N}}}
\newunicodechar{ℚ}{\ensuremath{\mathbb{Q}}}
\newunicodechar{𝔻}{\ensuremath{\mathbb{D}}}
\newunicodechar{α}{\ensuremath{\alpha}}
\newunicodechar{β}{\ensuremath{\beta}}
\newunicodechar{γ}{\ensuremath{\gamma}}
\newunicodechar{δ}{\ensuremath{\delta}}
\newunicodechar{ε}{\ensuremath{\varepsilon}}
\newunicodechar{θ}{\ensuremath{\theta}}
\newunicodechar{λ}{\ensuremath{\lambda}}
\newunicodechar{μ}{\ensuremath{\mu}}
\newunicodechar{σ}{\ensuremath{\sigma}}
\newunicodechar{τ}{\ensuremath{\tau}}
\newunicodechar{φ}{\ensuremath{\varphi}}
\newunicodechar{ψ}{\ensuremath{\psi}}
\newunicodechar{ω}{\ensuremath{\omega}}
\newunicodechar{Σ}{\ensuremath{\Sigma}}
% majuscules produites par \MakeUppercase dans les titres (α-aminés -> Α)
\newunicodechar{Α}{\ensuremath{\alpha}}
\newunicodechar{Β}{\ensuremath{\beta}}
\newunicodechar{Γ}{\ensuremath{\Gamma}}
\newunicodechar{Δ}{\ensuremath{\Delta}}
\newunicodechar{Ε}{\ensuremath{\varepsilon}}
\newunicodechar{Θ}{\ensuremath{\Theta}}
\newunicodechar{Λ}{\ensuremath{\Lambda}}
\newunicodechar{Μ}{\ensuremath{\mu}}
\newunicodechar{Τ}{\ensuremath{\tau}}
\newunicodechar{Φ}{\ensuremath{\Phi}}
\newunicodechar{Ψ}{\ensuremath{\Psi}}
\newunicodechar{Ω}{\ensuremath{\Omega}}
\newunicodechar{↦}{\ensuremath{\mapsto}}
\newunicodechar{∈}{\ensuremath{\in}}
\newunicodechar{∉}{\ensuremath{\notin}}
\newunicodechar{∧}{\ensuremath{\wedge}}
\newunicodechar{⇔}{\ensuremath{\Leftrightarrow}}
\newunicodechar{⟺}{\ensuremath{\Longleftrightarrow}}
\newunicodechar{⇒}{\ensuremath{\Rightarrow}}
\newunicodechar{⟹}{\ensuremath{\Longrightarrow}}
\newunicodechar{∘}{\ensuremath{\circ}}
\newunicodechar{∪}{\ensuremath{\cup}}
\newunicodechar{∩}{\ensuremath{\cap}}
\newunicodechar{≡}{\ensuremath{\equiv}}
\newunicodechar{⊂}{\ensuremath{\subset}}
\newunicodechar{⊥}{\ensuremath{\perp}}
\newunicodechar{∃}{\ensuremath{\exists}}
\newunicodechar{∀}{\ensuremath{\forall}}
\newunicodechar{∛}{\ensuremath{\sqrt[3]{\,}}}
\newunicodechar{∜}{\ensuremath{\sqrt[4]{\,}}}
\newunicodechar{⃗}{\ensuremath{{}^{\to}}}
\newunicodechar{ⁿ}{\ifmmode^{n}\else\textsuperscript{\ensuremath{n}}\fi}
\newunicodechar{ˣ}{\ifmmode^{x}\else\textsuperscript{\ensuremath{x}}\fi}
\newunicodechar{ᵇ}{\ifmmode^{b}\else\textsuperscript{\ensuremath{b}}\fi}
\newunicodechar{ʳ}{\ifmmode^{r}\else\textsuperscript{\ensuremath{r}}\fi}
\newunicodechar{ᵘ}{\ifmmode^{u}\else\textsuperscript{\ensuremath{u}}\fi}
\newunicodechar{ʸ}{\ifmmode^{y}\else\textsuperscript{\ensuremath{y}}\fi}
\newunicodechar{ᵃ}{\ifmmode^{a}\else\textsuperscript{\ensuremath{a}}\fi}
\newunicodechar{ᵉ}{\ifmmode^{e}\else\textsuperscript{\ensuremath{e}}\fi}
\newunicodechar{ᵏ}{\ifmmode^{k}\else\textsuperscript{\ensuremath{k}}\fi}
\newunicodechar{ᵖ}{\ifmmode^{p}\else\textsuperscript{\ensuremath{p}}\fi}
\newunicodechar{ᴺ}{\ifmmode^{N}\else\textsuperscript{\ensuremath{N}}\fi}
\newunicodechar{ᶜ}{\ifmmode^{c}\else\textsuperscript{\ensuremath{c}}\fi}
\newunicodechar{ʰ}{\ifmmode^{h}\else\textsuperscript{\ensuremath{h}}\fi}
\newunicodechar{ᵠ}{\ifmmode^{q}\else\textsuperscript{\ensuremath{q}}\fi}
\newunicodechar{ₙ}{\ifmmode_{n}\else\textsubscript{\ensuremath{n}}\fi}
\newunicodechar{ₐ}{\ifmmode_{a}\else\textsubscript{\ensuremath{a}}\fi}
\newunicodechar{ᵢ}{\ifmmode_{i}\else\textsubscript{\ensuremath{i}}\fi}
\newunicodechar{ᵣ}{\ifmmode_{r}\else\textsubscript{\ensuremath{r}}\fi}
\newunicodechar{ₖ}{\ifmmode_{k}\else\textsubscript{\ensuremath{k}}\fi}
\newunicodechar{ᵦ}{\ifmmode_{\beta}\else\textsubscript{\ensuremath{\beta}}\fi}
\newunicodechar{ₓ}{\ifmmode_{x}\else\textsubscript{\ensuremath{x}}\fi}
\newfontfamily\popdisplay{ArchivoBlack-Regular.ttf}[Path=../../../fonts/]
\newfontfamily\poplogo{SpaceGrotesk-Bold.ttf}[Path=../../../fonts/]
\newfontfamily\popsemi{PlusJakartaSans-SemiBold.ttf}[Path=../../../fonts/]
\newfontfamily\popxbold{PlusJakartaSans-ExtraBold.ttf}[Path=../../../fonts/]
\newfontfamily\popxboldls{PlusJakartaSans-ExtraBold.ttf}[Path=../../../fonts/,LetterSpace=3.0]

\setlist[enumerate]{leftmargin=1.7em,itemsep=5pt,topsep=4pt}
\setlist[itemize]{leftmargin=1.7em,itemsep=4pt,topsep=4pt,label={\color{poporange}\textbullet}}
\setlength{\parindent}{0pt}
\setlength{\parskip}{5pt}
\renewcommand{\arraystretch}{1.25}

% pgfplots : style Pop par défaut
\pgfplotsset{
  every axis/.append style={axis line style={popink,line width=1pt},tick style={popink},
    grid style={popline,line width=0.5pt},label style={font=\small},tick label style={font=\footnotesize}},
  popcurve/.style={poporange,line width=1.8pt,no markers,smooth},
}
% Même style disponible dans un \draw TikZ simple (hors pgfplots)
\tikzset{popcurve/.style={poporange,line width=1.8pt}}
\tikzset{popgrid/.style={popline,very thin}}
\colorlet{popcurve}{poporange} % le nom du style est parfois utilisé comme couleur

% ============================================================
% Pill / squiggle
% ============================================================
\newcommand{\poppill}[3]{%
  \tikz[baseline=-0.55ex]\node[draw=popink,line width=1.2pt,rounded corners=2.6pt,%
    fill=#2,inner xsep=6.5pt,inner ysep=3pt]{%
    \popxboldls\fontsize{7}{7}\selectfont\color{#3}\MakeUppercase{#1}};%
}
\newcommand{\popsquiggle}[1]{%
  \tikz[baseline=-0.4ex]{
    \draw[#1,line width=2.6pt,line cap=round]
      plot [smooth] coordinates {(0,0.09) (0.34,0.01) (0.68,0.21) (1.02,0.01) (1.36,0.10)};
  }%
}
% Mot-clé surligné (marqueur orange sous le texte)
\newcommand{\cle}[1]{{\popxbold\color{popdark}#1}}

% ============================================================
% En-tête / pied
% ============================================================
\pagestyle{fancy}
\fancyhf{}
\renewcommand{\headrulewidth}{0pt}
\renewcommand{\footrulewidth}{0pt}
\lhead{\raisebox{-0.15ex}{\poplogo\fontsize{13}{13}\selectfont\color{popink}OnBuch\color{poporange}+}}
\rhead{\poppill{\DOCMATIERE\ \textbullet\ \DOCNIVEAU\ \textbullet\ Leçon \DOCLECON}{white}{popink}}
\lfoot{\popxbold\fontsize{7.6}{7.6}\selectfont\color{popmuted}\MakeUppercase{Cours signé OnBuch+}\ \textbullet\ {\popsemi\DOCTITRE}}
\rfoot{\tikz[baseline=-0.6ex]\node[rounded corners=8.5pt,fill=popink,minimum height=17pt,minimum width=17pt,inner xsep=5pt,inner sep=0]{\popxbold\fontsize{8}{8}\selectfont\color{popcream}\thepage\,/\,\pageref*{LastPage}};}

% ============================================================
% Titres
% ============================================================
\newcommand{\chaptitre}[2]{%
  \begin{tcolorbox}[enhanced,colback=poporange,colframe=popink,boxrule=2.6pt,arc=11pt,
      drop shadow=popink,left=15pt,right=15pt,top=12pt,bottom=11pt,before skip=3pt,after skip=18pt]
    \poppill{#2}{popink}{popcream}\par\vspace{9pt}
    {\raggedright\hyphenpenalty=10000\exhyphenpenalty=10000\popdisplay\fontsize{22}{24}\selectfont\color{popcream}\MakeUppercase{#1}\par}
  \end{tcolorbox}
}
% Section numérotée : pastille orange avec le numéro + titre Archivo
\newcounter{popsec}
\newcommand{\coursec}[1]{%
  \stepcounter{popsec}\setcounter{popsub}{0}%
  \par\vspace{16pt}\noindent
  \begin{minipage}{\linewidth}
  \tikz[baseline=-0.7ex]\node[draw=popink,line width=1.6pt,fill=poporange,rounded corners=4pt,minimum size=22pt,inner sep=0]{\popdisplay\fontsize{12}{12}\selectfont\color{popcream}\thepopsec};%
  \hspace{8pt}{\popdisplay\fontsize{15}{17}\selectfont\color{popink}\MakeUppercase{#1}}\par
  \vspace{4pt}\noindent{\color{poporange}\rule{36pt}{3.6pt}}
  \end{minipage}\par\nopagebreak\vspace{8pt}
}
\newcounter{popsub}
\newcommand{\courssub}[1]{%
  \stepcounter{popsub}%
  \par\vspace{9pt}\noindent{\popxbold\fontsize{11.5}{13}\selectfont\color{popink}{\color{poporange}\thepopsec.\thepopsub}\hspace{6pt}#1}\par\nopagebreak\vspace{4pt}
}

% ============================================================
% Boîtes pédagogiques — même recette : fond blanc, bordure épaisse colorée,
% ombre dure encre, badge plein. Titre optionnel [#1].
% ============================================================
\tcbset{popbase/.style={breakable,enhanced,colback=white,boxrule=2.3pt,arc=9pt,
  shadow={3pt}{-3pt}{0pt}{popink},left=14pt,right=14pt,top=11pt,bottom=11pt,before skip=11pt,after skip=15pt,
  fontupper=\popsemi\fontsize{10.6}{14.5}\selectfont\color{popink},
  fontlower=\popsemi\fontsize{10.6}{14.5}\selectfont\color{popink}}}
% Coche dessinée (le glyphe ✓ n'existe pas dans les polices du document)
\newcommand{\popcheck}[1]{\tikz[baseline=-0.5ex]\draw[#1,line width=1.6pt,line cap=round,line join=round](-0.13,0.0)--(-0.04,-0.09)--(0.13,0.1);}
\providecommand{\coche}{\popcheck{popgreen}}
\providecommand{\resultat}[1]{\par\smallskip\noindent\fcolorbox{popgreen}{popgreenL}{\parbox{\dimexpr\linewidth-2\fboxsep-2\fboxrule\relax}{\textbf{Résultat.} #1}}\par\smallskip}
\newcommand{\popboxhead}[3]{\poppill{#1}{#2}{white}\ifx\relax#3\relax\else\hspace{6pt}{\popxbold\fontsize{10.6}{12}\selectfont\color{popink}#3}\fi\par\vspace{7pt}}

\newtcolorbox{aretenir}[1][]{popbase,colframe=popgreen,before upper={\popboxhead{À retenir}{popgreen}{#1}}}
\newtcolorbox{exemplebox}[1][]{popbase,colframe=poporange,before upper={\popboxhead{Exemple}{poporange}{#1}}}
\newtcolorbox{definition}[1][]{popbase,colframe=popblue,before upper={\popboxhead{Définition}{popblue}{#1}}}
\newtcolorbox{propriete}[1][]{popbase,colframe=poppurple,before upper={\popboxhead{Propriété}{poppurple}{#1}}}
\newtcolorbox{methode}[1][]{popbase,colframe=popgold,before upper={\popboxhead{Méthode}{popgold}{#1}}}
\newtcolorbox{attention}[1][]{popbase,colframe=poppink,before upper={\popboxhead{Attention}{poppink}{#1}}}
\newtcolorbox{experience}[1][]{popbase,colframe=popblue,colback=popblueL,before upper={\popboxhead{Expérience}{popblue}{#1}}}
% « Le savais-tu ? » : carte violette pleine (contexte, culture, Cameroun)
\newtcolorbox{savaistu}[1][]{popbase,colback=poppurple,colframe=popink,
  fontupper=\popsemi\fontsize{10.4}{14.2}\selectfont\color{white},
  before upper={\poppill{Le savais-tu\,?}{white}{poppurple}\ifx\relax#1\relax\else\hspace{6pt}{\popxbold\color{white}#1}\fi\par\vspace{7pt}}}
% Exercice résolu : énoncé en haut, \tcblower puis la solution sur fond crème
\newcounter{popexo}\newcounter{popexr}
\newtcolorbox{exoresolu}[1][]{popbase,colframe=poporange,colbacklower=popcream,
  segmentation style={popink,line width=1.2pt,dashed},
  before upper={\stepcounter{popexr}\popboxhead{Exercice résolu \thepopexr}{poporange}{#1}},
  before lower={\poppill{Solution}{popink}{popcream}\par\vspace{6pt}}}
% Exercice (non résolu en place) — la correction est dans la partie Corrigés
\newtcolorbox{exercice}[1][]{popbase,colframe=popink,
  before upper={\stepcounter{popexo}\popboxhead{Exercice \thepopexo}{popink}{#1}}}
% Corrigé
\newtcolorbox{corrige}[1][]{popbase,colframe=popgreen,colback=popgreenL,
  before upper={\popboxhead{Corrigé}{popgreen}{#1}}}
% Objectifs / prérequis / bilan
\newtcolorbox{objectifs}{popbase,colframe=popink,colback=poporangeL,
  before upper={\popboxhead{Objectifs de la leçon}{popink}{}}}
\newtcolorbox{prerequis}{popbase,colframe=popink,colback=popblueL,
  before upper={\popboxhead{Prérequis}{popblue}{}}}
\newtcolorbox{bilan}{popbase,colframe=popink,colback=white,boxrule=2.8pt,
  before upper={\popboxhead{Fiche bilan}{poporange}{L'essentiel à connaître pour l'examen}}}
% Figure encadrée avec légende
\newtcolorbox{popfigure}[1][]{enhanced,breakable=false,colback=white,colframe=popline,boxrule=1.4pt,arc=8pt,
  left=10pt,right=10pt,top=10pt,bottom=8pt,before skip=10pt,after skip=12pt,halign=center,
  lower separated=false,halign lower=center,
  fontlower=\popsemi\fontsize{9}{11.5}\selectfont\color{popmuted},
  #1}
\newcommand{\legende}[1]{\par\vspace{6pt}{\popsemi\fontsize{9}{11.5}\selectfont\color{popmuted}#1}}

% ============================================================
% Signature OnBuch+
% ============================================================
\newcommand{\poplogoinline}[1]{{\poplogo\fontsize{#1}{#1}\selectfont\color{popink}OnBuch\color{poporange}+}}
\newcommand{\signatureonbuch}{%
  \begin{tcolorbox}[enhanced,colback=popink,colframe=popink,boxrule=2.2pt,arc=10pt,
      drop shadow=poporange,left=14pt,right=14pt,top=10pt,bottom=10pt,before skip=0pt,after skip=0pt]
    \tikz[baseline=-0.7ex]\node[circle,fill=popgreen,draw=popcream,line width=1.3pt,minimum size=22pt,inner sep=0]{\popcheck{white}};%
    \hspace{9pt}\begin{minipage}[c]{0.62\linewidth}
      {\popxbold\fontsize{12}{13}\selectfont\color{popcream}Signé\ \textbullet\ {\poplogo OnBuch\color{poporange}+}}\par\vspace{2pt}
      {\raggedright\popsemi\fontsize{8}{10.5}\selectfont\color{popline}\MakeUppercase{Équipe pédagogique OnBuch+}\\\MakeUppercase{Conforme au programme officiel MINESEC}\par}
    \end{minipage}\hfill
    {\poplogo\fontsize{22}{22}\selectfont\color{popcream}OnBuch\color{poporange}+}
  \end{tcolorbox}%
}

% ============================================================
% Page de garde
% #1 titre · #2 matière · #3 niveau · #4 module · #5 n° leçon · #6 durée ·
% #7 description / objectifs (texte ou itemize)
% ============================================================
\newcommand{\pagegarde}[7]{%
  \thispagestyle{empty}
  \newgeometry{a4paper,margin=1.7cm,top=1.6cm,bottom=1.4cm}%
  \begin{tikzpicture}[remember picture,overlay]
    \fill[poporange!18] ([xshift=-3.2cm,yshift=-3.6cm]current page.north east) circle (4.6cm);
    \fill[poppurple!14] ([xshift=2.4cm,yshift=7.6cm]current page.south west) circle (3.8cm);
  \end{tikzpicture}%
  {\poplogo\fontsize{30}{30}\selectfont\color{popink}OnBuch\color{poporange}+}\hfill
  \raisebox{0.6ex}{\poppill{Cours signé \textbullet\ Premium}{popink}{popcream}}\par
  \vspace{1pt}\popsquiggle{poppurple}\par
  \vspace{8pt}
  \begin{tcolorbox}[enhanced,colback=poporange,colframe=popink,boxrule=2.8pt,arc=13pt,
      drop shadow=popink,left=18pt,right=18pt,top=12pt,bottom=13pt,after skip=10pt]
    \poppill{#2 \textbullet\ #3}{popink}{popcream}\hspace{5pt}\poppill{#4}{white}{popink}\par\vspace{8pt}
    {\popdisplay\fontsize{14}{14}\selectfont\color{popink}LEÇON #5}\par\vspace{4pt}
    {\raggedright\hyphenpenalty=10000\exhyphenpenalty=10000\popdisplay\fontsize{25}{27}\selectfont\color{popcream}\MakeUppercase{#1}\par}
    \vspace{8pt}
    {\popxbold\fontsize{9.5}{9.5}\selectfont\color{popink}\MakeUppercase{Durée conseillée : #6}}
  \end{tcolorbox}
  \begin{tcolorbox}[enhanced,colback=white,colframe=popink,boxrule=2.2pt,arc=10pt,
      drop shadow=popink,left=16pt,right=16pt,top=10pt,bottom=10pt,after skip=8pt]
    \poppill{Dans cette leçon}{poppurple}{white}\par\vspace{5pt}
    \popsemi\fontsize{9.3}{12.6}\selectfont\color{popink}#7
  \end{tcolorbox}
  \noindent\begin{minipage}[t]{0.487\linewidth}
    \begin{tcolorbox}[enhanced,colback=popgreen,colframe=popink,boxrule=2.2pt,arc=10pt,
        drop shadow=popink,left=11pt,right=11pt,top=10pt,bottom=10pt,height=3.1cm,valign=center]
      \tcbox[colback=white,colframe=popink,boxrule=1.3pt,arc=5pt,boxsep=0pt,nobeforeafter,
        left=3pt,right=3pt,top=3pt,bottom=3pt,valign=center]{\qrcode[height=1.9cm]{https://play.google.com/store/apps/details?id=com.luvvix.onbuch&hl=fr}}%
      \hspace{8pt}\begin{minipage}[c]{0.5\linewidth}
        {\popdisplay\fontsize{11}{12.5}\selectfont\color{white}\MakeUppercase{Télécharge OnBuch}}\par\vspace{4pt}
        {\raggedright\popsemi\fontsize{8.4}{11}\selectfont\color{white}Gratuit \textbullet\ Google Play\\Tuteur IA, annales, concours, résultats\par}
      \end{minipage}
    \end{tcolorbox}
  \end{minipage}\hfill
  \begin{minipage}[t]{0.487\linewidth}
    \begin{tcolorbox}[enhanced,colback=poppurple,colframe=popink,boxrule=2.2pt,arc=10pt,
        drop shadow=popink,left=11pt,right=11pt,top=10pt,bottom=10pt,height=3.1cm,valign=center]
      \tikz[baseline=-0.7ex]\node[circle,fill=white,draw=popink,line width=1.3pt,minimum size=1.6cm,inner sep=0]{\popdisplay\fontsize{24}{24}\selectfont\color{poppurple}+};%
      \hspace{8pt}\begin{minipage}[c]{0.6\linewidth}
        {\popdisplay\fontsize{11}{12.5}\selectfont\color{white}\MakeUppercase{Passe à OnBuch+}}\par\vspace{4pt}
        {\raggedright\popsemi\fontsize{8.4}{11}\selectfont\color{white}Tous les cours signés, Tuteur IA illimité, fiches PDF — onglet OnBuch+ de l'app\par}
      \end{minipage}
    \end{tcolorbox}
  \end{minipage}
  \vspace{8pt}
  \popcontenu
  \vfill
  \signatureonbuch
  \restoregeometry
  \newpage
}

% Rangée « Ce cours contient » (page de garde)
\newcommand{\poptile}[3]{% #1 couleur #2 gros texte #3 légende
  \begin{tcolorbox}[enhanced,colback=#1,colframe=popink,boxrule=2pt,arc=9pt,shadow={2.5pt}{-2.5pt}{0pt}{popink},
      left=6pt,right=6pt,top=9pt,bottom=9pt,halign=center,valign=center,height=1.75cm,width=\linewidth,nobeforeafter]
    {\popdisplay\fontsize{12.5}{13}\selectfont\color{white}\MakeUppercase{#2}}\par\vspace{3pt}
    {\centering\popsemi\fontsize{7.8}{9.5}\selectfont\color{white}#3\par}
  \end{tcolorbox}}
\newcommand{\popcontenu}{%
  \par\noindent{\popxboldls\fontsize{7.5}{7.5}\selectfont\color{popmuted}CE COURS SIGNÉ CONTIENT}\par\vspace{7pt}
  \noindent\begin{minipage}[t]{0.235\linewidth}\poptile{popblue}{Cours}{complet, illustré, pas à pas}\end{minipage}\hfill
  \begin{minipage}[t]{0.235\linewidth}\poptile{poporange}{Méthodes}{et exercices résolus}\end{minipage}\hfill
  \begin{minipage}[t]{0.235\linewidth}\poptile{poppink}{Exercices}{type examen + corrigés}\end{minipage}\hfill
  \begin{minipage}[t]{0.235\linewidth}\poptile{popgreen}{Fiche bilan}{l'essentiel à retenir}\end{minipage}\par}

% Sommaire
\newtcolorbox{sommaire}{enhanced,colback=white,colframe=popink,boxrule=2.2pt,arc=10pt,
  drop shadow=popink,left=18pt,right=18pt,top=13pt,bottom=13pt,before skip=6pt,after skip=20pt,
  before upper={\poppill{Sommaire}{poppurple}{white}\par\vspace{9pt}\popsemi\fontsize{10.6}{14.5}\selectfont\color{popink}}}

% Signature de fin de cours — poussée tout en bas de la dernière page
\newcommand{\findecours}{%
  \par\vfill\nopagebreak
  \begin{center}\popsquiggle{poporange}\end{center}\vspace{4pt}
  \signatureonbuch
}
\AtBeginDocument{\def\llbracket{\mathopen{[\mkern-3.5mu[}}\def\rrbracket{\mathclose{]\mkern-3.5mu]}}} % intervalles d'entiers (glyphe ⟦ absent de la police)
% Glyphes absents de Fira Math : redéfinis à partir de glyphes disponibles
\AtBeginDocument{%
  \def\setminus{\mathbin{\backslash}}\let\smallsetminus\setminus
  \def\vdots{\mathord{\vbox{\baselineskip3.2pt\lineskiplimit0pt\kern4pt\hbox{.}\hbox{.}\hbox{.}}}}%
  \def\ddots{\mathinner{\mkern1mu\raise6.5pt\hbox{.}\mkern2mu\raise3.6pt\hbox{.}\mkern2mu\raise0.7pt\hbox{.}\mkern1mu}}%
  \def\triangle{\mathord{\scalebox{0.9}{$\symup{\Delta}$}}}%
  \def\nabla{\mathord{\raisebox{\depth}{\scalebox{1}[-1]{$\symup{\Delta}$}}}}%
  \def\checkmark{\popcheck{popdark}}%
  \def\star{\mathbin{\ast}}\def\bigstar{\mathord{\ast}}%
  \def\diamond{\mathbin{\scalebox{0.6}{$\Diamond$}}}%
  \def\bigcup{\mathop{\mathchoice{\raisebox{-0.3em}{\scalebox{1.7}{$\cup$}}}{\scalebox{1.25}{$\cup$}}{\cup}{\cup}}\displaylimits}%
  \def\bigcap{\mathop{\mathchoice{\raisebox{-0.3em}{\scalebox{1.7}{$\cap$}}}{\scalebox{1.25}{$\cap$}}{\cap}{\cap}}\displaylimits}%
  \def\lesseqqgtr{\mathrel{\lessgtr}}\def\bigoplus{\mathop{\oplus}\displaylimits}%
}
\providecommand{\ssi}{si et seulement si} % abréviation fréquente
\AtBeginDocument{\providecommand{\wideparen}{\overparen}} % arc de cercle

%% FIGFIX-BEGIN v5 — correctifs globaux des figures (docs/onbuch-generation/tools/figfix)
\usepackage{adjustbox}\usepackage{etoolbox}\tcbuselibrary{raster}
% toute figure TikZ/pgfplots plus large que la ligne est réduite à la largeur disponible
\BeforeBeginEnvironment{tikzpicture}{\begin{adjustbox}{max width=\linewidth}}
\AfterEndEnvironment{tikzpicture}{\end{adjustbox}}
% légendes pgfplots lisibles par-dessus les courbes
\pgfplotsset{every axis/.append style={legend style={fill=white,fill opacity=0.92,text opacity=1,draw=black!15,inner sep=3pt}}}
% étiquettes d'arêtes (syntaxe edge["..."]) avec fond blanc
\tikzset{every edge quotes/.append style={fill=white,inner sep=1.2pt}}
%% FIGFIX-END
\begin{document}

\pagegarde{Établissement de la chronologie des évènements géologiques d'une région}{SVT}{1ère D}{Module 3 : L'Éducation à l'Environnement et au Développement Durable}{ND13}{20 h}{Cette leçon présente les principes de la stratigraphie (superposition, continuité, recoupement) et l'usage des fossiles stratigraphiques pour établir la chronologie relative des événements géologiques d'une région. À partir de coupes géologiques, l'élève apprend à ordonner dépôts, plis, failles et discordances pour reconstituer l'histoire géologique d'un terrain.
\vspace{4pt}
\begin{multicols}{2}\small
\begin{itemize}
\item À la fin de cette leçon, je sais définir la stratigraphie, une couche (strate) et une formation géologique.
\item À la fin de cette leçon, je sais énoncer et appliquer le principe de superposition pour ordonner des couches sédimentaires.
\item À la fin de cette leçon, je sais appliquer le principe de continuité latérale pour corréler des couches entre deux affleurements.
\item À la fin de cette leçon, je sais appliquer le principe de recoupement pour dater relativement une faille, un pli ou une intrusion.
\item À la fin de cette leçon, je sais reconnaître et interpréter une discordance stratigraphique sur une coupe.
\item À la fin de cette leçon, je sais définir un fossile stratigraphique et citer ses critères.
\end{itemize}\end{multicols}}

\begin{sommaire}
\begin{multicols}{2}
\begin{enumerate}
\item Stratigraphie et couches sédimentaires
\item Les principes de la stratigraphie
\item Les discordances stratigraphiques
\item Les fossiles stratigraphiques
\item Reconstitution de la chronologie des événements géologiques
\item Activité d'intégration
\item Méthodes et exercices résolus
\item Exercices
\item Corrigés des exercices
\item Fiche bilan
\end{enumerate}
\end{multicols}
\end{sommaire}

\begin{objectifs}
\begin{itemize}
\item À la fin de cette leçon, je sais définir la stratigraphie, une couche (strate) et une formation géologique.
\item À la fin de cette leçon, je sais énoncer et appliquer le principe de superposition pour ordonner des couches sédimentaires.
\item À la fin de cette leçon, je sais appliquer le principe de continuité latérale pour corréler des couches entre deux affleurements.
\item À la fin de cette leçon, je sais appliquer le principe de recoupement pour dater relativement une faille, un pli ou une intrusion.
\item À la fin de cette leçon, je sais reconnaître et interpréter une discordance stratigraphique sur une coupe.
\item À la fin de cette leçon, je sais définir un fossile stratigraphique et citer ses critères.
\item À la fin de cette leçon, je sais exploiter des données paléontologiques pour dater relativement une formation géologique.
\item À la fin de cette leçon, je sais établir la chronologie relative complète des événements géologiques d'une région à partir d'une coupe (dépôts, plissement, failles, érosion, discordances).
\end{itemize}
\end{objectifs}

\begin{prerequis}
\begin{itemize}
\item Notion de roche sédimentaire et de sédimentation (classes antérieures et début du Module 3)
\item Notion de fossile et de fossilisation
\item Notions de pli et de faille (déformations de la couverture sédimentaire)
\item Lecture et interprétation d'un schéma ou d'une coupe géologique simple
\item Grandes divisions de l'échelle des temps géologiques (ère primaire, secondaire, tertiaire, quaternaire)
\end{itemize}
\end{prerequis}

\newpage

\coursec{Stratigraphie et couches sédimentaires}

\textbf{Situation d'accroche.} Lors des travaux de terrassement de la route Yaoundé--Nsimalen, les engins ont mis au jour des couches de roches superposées, certaines plissées, d'autres coupées par des fractures, et l'on y a même trouvé des coquillages fossilisés... en pleine forêt continentale ! Comment des coquillages marins ont-ils pu se retrouver à des centaines de kilomètres des côtes actuelles ? Quelles couches se sont déposées en premier ? Quand la mer a-t-elle recouvert la région ? Quels événements (plissement, fracturation, érosion) se sont succédé, et dans quel ordre ?

\textbf{Problématique.} Comment les géologues peuvent-ils, à partir d'une simple coupe de terrain, reconstituer l'histoire géologique d'une région et ordonner les événements qui l'ont affectée, sans connaître leur âge exact en années ?

Pour répondre à cette question, il faut d'abord comprendre ce qu'est une couche sédimentaire, comment elle se forme, et ce que l'on peut lire en elle. C'est l'objet de cette première partie.

\courssub{La stratigraphie : la science des couches}

\textbf{Observation.} Sur le flanc d'une carrière, le long d'une tranchée routière ou sur une falaise littorale (par exemple à Kribi ou à Limbé), on observe souvent des roches disposées en bandes horizontales superposées, comme les feuilles d'un livre ou les étages d'un gâteau. Chaque bande diffère de ses voisines par sa couleur, sa dureté ou sa composition. Cette organisation n'est pas un hasard : elle résulte du dépôt successif de sédiments au cours du temps.

\begin{definition}[Stratigraphie]
La \cle{stratigraphie} (du latin \emph{stratum}, couche, et du grec \emph{graphein}, écrire) est la branche de la géologie qui étudie les \cle{couches} (ou \cle{strates}) des roches, principalement sédimentaires, leurs relations de superposition et leurs contenus, afin de reconstituer la chronologie des événements géologiques d'une région.
\end{definition}

\begin{definition}[Strate ou couche sédimentaire]
Une \cle{strate} (ou couche) est une unité de roche sédimentaire déposée dans des conditions relativement constantes, limitée par deux surfaces à peu près parallèles : la surface supérieure appelée \cle{toit} et la surface inférieure appelée \cle{mur}. Une strate est individualisée par ses caractères propres (composition, couleur, épaisseur, fossiles) qui la distinguent des couches voisines.
\end{definition}

\begin{attention}[Vocabulaire à ne pas confondre]
Ne confonds pas \emph{strate} et \emph{formation géologique}. Une \cle{formation géologique} est un ensemble de plusieurs strates présentant des caractères communs (même nature de roche, mêmes conditions de dépôt) et suffisamment important pour être cartographié. La strate est l'unité élémentaire ; la formation est un regroupement de strates. De même, le \cle{toit} est la limite supérieure de la couche (celle qui est en contact avec la couche plus récente), le \cle{mur} est sa limite inférieure (en contact avec la couche plus ancienne).
\end{attention}

\courssub{Affleurement et coupe géologique : observer les couches}

Pour étudier les couches, le géologue a besoin de les voir. Deux situations se présentent :

\begin{itemize}
\item Un \cle{affleurement} est une zone où les roches du sous-sol apparaissent directement à la surface du sol, sans couverture végétale ni sol épais : falaise, paroi de carrière, tranchée de route, berge de rivière. C'est le « terrain d'observation » du géologue.
\item Une \cle{coupe géologique} est une représentation schématique en section verticale de l'organisation des couches dans le sous-sol, dressée à partir des observations d'affleurements et, éventuellement, de sondages. Elle indique la nature, l'épaisseur, l'attitude (horizontale, plissée, inclinée) et les relations des couches.
\end{itemize}

\begin{exemplebox}[La tranchée de Yaoundé--Nsimalen]
Sur la tranchée de la route Yaoundé--Nsimalen, le géologue relève de bas en haut : une couche de sable grossier, une couche d'argile, une couche de calcaire contenant des coquillages, puis une couche de sable fin. Il reporte ces observations sur une coupe géologique en respectant l'ordre de superposition et l'épaisseur relative des couches. Cette coupe sera le document de base de toute l'étude chronologique.
\end{exemplebox}

\courssub{Comment se forme une couche sédimentaire ?}

\textbf{Observation et questionnement.} À l'embouchure de la Sanaga ou du Wouri, l'eau du fleuve, brune et chargée de particules, se jette dans la mer et devient progressivement claire : les particules transportées se déposent. D'où viennent ces particules, et que deviennent-elles une fois déposées ?

\begin{experience}[De la vase à la roche : démarche d'investigation]
\textbf{Observation :} dans un verre d'eau où l'on agite de la terre et du sable, on voit, après repos, le sable se déposer d'abord, puis les particules fines (argiles) par-dessus : un dépôt en couches se forme spontanément.

\textbf{Hypothèse :} les roches sédimentaires observées en couches dans la nature résultent du dépôt de particules transportées par l'eau, suivi d'une transformation de ces dépôts en roche.

\textbf{Données :} les sédiments proviennent de l'\cle{érosion} et de l'\cle{altération} des roches préexistantes des continents ; ils sont transportés par les cours d'eau, le vent ou les glaciers (\cle{transport}), puis déposés (\cle{sédimentation}) lorsque l'énergie du milieu diminue, le plus souvent en \cle{milieu aquatique} : mer, lac, fleuve, lagune. Les dépôts s'accumulent couche après couche, les plus récents recouvrant les plus anciens.

\textbf{Interprétation :} sous le poids des dépôts accumulés, les sédiments meubles subissent la \cle{diagénèse} : ensemble des transformations physico-chimiques (compaction, cimentation des grains par des sels minéraux dissous) qui transforment le sédiment meuble en roche sédimentaire consolidée (sable $\rightarrow$ grès ; vase $\rightarrow$ argile ou schiste ; boue calcaire $\rightarrow$ calcaire).

\textbf{Conclusion :} une couche sédimentaire est le résultat d'un cycle : érosion $\rightarrow$ transport $\rightarrow$ dépôt en milieu aquatique $\rightarrow$ diagénèse. Chaque couche enregistre donc les conditions du milieu au moment de son dépôt.
\end{experience}

\begin{aretenir}[Conditions de formation d'une couche sédimentaire]
\begin{enumerate}
\item Production de particules par érosion et altération des roches du continent ;
\item Transport de ces particules (eau, vent, glace) ;
\item Dépôt (sédimentation) en milieu aquatique (mer, lac, fleuve), les couches se superposant au fil du temps ;
\item Diagénèse : compaction et cimentation transformant le sédiment en roche.
\end{enumerate}
La présence de coquillages marins fossilisés dans une couche indique donc que cette couche s'est déposée \emph{dans la mer} : c'est ainsi que l'on comprend que la région de Yaoundé a été recouverte par la mer à certaines époques de son histoire.
\end{aretenir}

\courssub{Les caractères d'une couche sédimentaire}

Pour décrire et identifier une couche, le géologue relève plusieurs caractères :

\begin{itemize}
\item \textbf{Ses limites :} le \cle{mur} (base, en contact avec la couche sous-jacente plus ancienne) et le \cle{toit} (sommet, en contact avec la couche sus-jacente plus récente). La connaissance du mur et du toit définit la \cle{polarité} de la couche, c'est-à-dire le sens « ancien $\rightarrow$ récent ».
\item \textbf{Son épaisseur :} distance entre le mur et le toit, mesurée perpendiculairement aux limites ; elle peut varier latéralement (la couche peut s'amincir jusqu'à disparaître).
\item \textbf{Sa composition (lithologie) :} nature des particules et du ciment : sable/grès, argile, calcaire, etc. Elle renseigne sur le milieu de dépôt (un calcaire à débris de coquilles indique un milieu marin ; un sable grossier, un milieu fluviatile ou littoral).
\item \textbf{Son contenu paléontologique :} les \cle{fossiles} (restes ou traces d'organismes) qu'elle contient. Ils renseignent à la fois sur le milieu de dépôt et sur l'âge relatif de la couche (nous y reviendrons dans la suite de la leçon avec les fossiles stratigraphiques).
\end{itemize}

\begin{popfigure}
\begin{center}
\begin{tikzpicture}[scale=1.0]
% Couches
\fill[popgoldL] (0,0) rectangle (9,1.2);
\fill[popgreenL] (0,1.2) rectangle (9,2.4);
\fill[popblueL] (0,2.4) rectangle (9,3.8);
\fill[poporangeL] (0,3.8) rectangle (9,4.8);
% Contours
\draw[thick,popdark] (0,0) rectangle (9,4.8);
\draw[thick,popdark] (0,1.2) -- (9,1.2);
\draw[thick,popdark] (0,2.4) -- (9,2.4);
\draw[thick,popdark] (0,3.8) -- (9,3.8);
% Noms des couches
\node at (4.5,0.6) {\small Couche 1 : grès (la plus ancienne)};
\node at (4.5,1.8) {\small Couche 2 : argile};
\node at (4.5,3.1) {\small Couche 3 : calcaire à coquillages};
\node at (4.5,4.3) {\small Couche 4 : sable fin (la plus récente)};
% Fossiles dans le calcaire
\draw[popdark] (2.2,3.0) circle (0.12);
\draw[popdark] (6.5,3.3) circle (0.12);
% Toit et mur
\draw[thick,popink] (9,3.8) -- (9.7,3.8);
\node[right,popink] at (9.7,3.8) {\small toit de la couche 3};
\draw[thick,popink] (9,2.4) -- (9.7,2.4);
\node[right,popink] at (9.7,2.4) {\small mur de la couche 3};
% Épaisseur
\draw[<->,thick,popblue] (-0.5,2.4) -- (-0.5,3.8);
\node[left,popblue] at (-0.5,3.1) {\small épaisseur};
% Flèche du temps
\draw[->,very thick,poporange] (10.6,0.4) -- (10.6,4.6);
\node[rotate=90,poporange] at (11.0,2.5) {\small du plus ancien au plus récent};
\end{tikzpicture}
\end{center}
\legende{Coupe géologique schématique montrant quatre couches sédimentaires horizontales superposées. Chaque couche est limitée par un mur (base) et un toit (sommet). Dans une série non perturbée, l'âge des couches augmente du toit vers le mur, c'est-à-dire de haut en bas.}
\end{popfigure}

\courssub{La datation relative : ordonner sans chronométrer}

\textbf{Situation concrète.} Dans une famille, on peut dire sans hésiter que le grand-père est né avant le père, et le père avant le fils, sans connaître les dates de naissance exactes de chacun. On établit un \emph{ordre}, pas des dates. La datation relative en géologie procède exactement de la même manière.

\begin{definition}[Datation relative]
La \cle{datation relative} est l'ensemble des méthodes qui permettent de déterminer l'\cle{ordre chronologique} des événements géologiques (dépôt des couches, plissement, fracturation, érosion, intrusion magmatique) les uns par rapport aux autres, en précisant lequel est le plus ancien et lequel est le plus récent, \emph{sans} connaître leur âge exact en années. Elle s'oppose à la \cle{datation absolue}, qui attribue un âge chiffré (en millions d'années) grâce notamment à la radioactivité des roches.
\end{definition}

\begin{popfigure}
\begin{center}
\begin{tikzpicture}[scale=1.0]
% Datation relative
\node[font=\small\bfseries,popblue] at (3.2,2.9) {Datation relative};
\draw[->,thick,popblue] (0,2.2) -- (6.4,2.2);
\foreach \x/\t in {0.4/{dépôt A}, 2.4/{dépôt B}, 4.4/{plissement}} {
  \fill[popblue] (\x,2.2) circle (0.07);
  \node[above,popblue,font=\scriptsize] at (\x,2.2) {\t};
}
\node[below,popblue,font=\scriptsize,align=center] at (3.2,2.0) {ordre seul : A avant B avant le plissement\\ (pas d'âge en années)};
% Datation absolue
\node[font=\small\bfseries,poporange] at (3.2,0.9) {Datation absolue};
\draw[->,thick,poporange] (0,0.2) -- (6.4,0.2);
\foreach \x/\t in {0.4/{540 Ma}, 2.4/{480 Ma}, 4.4/{300 Ma}} {
  \fill[poporange] (\x,0.2) circle (0.07);
  \node[below,poporange,font=\scriptsize] at (\x,0.2) {\t};
}
\node[above,poporange,font=\scriptsize] at (3.2,0.35) {âges chiffrés en millions d'années (Ma)};
\end{tikzpicture}
\end{center}
\legende{Comparaison des deux démarches de datation. La datation relative ordonne les événements les uns par rapport aux autres ; la datation absolue leur attribue un âge chiffré. Les deux démarches sont complémentaires.}
\end{popfigure}

\begin{attention}[Piège classique : la couche la plus récente n'est pas toujours en haut !]
Dans une série sédimentaire \emph{non perturbée}, la couche la plus récente est bien au sommet. Mais si le terrain a été renversé par un \cle{pli couché}, la série peut être retournée : la couche la plus ancienne se retrouve alors au-dessus de la plus récente ! Avant toute conclusion, il faut donc vérifier la \cle{polarité} des couches grâce aux critères de polarité (position du mur et du toit, figures sédimentaires, contenu fossilifère). Ne jamais appliquer mécaniquement « en haut = récent » sans contrôle.
\end{attention}

\courssub{Intérêt pratique de la stratigraphie}

Pourquoi étudier les couches ? Au-delà de la connaissance scientifique, la stratigraphie a des applications concrètes majeures :

\begin{itemize}
\item \textbf{Reconstituer l'histoire géologique d'une région :} transgressions et régressions marines, phases de plissement, de fracturation, d'érosion — autant d'événements qui expliquent le paysage actuel (relief, répartition des roches).
\item \textbf{Recherche des ressources naturelles :} le \cle{pétrole} et le gaz s'accumulent dans certaines couches sédimentaires poreuses (réservoirs) recouvertes de couches imperméables (couvertures) ; l'\cle{eau souterraine} circule dans les aquifères sableux ; les matériaux de construction (sable, calcaire, argile) s'exploitent dans des couches précises. Connaître l'ordre et la géométrie des couches permet de localiser ces ressources.
\end{itemize}

\begin{savaistu}[Le pétrole du Rio del Rey]
Le bassin du Rio del Rey, au large du Littoral camerounais (entre Douala et la frontière nigériane), est la principale zone de production pétrolière du Cameroun. Le pétrole y est piégé dans des couches sédimentaires du delta du Niger, dont la chronologie a été établie par stratigraphie : les géologues ont ordonné les dépôts successifs du delta, identifié les couches réservoirs (sables poreux) et les couches couvertures (argiles imperméables) qui emprisonnent les hydrocarbures. Sans la connaissance de l'ordre et de la géométrie des couches, aucun forage pétrolier ne serait possible !
\end{savaistu}

\begin{exoresolu}[Identifier les caractères d'une couche sur un affleurement]
\textbf{Énoncé.} Sur une falaise littorale de la région de Kribi, un géologue observe de bas en haut : une couche A de sable consolidé (grès) de 2~m d'épaisseur, puis une couche B de calcaire de 1~m d'épaisseur contenant de nombreux restes de coquillages marins. (1) Indiquer le mur et le toit de la couche B. (2) Quel milieu de dépôt peut-on déduire de la couche B ? Justifier. (3) Quelle couche est la plus ancienne ? On précise que la série n'est pas renversée.
\tcblower
\textbf{Solution détaillée.}
\begin{enumerate}
\item Le \cle{mur} de la couche B est sa surface inférieure, en contact avec la couche A (grès) ; son \cle{toit} est sa surface supérieure, en contact avec l'air ou la couche sus-jacente.
\item La couche B contient des restes de \cle{coquillages marins} : or les fossiles indiquent le milieu dans lequel vivaient les organismes. La couche B s'est donc déposée en \cle{milieu marin} : la mer recouvrait la région de Kribi au moment de ce dépôt.
\item La série n'étant pas renversée, la couche située en dessous est la plus ancienne : la couche A (grès) est plus ancienne que la couche B (calcaire). Il s'agit d'une application du principe de superposition, que nous étudierons en détail dans la section suivante.
\end{enumerate}
\end{exoresolu}

\begin{aretenir}[L'essentiel de la section]
\begin{itemize}
\item La \cle{stratigraphie} étudie les couches (strates) des roches sédimentaires et leurs relations pour reconstituer l'histoire géologique.
\item Une couche se forme par dépôt de sédiments en milieu aquatique, suivi de la \cle{diagénèse} ; elle est caractérisée par son \cle{mur}, son \cle{toit}, son épaisseur, sa composition et ses fossiles.
\item La \cle{datation relative} ordonne les événements géologiques les uns par rapport aux autres, sans âge chiffré ; elle se fonde sur des principes (superposition, continuité, recoupement) et sur les fossiles, que nous étudierons dans les sections suivantes.
\item Attention : en terrain plissé ou renversé, toujours vérifier la \cle{polarité} des couches avant de conclure.
\end{itemize}
\end{aretenir}

\coursec{Les principes de la stratigraphie}

Dans la section précédente, nous avons découvert ce qu'est une \cle{couche sédimentaire} et comment les sédiments s'accumulent, se lithifient et forment des séries stratigraphiques que le géologue peut observer en affleurement — le long d'une falaise, d'une carrière ou d'une tranchée de route. Mais observer des couches ne suffit pas : il faut pouvoir \emph{lire} leur histoire. Comment savoir laquelle est la plus ancienne ? Comment affirmer qu'une faille est plus récente que les roches qu'elle traverse ? Comment relier entre eux des affleurements distants de plusieurs kilomètres ?

C'est toute l'ambition des \cle{principes de la stratigraphie} : trois règles simples, énoncées dès le XVII\textsuperscript{e} siècle par le savant danois Nicolas Sténon (Niels Stensen), qui permettent de transformer une coupe géologique muette en un véritable récit chronologique. Ces principes sont la \emph{grammaire} de la lecture des terrains. Nous allons les énoncer, les justifier logiquement, puis apprendre à les appliquer ensemble.

\courssub{Le principe de superposition}

\textbf{Une observation de la vie courante.} Imagine une pile de dossiers sur le bureau d'un directeur d'hôpital à Yaoundé : chaque jour, la secrétaire dépose les nouveaux dossiers \emph{par-dessus} les précédents. Sans même lire les dates, tu sais que le dossier du fond de la pile est le plus ancien, et celui du sommet le plus récent. La nature procède exactement de la même manière avec les sédiments.

\textbf{Observation géologique.} Sur une falaise ou une carrière, on observe des couches empilées les unes sur les autres. Aucune couche ne peut être déposée \emph{sous} une couche déjà existante : le sédiment tombe toujours \emph{sur} ce qui existe déjà. Cette remarque, d'une simplicité enfantine, fonde toute la datation relative.

\begin{propriete}[Principe de superposition]
Dans une série stratigraphique \cle{non renversée} (c'est-à-dire restée dans sa position initiale de dépôt), une couche est \textbf{plus ancienne} que la couche qui la \textbf{recouvre} (au-dessus d'elle) et \textbf{plus récente} que la couche qui la \textbf{supporte} (en dessous d'elle). L'âge des couches \emph{diminue donc de la base vers le sommet} de la série.
\end{propriete}

\begin{experience}[Démonstration guidée : justification logique du principe de superposition]
\textbf{Problème :} pourquoi peut-on affirmer que, dans une série non renversée, la couche du bas est la plus ancienne ?

\textbf{Démarche :}
\begin{enumerate}
\item \textbf{Observation :} dans un vase ou un bac transparent, versons successivement du sable grossier, puis du sable fin, puis de l'argile, en laissant chaque apport se déposer avant le suivant.
\item \textbf{Constat :} le premier apport (sable grossier) forme la couche du \emph{fond} ; le dernier (argile) forme la couche du \emph{sommet}. Chaque nouvelle couche ne peut se déposer que \emph{sur} une surface préexistante.
\item \textbf{Raisonnement :} soit une couche $C_2$ recouvrant une couche $C_1$. Pour que $C_2$ puisse se déposer sur $C_1$, il faut que $C_1$ \emph{existe déjà}. Donc $C_1$ est antérieure à $C_2$. Par récurrence, si une série comporte les couches $C_1, C_2, \ldots, C_n$ de bas en haut, alors $C_1$ est antérieure à $C_2$, elle-même antérieure à $C_3$, et ainsi de suite jusqu'à $C_n$.
\item \textbf{Conclusion :} dans une série non renversée, l'ordre vertical des couches (du bas vers le haut) reproduit l'ordre chronologique de leur dépôt (du plus ancien au plus récent). Le principe de superposition est ainsi une conséquence directe et nécessaire du processus de dépôt successif des sédiments.
\end{enumerate}
\end{experience}

\begin{popfigure}
\begin{tikzpicture}[scale=1]
% Couches
\fill[popgoldL] (0,0) rectangle (7,1);
\fill[poporangeL] (0,1) rectangle (7,2);
\fill[popgreenL] (0,2) rectangle (7,3);
\fill[popblueL] (0,3) rectangle (7,4);
\fill[poppinkL] (0,4) rectangle (7,5);
\draw[thick] (0,0) rectangle (7,5);
\foreach \y in {1,2,3,4} \draw[thick] (0,\y) -- (7,\y);
% Numéros
\node at (3.5,0.5) {\textbf{1}};
\node at (3.5,1.5) {\textbf{2}};
\node at (3.5,2.5) {\textbf{3}};
\node at (3.5,3.5) {\textbf{4}};
\node at (3.5,4.5) {\textbf{5}};
% Flèche âge
\draw[-{Stealth[length=3mm]},very thick,popink] (7.6,0.3) -- (7.6,4.7);
\node[rotate=90,popink] at (8.15,2.5) {\textbf{du plus ancien au plus récent}};
% Légendes latérales
\draw[thin,popmuted] (0,0.5) -- (-0.4,0.5) node[left,align=right] {couche 1 : la plus ancienne};
\draw[thin,popmuted] (0,4.5) -- (-0.4,4.5) node[left,align=right] {couche 5 : la plus récente};
\end{tikzpicture}
\legende{Série stratigraphique de cinq couches non renversées. D'après le principe de superposition, la couche 1 (base) est la plus ancienne et la couche 5 (sommet) la plus récente : l'âge diminue de la base vers le sommet.}
\end{popfigure}

\begin{attention}[La condition « non renversée » est indispensable]
Le principe de superposition ne s'applique que si la série a conservé sa \textbf{polarité initiale}. Or, les déformations tectoniques peuvent \emph{renverser} une série entière (plis couchés, charriages) : la couche la plus ancienne se retrouve alors \emph{au-dessus} ! Avant d'appliquer le principe, il faut donc vérifier que la série est « à l'endroit », grâce à des critères de polarité (figures sédimentaires, contenu fossilifère). Au Probatoire, sauf indication contraire de l'énoncé, les séries sont considérées non renversées.
\end{attention}

\courssub{Le principe de continuité latérale}

\textbf{Situation concrète.} Sur la route de l'Ouest camerounais, on observe parfois, de part et d'autre d'une vallée, des falaises présentant des couches étonnamment semblables : même couleur, même nature de roche, même succession. S'agit-il d'une coïncidence ? Non : ces couches étaient autrefois \emph{réunies} en une seule nappe continue, que l'érosion a ensuite découpée.

\textbf{Justification.} Lors de son dépôt, un sédiment s'étale latéralement sur toute la surface du bassin de sédimentation, comme une nappe d'eau recouvre uniformément le fond d'une cuvette. Une couche possède donc, au moment de sa formation, une \cle{extension latérale continue}, limitée seulement par les bords du bassin. Si l'on retrouve aujourd'hui, en deux points distants, des couches de \emph{même nature} (même roche) et de \emph{même contenu fossilifère} (mêmes fossiles), il est logique de considérer qu'il s'agit de \emph{la même couche}, déposée au \emph{même moment}, puis séparée par l'érosion ou masquée sous la végétation.

\begin{propriete}[Principe de continuité latérale (admis)]
Une couche sédimentaire possède, au moment de son dépôt, une extension latérale continue dans tout le bassin de sédimentation. Par conséquent, deux couches de \textbf{même nature lithologique} et de \textbf{même contenu fossilifère}, observées en deux affleurements distincts, peuvent être \cle{corrélées} : elles correspondent au même niveau stratigraphique et ont le \textbf{même âge}. Ce principe, admis, est justifié par l'observation directe des affleurements que l'on peut suivre latéralement sur le terrain.
\end{propriete}

\begin{popfigure}
\begin{tikzpicture}[scale=1]
% Affleurement A (gauche)
\fill[popgoldL] (0,0) rectangle (2.6,0.8);
\fill[popgreenL] (0,0.8) rectangle (2.6,1.6);
\fill[popblueL] (0,1.6) rectangle (2.6,2.4);
\draw[thick] (0,0) rectangle (2.6,2.4);
\draw[thick] (0,0.8) -- (2.6,0.8);
\draw[thick] (0,1.6) -- (2.6,1.6);
\node at (1.3,0.4) {\small $a$};
\node at (1.3,1.2) {\small $b$};
\node at (1.3,2.0) {\small $c$};
\node[below] at (1.3,0) {\small Affleurement A};
% Vallée (végétation)
\fill[poppinkL] (2.6,0) rectangle (5.4,0.5);
\node[popmuted] at (4,0.25) {\small zone masquée};
% Affleurement B (droite)
\fill[popgoldL] (5.4,0) rectangle (8,0.8);
\fill[popgreenL] (5.4,0.8) rectangle (8,1.6);
\fill[popblueL] (5.4,1.6) rectangle (8,2.4);
\draw[thick] (5.4,0) rectangle (8,2.4);
\draw[thick] (5.4,0.8) -- (8,0.8);
\draw[thick] (5.4,1.6) -- (8,1.6);
\node at (6.7,0.4) {\small $a'$};
\node at (6.7,1.2) {\small $b'$};
\node at (6.7,2.0) {\small $c'$};
\node[below] at (6.7,0) {\small Affleurement B};
% Lignes de corrélation
\draw[dashed,thick,popink] (2.6,0.8) -- (5.4,0.8);
\draw[dashed,thick,popink] (2.6,1.6) -- (5.4,1.6);
\draw[dashed,thick,popink] (2.6,2.4) -- (5.4,2.4);
\node[popink,above] at (4,2.4) {\small corrélations};
\end{tikzpicture}
\legende{Principe de continuité latérale. Les couches $a$, $b$, $c$ de l'affleurement A et les couches $a'$, $b'$, $c'$ de l'affleurement B ont même nature et même contenu fossilifère : les lignes pointillées matérialisent leur corrélation. Les couches $a$ et $a'$ (de même $b/b'$ et $c/c'$) ont le même âge.}
\end{popfigure}

\begin{exemplebox}[Corréler deux carrières]
Deux carrières situées à 15 km l'une de l'autre montrent chacune, à leur base, un calcaire gris renfermant les mêmes coquilles d'ammonites, surmonté d'une argile rouge. D'après le principe de continuité latérale, le calcaire gris de la carrière 1 et celui de la carrière 2 forment \emph{une seule et même couche}, déposée au même instant dans un même bassin ; il en va de même pour l'argile rouge. On peut donc reconstituer une colonne stratigraphique unique valable pour toute la région.
\end{exemplebox}

\begin{attention}[Les conditions de la corrélation]
On ne peut corréler deux couches que si \textbf{deux critères} sont réunis : même nature de roche \emph{et} même contenu fossilifère. Une simple ressemblance de couleur ou de texture ne suffit pas : deux dépôts de sable d'âges très différents peuvent se ressembler ! De plus, une même couche peut changer de nature latéralement (\emph{faciès latéral}) : le principe s'applique alors avec prudence.
\end{attention}

\courssub{Le principe de recoupement}

\textbf{Intuition.} Pour casser un morceau de pain, encore faut-il que le pain soit déjà cuit : on ne peut pas rompre ce qui n'existe pas. De même, une faille ne peut fracturer que des couches \emph{déjà formées}, et un filon de roche magmatique ne peut s'injecter que dans des roches \emph{déjà en place}. Cette évidence logique constitue le troisième grand principe de la stratigraphie.

\begin{propriete}[Principe de recoupement]
Toute structure géologique (\cle{faille}, \cle{filon}, \cle{pli}) qui \textbf{recoupe} des couches sédimentaires est \textbf{postérieure} (plus récente) à ces couches. Réciproquement, une structure \textbf{recoupée} par une autre est \textbf{antérieure} (plus ancienne) à celle qui la recoupe.
\end{propriete}

\begin{experience}[Démonstration guidée : la faille ne peut fracturer que des couches déjà existantes]
\textbf{Problème :} une faille $F$ décale trois couches $C_1$, $C_2$, $C_3$. Peut-on dater la faille par rapport aux couches ?

\textbf{Démarche :}
\begin{enumerate}
\item \textbf{Observation :} sur la coupe, la faille $F$ traverse les trois couches et les décale verticalement : les bords de chaque couche ne sont plus en regard de part et d'autre de la faille.
\item \textbf{Hypothèse à tester :} la faille serait-elle \emph{antérieure} aux couches ? Si tel était le cas, les couches se seraient déposées \emph{après} la fracturation : elles recouvriraient la faille sans être décalées par elle, et leurs bords resteraient continus de part et d'autre.
\item \textbf{Confrontation aux faits :} or on observe que les couches \emph{sont décalées} par la faille. L'hypothèse est donc rejetée.
\item \textbf{Conclusion :} les couches $C_1$, $C_2$, $C_3$ existaient \emph{déjà} lorsque la fracturation s'est produite. La faille $F$ est donc \textbf{postérieure} aux trois couches qu'elle recoupe. Le raisonnement est identique pour un filon (injection de magma dans des fissures de roches préexistantes) ou un pli (déformation de couches déjà déposées).
\end{enumerate}
\end{experience}

\begin{popfigure}
\begin{tikzpicture}[scale=1]
% Bloc gauche de la faille
\fill[popgoldL] (0,0) -- (3.4,0) -- (3.4,1.0) -- (0,1.0) -- cycle;
\fill[popgreenL] (0,1.0) -- (3.4,1.0) -- (3.4,2.0) -- (0,2.0) -- cycle;
\fill[popblueL] (0,2.0) -- (3.4,2.0) -- (3.4,3.0) -- (0,3.0) -- cycle;
% Bloc droit (décalé vers le haut)
\fill[popgoldL] (3.4,0.7) -- (7,0.7) -- (7,1.7) -- (3.4,1.7) -- cycle;
\fill[popgreenL] (3.4,1.7) -- (7,1.7) -- (7,2.7) -- (3.4,2.7) -- cycle;
\fill[popblueL] (3.4,2.7) -- (7,2.7) -- (7,3.7) -- (3.4,3.7) -- cycle;
% Contours
\draw[thick] (0,0) rectangle (3.4,3.0);
\draw[thick] (0,1.0) -- (3.4,1.0);
\draw[thick] (0,2.0) -- (3.4,2.0);
\draw[thick] (3.4,0.7) rectangle (7,3.7);
\draw[thick] (3.4,1.7) -- (7,1.7);
\draw[thick] (3.4,2.7) -- (7,2.7);
% Faille
\draw[very thick,popink] (3.4,-0.1) -- (3.4,3.9);
\node[popink,above] at (3.4,3.9) {\textbf{faille $F$}};
% Étiquettes couches
\node at (1.7,0.5) {\small $C_1$};
\node at (1.7,1.5) {\small $C_2$};
\node at (1.7,2.5) {\small $C_3$};
\node at (5.2,1.2) {\small $C_1$};
\node at (5.2,2.2) {\small $C_2$};
\node at (5.2,3.2) {\small $C_3$};
% Flèches de rejet
\draw[-{Stealth[length=2.5mm]},thick,popdark] (3.1,3.5) -- (3.1,4.3);
\draw[-{Stealth[length=2.5mm]},thick,popdark] (3.7,3.2) -- (3.7,2.4);
% Relations d'antériorité
\node[align=center,below] at (1.7,-0.15) {\small couches $C_1, C_2, C_3$ : \textbf{antérieures}};
\node[align=center,below] at (5.2,0.45) {\small faille $F$ : \textbf{postérieure}};
\end{tikzpicture}
\legende{Principe de recoupement. La faille $F$ recoupe et décale les couches $C_1$, $C_2$, $C_3$ (les flèches indiquent le mouvement relatif des blocs). La faille est donc postérieure aux trois couches recoupées, qui lui sont antérieures.}
\end{popfigure}

\begin{methode}[Appliquer le principe de recoupement sur une coupe]
\begin{enumerate}
\item \textbf{Repérer} toutes les structures de la coupe : couches, failles, filons, plis.
\item Pour chaque structure, \textbf{identifier ce qui est recoupé} : toute couche ou structure traversée, décalée ou déformée par une autre est \emph{plus ancienne}.
\item \textbf{Identifier ce qui recoupe} : la faille, le filon ou le pli qui traverse les autres éléments est \emph{plus récent} que chacun d'eux.
\item \textbf{Hiérarchiser} : classer tous les éléments du plus ancien au plus récent en combinant recoupement et superposition.
\item \textbf{Vérifier} la cohérence : un même élément ne peut pas être à la fois antérieur et postérieur à un autre.
\end{enumerate}
\end{methode}

\begin{attention}[Erreur fréquente : la faille n'est pas plus récente que TOUTES les couches]
Beaucoup d'élèves écrivent : « la faille est l'événement le plus récent de la coupe ». C'est \textbf{faux} en général ! Une faille n'est plus récente \emph{que des couches qu'elle recoupe effectivement}. Si une couche $C_4$ recouvre la faille \emph{sans être décalée} par elle, alors $C_4$ est \emph{postérieure} à la faille : le dépôt de $C_4$ s'est fait après la fracturation, sur un relief déjà érodé. La faille se situe alors \emph{entre} le dépôt de la dernière couche recoupée et celui de la première couche non recoupée. Observe toujours si la faille atteint le sommet de la série ou s'arrête sous certaines couches.
\end{attention}

\courssub{Application combinée des trois principes}

Sur une coupe réelle, les trois principes s'utilisent \emph{ensemble} : la superposition ordonne les couches entre elles, la continuité latérale relie les affleurements, et le recoupement insère les structures (failles, filons, plis) dans la chronologie. Entraînons-nous sur un cas complet.

\begin{exoresolu}[Hiérarchiser les événements d'une coupe]
\textbf{Énoncé.} Une coupe géologique montre, de bas en haut, les couches sédimentaires $C_1$, $C_2$ et $C_3$ (série non renversée). Une faille $F$ recoupe $C_1$ et $C_2$ mais s'arrête sous $C_3$, qui n'est pas décalée. Un filon de granite $G$ recoupe $C_1$, $C_2$, $C_3$ et la faille $F$. Établir la chronologie relative de tous ces événements, du plus ancien au plus récent, en justifiant chaque étape.
\tcblower
\textbf{Solution détaillée.}
\begin{enumerate}
\item \textbf{Ordre des couches (superposition) :} la série n'est pas renversée, donc $C_1$ est plus ancienne que $C_2$, elle-même plus ancienne que $C_3$ : $C_1 \rightarrow C_2 \rightarrow C_3$.
\item \textbf{Position de la faille (recoupement) :} $F$ recoupe $C_1$ et $C_2$, donc $F$ est \emph{postérieure} à $C_1$ et $C_2$. Mais $C_3$ recouvre $F$ sans être recoupée, donc $C_3$ est \emph{postérieure} à $F$. La faille se place donc entre $C_2$ et $C_3$ : $C_1 \rightarrow C_2 \rightarrow F \rightarrow C_3$.
\item \textbf{Position du filon (recoupement) :} $G$ recoupe $C_1$, $C_2$, $C_3$ \emph{et} la faille $F$ : il est postérieur à chacun de ces éléments. C'est donc l'événement le plus récent de la coupe.
\item \textbf{Chronologie finale (du plus ancien au plus récent) :}
\[ C_1 \;\rightarrow\; C_2 \;\rightarrow\; \text{faille } F \;\rightarrow\; C_3 \;\rightarrow\; \text{filon } G \]
\end{enumerate}
\textbf{Vérification :} chaque structure est bien postérieure à tout ce qu'elle recoupe et antérieure à tout ce qui la recoupe ; la chronologie est cohérente.
\end{exoresolu}

\courssub{Limites des principes de stratigraphie}

Ces trois principes sont puissants, mais leur application exige de la vigilance : le terrain réel est plus compliqué que les schémas.

\begin{itemize}
\item \textbf{Les séries renversées.} Les plis couchés et les charriages peuvent retourner une série entière : la superposition donne alors un ordre \emph{inverse} de la chronologie. Il faut d'abord déterminer la polarité des couches (grâce aux figures sédimentaires ou aux fossiles) avant d'appliquer le principe.
\item \textbf{Les lacunes de sédimentation.} La sédimentation n'est pas continue : pendant de longues périodes, aucun dépôt ne se produit (ou les dépôts sont ensuite érodés). Une série peut donc « sauter » des millions d'années sans qu'aucune couche n'en témoigne : deux couches superposées ne sont pas forcément d'âges proches.
\item \textbf{L'érosion.} L'érosion peut faire disparaître des couches entières avant le dépôt des suivantes : la chronologie lue sur la coupe est alors \emph{incomplète}. C'est l'origine des \cle{discordances}, que nous étudierons dans la section suivante.
\item \textbf{Les limites du recoupement.} Une structure ne date que ce qu'elle recoupe : si une faille n'affleure que dans les couches inférieures, on ne peut rien dire de son âge par rapport aux couches qu'elle n'atteint pas.
\end{itemize}

\begin{aretenir}[L'essentiel de la section]
\begin{itemize}
\item \textbf{Superposition :} dans une série non renversée, l'âge des couches diminue de la base vers le sommet ; chaque couche est plus ancienne que celle qui la recouvre et plus récente que celle qui la supporte.
\item \textbf{Continuité latérale :} deux couches de même nature et de même contenu fossilifère, observées en deux points, peuvent être corrélées et ont le même âge.
\item \textbf{Recoupement :} ce qui recoupe est plus récent ; ce qui est recoupé est plus ancien. Une faille n'est plus récente que des couches qu'elle recoupe effectivement.
\item La combinaison des trois principes permet d'établir la \textbf{chronologie relative complète} d'une région, mais il faut se méfier des séries renversées, des lacunes de sédimentation et de l'érosion.
\end{itemize}
\end{aretenir}

Ces trois principes nous donnent l'ordre des événements, mais pas encore de repères précis dans le temps. Pour affiner la datation relative et dater une formation avec plus de précision, les géologues disposent d'un outil remarquable : les \emph{fossiles stratigraphiques}, que nous étudierons dans la section 4 — après avoir examiné, dans la section 3, ces archives spectaculaires de l'histoire de la Terre que sont les discordances.

\coursec{Les discordances stratigraphiques}

Dans la section précédente, nous avons vu que les principes de superposition, de continuité et de recoupement permettent d'ordonner des couches \emph{tant qu'elles sont parallèles et continues}. Mais que se passe-t-il lorsque l'on observe, sur une coupe géologique, des couches inclinées, voire pliées, qui sont brutalement tronquées par une surface plane ou ondulée, surmontée de nouvelles couches parfaitement horizontales ? L'histoire géologique ne s'est pas déroulée en un seul tenant : il y a eu une \textbf{interruption}, une \textbf{lacune} dans l'enregistrement sédimentaire. C'est ce que l'on appelle une \textbf{discordance stratigraphique}. Comprendre ces surfaces est la clé pour reconstituer des histoires géologiques complexes, comme celle des bassins sédimentaires du Cameroun (Douala, Garoua) qui ont connu des phases de déformation et d'érosion avant de recevoir de nouveaux sédiments.

\courssub{Définition et nature d'une discordance}

\begin{definition}[Discordance stratigraphique]
Une \cle{discordance stratigraphique} est une surface d'érosion (ou de non-dépôt) qui sépare deux ensembles de couches sédimentaires d'âges différents et dont les \textbf{orientations (pendages) sont distinctes}. Elle matérialise une \textbf{lacune sédimentaire} (un intervalle de temps non enregistré par des roches) correspondant à une phase d'émersion, d'érosion et/ou de déformation tectonique.
\end{definition}

\begin{attention}[Discordance vs Faille : ne pas confondre !]
Sur une coupe, une faille et une discordance peuvent toutes deux montrer une rupture brutale. Mais :
\begin{itemize}
    \item La \textbf{faille} est une \cle{fracture} avec déplacement relatif des compartiments (rejet). Les couches de part et d'autre sont \emph{de même âge} et \emph{même nature} (continuité lithologique rompue). On observe souvent un miroir de faille, des striations, une brèche de faille.
    \item La \textbf{discordance} est une \cle{surface d'érosion}. Les couches \emph{sous-jacentes} sont plus anciennes, souvent déformées (plissées, redressées), et \emph{tronquées} (leurs extrémités sont coupées). Les couches \emph{suprajacentes} sont plus récentes, généralement horizontales, et viennent \emph{recouvrir} la surface d'érosion en l'épousant (transgression). Il n'y a pas de déplacement tectonique \emph{au moment} de la mise en place des couches supérieures.
\end{itemize}
\end{attention}

\courssub{Genèse d'une discordance angulaire : une histoire en quatre temps}

La discordance la plus parlante est la \cle{discordance angulaire} (ou discordance de pendage). Sa formation suit un scénario type que l'on peut reconstituer par une démarche d'investigation à partir de l'observation d'une coupe.

\begin{experience}[Modélisation de la formation d'une discordance angulaire]
\textbf{Matériel :} Bac transparent, sable de deux couleurs (ex. jaune et rouge), eau, règle, spatule, bloc de bois pour simuler la poussée tectonique.
\textbf{Protocole :}
\begin{enumerate}
    \item \textbf{Étape 1 (Dépôt initial) :} Déposer une couche de sable jaune horizontalement au fond du bac (simule une mer peu profonde, sédimentation continue). Tasser légèrement.
    \item \textbf{Étape 2 (Déformation tectonique) :} Pousser le bloc de bois contre le sable jaune pour le plisser/redresser (simule une compression orogénique, ex. chaîne de montagne).
    \item \textbf{Étape 3 (Émersion et érosion) :} Retirer l'eau. Avec la spatule, « éroder » le sommet des plis pour créer une surface plane ou légèrement ondulée (surface de \cle{planation}). Observer les couches jaunes tronquées.
    \item \textbf{Étape 4 (Transgression et nouveau dépôt) :} Remettre l'eau. Déposer le sable rouge horizontalement sur la surface érodée.
\end{enumerate}
\textbf{Observations :} Les couches jaunes sont plissées et leurs sommets sont coupés. La surface de contact est irrégulière (paléorelief). Les couches rouges sont horizontales, parallèles à la surface de contact, et reposent sur les tronçons des couches jaunes.
\textbf{Interprétation :} La surface de contact est une discordance angulaire. Elle sépare deux cycles sédimentaires distincts séparés par une phase orogénique (plissement) et une phase d'érosion.
\textbf{Conclusion :} Une discordance angulaire témoigne toujours d'un cycle \textbf{dépôt $\rightarrow$ déformation $\rightarrow$ érosion $\rightarrow$ nouveau dépôt}.
\end{experience}

\begin{popfigure}
\begin{tikzpicture}[scale=0.9, every node/.style={font=\small}]
    % Styles
    \tikzset{
        layer1/.style={fill=poporange!40, draw=poporange!80!black, pattern=north east lines, pattern color=poporange!60!black},
        layer2/.style={fill=popblue!30, draw=popblue!80!black, pattern=north west lines, pattern color=popblue!60!black},
        erosion/.style={draw=popdark, thick, decorate, decoration={zigzag, segment length=4pt, amplitude=1.5pt}},
        arrow/.style={->, >=Stealth, thick, popdark},
        labelbox/.style={rectangle, rounded corners, fill=white, draw=popdark, inner sep=3pt, font=\small\bfseries}
    }
    \usetikzlibrary{patterns, decorations.pathmorphing, arrows.meta, positioning, calc}

    % Coordonnées de base
    \def\basey{0}
    \def\toplayerone{2.5}
    \def\erosiontop{3.0}
    \def\toplayerTwo{5.0}
    \def\leftx{0}
    \def\rightx{6}

    % --- ÉTAPE 1 : Dépôt initial ---
    \begin{scope}[xshift=0cm, yshift=0cm]
        \node[labelbox, anchor=north] at (3, 0.3) {Étape 1 : Dépôt marin};
        \fill[layer1] (\leftx,\basey) rectangle (\rightx,\toplayerone);
        \draw[thin, poporange!60!black] (\leftx,0.5) -- (\rightx,0.5);
        \draw[thin, poporange!60!black] (\leftx,1.0) -- (\rightx,1.0);
        \draw[thin, poporange!60!black] (\leftx,1.5) -- (\rightx,1.5);
        \draw[thin, poporange!60!black] (\leftx,2.0) -- (\rightx,2.0);
        \node[align=center, text width=5cm] at (3, -1.2) {Sédimentation continue\\dans une mer calme\\Couches horizontales};
    \end{scope}

    % --- ÉTAPE 2 : Plissement ---
    \begin{scope}[xshift=8cm, yshift=0cm]
        \node[labelbox, anchor=north] at (3, 0.3) {Étape 2 : Plissement (Orogenèse)};
        % Couches plissées (simplifié : sinusoides)
        \foreach \ybase in {0.25, 0.75, 1.25, 1.75, 2.25} {
            \draw[poporange!80!black, thick, domain=\leftx:\rightx, smooth, variable=\x, samples=100]
                plot ({\x}, {\ybase + 0.4*sin(180*\x/3)});
        }
        \fill[layer1, opacity=0.5] plot[domain=\leftx:\rightx, smooth, variable=\x, samples=100] ({\x}, {0.25 + 0.4*sin(180*\x/3)}) -- (\rightx,\basey) -- (\leftx,\basey) -- cycle;
        \node[align=center, text width=5cm] at (3, -1.2) {Compression tectonique\\Plissement des couches\\Émersion du secteur};
    \end{scope}

    % --- ÉTAPE 3 : Érosion ---
    \begin{scope}[xshift=0cm, yshift=-7cm]
        \node[labelbox, anchor=north] at (3, 0.3) {Étape 3 : Érosion (Planation)};
        % Couches plissées tronquées
        \foreach \ybase in {0.25, 0.75, 1.25, 1.75, 2.25} {
            \draw[poporange!80!black, thick, domain=\leftx:\rightx, smooth, variable=\x, samples=100]
                plot ({\x}, {\ybase + 0.4*sin(180*\x/3)});
        }
        % Surface d'érosion (zigzag)
        \draw[erosion] (\leftx,\erosiontop) -- (\rightx,\erosiontop);
        % Hachures au-dessus pour montrer le vide/air
        \foreach \i in {1,...,10} {
            \draw[popmuted, thin] ({rand*0.5+3}, {3.2+rand*0.5}) -- ++({rand*0.3}, {rand*0.3});
        }
        \node[align=center, text width=5cm] at (3, -1.2) {Émersion, érosion atmosphérique\\Troncature des plis\\Surface de planation (pénéplaine)};
    \end{scope}

    % --- ÉTAPE 4 : Transgression et nouveau dépôt ---
    \begin{scope}[xshift=8cm, yshift=-7cm]
        \node[labelbox, anchor=north] at (3, 0.3) {Étape 4 : Transgression \& Nouveau dépôt};
        % Couches inférieures (tronquées)
        \foreach \ybase in {0.25, 0.75, 1.25, 1.75, 2.25} {
            \draw[poporange!80!black, thick, domain=\leftx:\rightx, smooth, variable=\x, samples=100]
                plot ({\x}, {\ybase + 0.4*sin(180*\x/3)});
        }
        % Surface de discordance
        \draw[popdark, ultra thick] (\leftx,\erosiontop) -- (\rightx,\erosiontop);
        % Couches supérieures horizontales
        \fill[layer2] (\leftx,\erosiontop) rectangle (\rightx,\toplayerTwo);
        \draw[thin, popblue!60!black] (\leftx,3.5) -- (\rightx,3.5);
        \draw[thin, popblue!60!black] (\leftx,4.0) -- (\rightx,4.0);
        \draw[thin, popblue!60!black] (\leftx,4.5) -- (\rightx,4.5);
        \node[align=center, text width=5cm] at (3, -1.2) {Avancée de la mer (transgression)\\Dépôt horizontal concordant\\sur la surface d'érosion};
    \end{scope}

    % Flèches entre étapes
    \draw[arrow] (6.5, 1.25) -- (7.5, 1.25) node[midway, above] {Poussée tectonique};
    \draw[arrow] (3, -3.5) -- (3, -4.5) node[midway, left] {Érosion};
    \draw[arrow] (10.5, -3.5) -- (10.5, -4.5) node[midway, right] {Transgression};
    \draw[arrow] (6.5, -5.75) -- (7.5, -5.75) node[midway, above] {Temps géologique};
\end{tikzpicture}
\legende{Genèse d'une discordance angulaire en quatre étapes : (1) dépôt horizontal initial, (2) plissement par compression tectonique, (3) émersion et érosion créant une surface de planation qui tronque les plis, (4) transgression marine et dépôt de nouvelles couches horizontales recouvrant la surface de discordance.}
\end{popfigure}

\courssub{Les notions de transgression et de régression}

L'étape finale de formation d'une discordance (le nouveau dépôt) est intimement liée aux variations du niveau de la mer par rapport au continent.

\begin{definition}[Transgression et Régression marines]
\begin{itemize}
    \item La \cle{transgression marine} est l'avancée de la mer sur le continent (élévation du niveau marin relatif). Elle se traduit par le dépôt de sédiments marins \textbf{sur des roches plus anciennes} (souvent une surface d'érosion continentale). Les faciès deviennent plus profonds vers le haut (ex : grès $\rightarrow$ calcaires $\rightarrow$ marnes).
    \item La \cle{régression marine} est le retrait de la mer (abaissement du niveau marin relatif). Elle se traduit par une érosion des sédiments marins ou le dépôt de faciès continentaux \textbf{sur des faciès marins}. Les faciès deviennent moins profonds vers le haut.
\end{itemize}
\end{definition}

\begin{propriete}[Marqueurs de transgression sur une discordance]
La présence d'un \cle{conglomérat basal} (galets roulés, fragments de roches sous-jacentes) au contact de la discordance, suivi de sédiments marins fins (grès, puis calcaires/marnes), est la signature classique d'une \textbf{transgression marine} postérieure à l'érosion. Ce conglomérat représente le rivage ancien (paléorivage).
\end{propriete}

\courssub{Lecture d'une discordance sur une coupe géologique}

Sur une coupe ou une carte géologique, la discordance se reconnaît à un faisceau d'indices concordants.

\begin{methode}[Reconnaître une discordance sur une coupe]
\begin{enumerate}
    \item \textbf{Observer la géométrie des couches :} Repérer deux ensembles de couches aux pendages différents (ex. couches inférieures redressées/versées vs couches supérieures horizontales).
    \item \textbf{Identifier la surface de contact :} C'est une surface souvent irrégulière (paléorelief), matérialisée par un trait épais ou en zigzag sur les schémas.
    \item \textbf{Vérifier la troncature :} Les couches inférieures s'arrêtent brutalement contre cette surface (leurs extrémités sont coupées). Elles ne la traversent pas.
    \item \textbf{Vérifier le parallélisme suprajacent :} Les couches supérieures sont parallèles à la surface de contact (elles l'épousent), témoignant d'un dépôt \emph{postérieur} à la mise en place de cette surface.
    \item \textbf{Chercher un conglomérat basal :} Sa présence au sommet de la surface de contact confirme l'érosion et la transgression.
\end{enumerate}
\end{methode}

\begin{popfigure}
\begin{tikzpicture}[scale=1.1, every node/.style={font=\small}]
    \usetikzlibrary{patterns, decorations.pathmorphing, arrows.meta, calc}
    
    % Styles
    \tikzset{
        paleozoic/.style={fill=poporange!30, draw=poporange!80!black, pattern=north east lines, pattern color=poporange!60!black},
        mesozoic/.style={fill=popblue!20, draw=popblue!80!black, pattern=north west lines, pattern color=popblue!60!black},
        cenozoic/.style={fill=popgreen!20, draw=popgreen!80!black, pattern=horizontal lines, pattern color=popgreen!60!black},
        discord/.style={draw=popdark, ultra thick, decorate, decoration={zigzag, segment length=5pt, amplitude=2pt}},
        fault/.style={draw=poppink, very thick, dashed},
        labelstrat/.style={font=\small\bfseries, text width=2.5cm, align=center},
        legendbox/.style={rectangle, draw=popdark, fill=white, rounded corners, inner sep=4pt, font=\small}
    }

    % Coupe géologique simplifiée (Largeur 12cm, Hauteur 8cm)
    \def\surfbase{0}
    \def\surftop{7}
    
    % --- SOCLE / PALEOZOIQUE (Plissé, tronqué) ---
    % On dessine des plis tronqués par la discordance à y=3.5
    \def\discy{3.5}
    
    % Couches paléozoïques (plissées)
    \foreach \i/\ybase in {1/0.3, 2/0.9, 3/1.5, 4/2.1, 5/2.7} {
        % Plis : sinusoïde
        \draw[paleozoic, thick, domain=0:12, smooth, variable=\x, samples=200] 
            plot ({\x}, {\ybase + 0.5*sin(150*\x)});
        % Remplissage sous la courbe jusqu'en bas (simplifié pour le visuel)
        % On ne remplit que la partie visible sous la discordance
        \clip (0,0) rectangle (12,\discy);
        \fill[paleozoic, opacity=0.4] plot[domain=0:12, smooth, variable=\x, samples=200] ({\x}, {\ybase + 0.5*sin(150*\x)}) -- (12,0) -- (0,0) -- cycle;
        \clip (0,0) rectangle (12,8); % reset clip
    }
    
    % SURFACE DE DISCORDANCE PRINCIPALE (Angulaire)
    \draw[discord] (0,\discy) -- (12,\discy);
    \node[labelstrat, anchor=south, rotate=0] at (6, \discy+0.1) {\textbf{Surface de discordance angulaire} (Érosion / Planation)};
    
    % CONGLOMERAT BASAL (mince)
    \fill[popgold!50!poporange, draw=popgold!80!black, thick] (0,\discy) rectangle (12,\discy+0.15);
    \node[font=\tiny, popdark] at (1, \discy+0.07) {Conglomérat basal};
    
    % --- MESOZOIQUE (Horizontal, transgressif) ---
    \foreach \y in {3.8, 4.3, 4.8, 5.3, 5.8, 6.3} {
        \draw[popblue!80!black, thick] (0,\y) -- (12,\y);
        \fill[popblue!15] (0,\y-0.25) rectangle (12,\y); % bandes
    }
    \node[labelstrat, anchor=south] at (6, 6.5) {\textbf{Série Mésozoïque} (Horizontalité, Transgression)};
    
    % --- FAILLE RECUPANT TOUT (Exemple pour attention) ---
    \draw[fault] (10, 0) -- (10.5, 8);
    \node[poppink, font=\small\bfseries, rotate=90, anchor=south] at (10.6, 4) {Faille normale (Postérieure)};
    
    % Légendes latérales
    \node[legendbox, anchor=north west] at (0.2, 7.8) {
        \textbf{Légende :}\par
        \tikz\fill[paleozoic, opacity=0.5] (0,0) rectangle (0.3,0.2); Paléozoïque (Plissé, tronqué)\par
        \tikz\fill[popgold!50!poporange] (0,0) rectangle (0.3,0.2); Conglomérat basal\par
        \tikz\fill[popblue!20] (0,0) rectangle (0.3,0.2); Mésozoïque (Horizontal)\par
        \tikz\draw[discord] (0,0.1) -- (0.3,0.1); Discordance\par
        \tikz\draw[fault] (0,0.1) -- (0.3,0.1); Faille
    };
    
    % Échelle et Nord
    \node[anchor=south east] at (11.8, 0.2) {\textbf{N}};
    \draw[|-|, thick] (9, 0.3) -- (11, 0.3) node[midway, above] {2 km};
    
    % Annotations chronologiques
    \draw[<-, thick, popdark] (0.5, 1.5) -- (-1, 1.5) node[left, labelstrat, anchor=east, align=center]{Ancien\\$\downarrow$};
    \draw[<-, thick, popdark] (0.5, 5.5) -- (-1, 5.5) node[left, labelstrat, anchor=east, align=center]{Récent\\$\uparrow$};
\end{tikzpicture}
\legende{Coupe géologique type montrant une discordance angulaire. Les couches paléozoïques (orangé, plissées) sont tronquées par la surface de discordance (trait zigzag épais). Un conglomérat basal (or) marque le rivage de la transgression mésozoïque. Les couches mésozoïques (bleu) sont horizontales et concordantes entre elles. Une faille normale (rose, pointillés) recoupe l'ensemble, elle est donc postérieure à la discordance.}
\end{popfigure}

\courssub{Valeur chronologique de la discordance : un repère temporel majeur}

La discordance n'est pas seulement une surface géométrique ; c'est un \textbf{repère chronostratigraphique} fondamental.

\begin{propriete}[Chronologie relative imposée par une discordance]
Soit une discordance séparant une \textbf{Série Inférieure} (SI) et une \textbf{Série Supérieure} (SS).
\begin{enumerate}
    \item Tous les événements affectant la SI (dépôt, diagenèse, plissement, métamorphisme, érosion) sont \textbf{antérieurs} à la mise en place de la surface de discordance.
    \item La surface de discordance elle-même représente l'intervalle de temps de l'érosion (lacune).
    \item Tous les événements affectant la SS (dépôt, diagenèse, déformation ultérieure) sont \textbf{postérieurs} à la surface de discordance.
    \item \textbf{Conséquence :} La discordance permet de \textbf{séparer l'histoire géologique en deux cycles distincts} (Cycle 1 = SI ; Cycle 2 = SS).
\end{enumerate}
\end{propriete}

\begin{exemplebox}[Ordre des événements autour d'une discordance : Cas pratique]
\textbf{Document :} Coupe schématique montrant :
\begin{itemize}
    \item Série A : Grès et schistes, plissés (plis couchés), tronqués au sommet.
    \item Surface S : Discordance angulaire avec conglomérat basal.
    \item Série B : Calcaires et marnes horizontaux, non déformés.
    \item Dyke (filon) doléritique recoupant la Série A et la Surface S, mais s'arrêtant à la base de la Série B (ou la recoupant ?).
\end{itemize}
\textbf{Question :} Établir la chronologie relative des événements (E1 à E6).
\textbf{Raisonnement guidé :}
\begin{enumerate}
    \item \textbf{Dépôt Série A} (Grès/Schistes) : Plus ancien (principe de superposition dans SI).
    \item \textbf{Plissement Série A} : Recoupe la stratification de A $\rightarrow$ Postérieur au dépôt A.
    \item \textbf{Érosion / Planation (Surface S)} : Tronque les plis de A $\rightarrow$ Postérieur au plissement.
    \item \textbf{Dépôt Conglomérat basal + Série B} : Reposent sur S, horizontaux $\rightarrow$ Postérieurs à l'érosion (Transgression).
    \item \textbf{Mise en place du Dyke} : Recoupe A et S. S'il recoupe B $\rightarrow$ Postérieur à B. S'il s'arrête sous B $\rightarrow$ Antérieur à B (synchrone de l'érosion ou juste avant). \textbf{Hypothèse standard :} Il recoupe A et S, mais pas B $\rightarrow$ Il est antérieur au dépôt de B (magmatisme lié à l'orogenèse ou à l'extension pré-transgressive).
\end{enumerate}
\textbf{Chronologie finale :} Dépôt A $<$ Plissement A $<$ Érosion (Surface S) $<$ \textbf{[Dyke ?]} $<$ Dépôt Conglomérat/B $<$ Éventuelle déformation finale.
\end{exemplebox}

\courssub{Exemples régionaux : Les discordances des bassins sédimentaires camerounais}

Le Cameroun offre des exemples textbook de discordances majeures, témoins de l'ouverture de l'Atlantique Sud et de l'histoire du Gondwana.

\begin{savaistu}[Discordances dans les bassins du Cameroun]
\begin{itemize}
    \item \textbf{Bassin de Douala (Côtier) :} La discordance majeure est la \textbf{discordance albienne (Crétacé inférieur)}. Elle sépare les séries \textbf{prérift / synrift} (Grès de Mundeck, Argiles de Logbaba - souvent continentaux, lacustres, déformés, faillés) des séries \textbf{post-rift} (Calcaires du Crétacé supérieur, Tertiaire - marins, sub-horizontaux). Cette discordance marque la \textbf{transgression marine albienne} majeure qui a envahi le rift en ouverture. On observe localement des discordances angulaires mineures liées à la tectonique halocinétique (mobilité du sel).
    \item \textbf{Bassin de Garoua (Intérieur, Bénoué) :} Discordance majeure à la base du \textbf{Crétacé} (Continental Intercalaire) sur le \textbf{Socle Précambrien} (granites, gneiss, schistes). C'est une discordance \textbf{non angulaire} (parallèle) mais à fort relief (paléorelief granitique érodé). Le socle est plissé/métamorphisé (cycle panafricain), érodé (surface de planation), puis recouvert par les grès fluviatiles crétacés. C'est une \textbf{discordance socle/couverture} classique.
    \item \textbf{Bassin du Logone-Birni (Extrême-Nord) :} Discordances multiples liées aux variations du Lac Mega-Chad (transgressions/régressions quaternaires) et à la tectonique distensive récente.
\end{itemize}
\end{savaistu}

\begin{exoresolu}[Analyse d'une coupe du bassin de Douala (simplifiée)]
\textbf{Énoncé :} On observe sur une coupe SE-NW du bassin de Douala (document fourni) :
\begin{itemize}
    \item Unité 1 (Base) : Alternances grès/argiles (Mundeck/Logbaba), fortement faillées (failles normales listriques), pendage variable $20-40^\circ$ vers le NW.
    \item Surface D : Surface plane, continue, recoupant les failles et les couches de l'Unité 1. Conglomérat basal présent par endroits.
    \item Unité 2 (Sommet) : Calcaires bioclastiques, marnes (Série Crétacé Supérieur/Tertiaire), pendage $< 5^\circ$ (sub-horizontaux), recouvrant la surface D en transgression.
    \item Un filon doléritique recoupe l'Unité 1 et la Surface D, mais est recouvert par l'Unité 2.
\end{itemize}
\textbf{Établir la chronologie relative des événements géologiques (E1 à E6).}
\tcblower
\textbf{Solution détaillée :}
\begin{enumerate}
    \item \textbf{E1 : Dépôt de l'Unité 1 (Séries prérift/synrift : Mundeck/Logbaba).} Âge : Jurassique supérieur - Crétacé inférieur (Néocomien-Barrémien). Milieu : Lacustre/Fluvial. Principe : Superposition (base la plus ancienne).
    \item \textbf{E2 : Tectonique extensive (Rifting) : Mise en place des failles normales listriques.} Ces failles recoupent la stratification de l'Unité 1 $\rightarrow$ Postérieures au dépôt E1. Elles contrôlent la subsidence et l'épaisseur des couches (croissance synsédimentaire possible).
    \item \textbf{E3 : Érosion / Planation (Mise en place de la Surface D - Discordance Albienne).} La surface D recoupe \textbf{à la fois} les couches de l'Unité 1 \textbf{et} les plans de faille de E2 $\rightarrow$ Elle est postérieure à la tectonique extensive. Le conglomérat basal témoigne d'une érosion continentale/fluviale avant l'arrivée de la mer.
    \item \textbf{E4 : Magmatisme : Mise en place du filon doléritique.} Le filon recoupe l'Unité 1 et la Surface D $\rightarrow$ Postérieur à E3. Mais il est recouvert par l'Unité 2 $\rightarrow$ Antérieur au dépôt de l'Unité 2. Ce magmatisme est souvent lié à l'ouverture de l'Atlantique (point chaud du Mont Cameroun / ligne du Cameroun) juste avant/après la transgression albienne.
    \item \textbf{E5 : Transgression marine Albienne et Dépôt de l'Unité 2 (Séries post-rift).} L'Unité 2 repose en concordance (parallélisme) sur la Surface D $\rightarrow$ Postérieure à E3 et E4. Milieu : Marin peu profond (calcaires) à plus profond (marnes). Horizontalité = fin de la tectonique extensive majeure.
    \item \textbf{E6 : Déformation mineure post-dépositionnelle (Plissement doux / Failles de croissance tertiaires).} L'Unité 2 montre un pendage très faible ($<5^\circ$) et quelques failles mineures $\rightarrow$ Tectonique post-dépositionnelle mineure (compaction, glissement gravitaire sur sel, ou compression alpine lointaine).
\end{enumerate}
\textbf{Chronologie synthétique :} E1 (Dépôt Synrift) $<$ E2 (Failles Normales) $<$ E3 (Érosion/Discordance Albienne) $<$ E4 (Filon Doléritique) $<$ E5 (Transgression/Dépôt Post-rift) $<$ E6 (Déformation Mineure).
\end{exoresolu}

\courssub{Autres types de discordances (pour aller plus loin)}

\begin{definition}[Typologie des discordances]
\begin{itemize}
    \item \textbf{Discordance angulaire (ou de pendage) :} Cas général traité ci-dessus. Pendage(SI) $\neq$ Pendage(SS). Témoin d'une déformation tectonique (plissement, basculement) entre les deux dépôts.
    \item \textbf{Discordance parallèle (ou disconformité) :} Pendage(SI) = Pendage(SS) (souvent horizontal). Aucune déformation tectonique visible. Seule une \textbf{érosion} (surface de planation, karst, paléosol) ou une \textbf{non-dépôt} (condensation) sépare les deux séries. Détectable seulement par une lacune paléontologique (fossiles manquants) ou un conglomérat basal/paléosol.
    \item \textbf{Non-conformité (ou discordance socle/couverture) :} La Série Inférieure est un \textbf{socle métamorphique ou magmatique} (non stratifié, ou schistosité non parallèle à la stratification). La Série Supérieure est sédimentaire. C'est le cas type du bassin de Garoua sur le socle précambrien.
    \item \textbf{Paraconformité :} Discordance parallèle sans érosion visible ni conglomérat, détectée uniquement par une lacune biostratigraphique (zone à fossiles manquants). Difficile à voir sur le terrain sans paléontologie.
\end{itemize}
\end{definition}

\begin{attention}[Pièges classiques au Probatoire]
\begin{enumerate}
    \item \textbf{Inverser l'ordre :} Dire « les couches du haut sont plus anciennes car elles sont pliées ». \textbf{Faux.} Le pendage ne donne pas l'âge. C'est la position relative (superposition) et la troncature qui donnent l'âge. Les couches pliées sont \emph{sous} la discordance, donc plus anciennes.
    \item \textbf{Confondre faille et discordance :} Voir boîte \textbf{Attention} plus haut. Vérifier : y a-t-il troncature des couches inférieures ? Les couches supérieures sont-elles parallèles à la surface de contact ? $\rightarrow$ Discordance. Y a-t-il rejet de couches de même âge ? $\rightarrow$ Faille.
    \item \textbf{Oublier la lacune :} Une discordance représente du \textbf{temps manquant}. Ne pas écrire « dépôt continu ». Préciser : « Érosion de X millions d'années non enregistrée ».
    \item \textbf{Négliger le conglomérat basal :} C'est la preuve terrain de l'érosion et de la transgression. Toujours le mentionner s'il est présent sur le document.
\end{enumerate}
\end{attention}

\begin{exercice}[Application : Lecture de coupe avec discordance multiple]
\textbf{Document :} Coupe schématique (non fournie ici, à imaginer ou dessiner) montrant de la base au sommet :
1. Granite (Socle).
2. Surface D1 (irrégulière, érodée).
3. Grès rouge (Buntsandstein), horizontal.
4. Calcaire (Muschelkalk), horizontal.
5. Surface D2 (plane, recoupant le calcaire).
6. Argiles à gypses (Keuper), horizontal.
7. Faille normale F recoupant tout (1 à 6) avec rejet de 50 m.
\textbf{Questions :}
\begin{enumerate}
    \item Qualifier les surfaces D1 et D2 (type de discordance).
    \item Établir la chronologie relative complète (7 événements minimum).
    \item Quel événement est le plus récent ? Justifier.
\end{enumerate}
\end{exercice}

\begin{corrige}[Exercice : Lecture de coupe avec discordance multiple]
\begin{enumerate}
    \item \textbf{D1 : Non-conformité (Discordance socle/couverture).} Socle granitique (non stratifié) recouvert par des sédiments horizontaux. \textbf{D2 : Disconformité (Discordance parallèle).} Calcaire et Argiles ont même pendage (horizontal). D2 est une surface d'érosion/planation (éventuellement karstique sur calcaire) sans déformation tectonique intermédiaire.
    \item \textbf{Chronologie :}
    \begin{enumerate}
        \item Mise en place du Granite (refroidissement/cristallisation).
        \item Érosion du Granite $\rightarrow$ Surface D1 (Pénéplaine).
        \item Transgression / Dépôt Grès (Buntsandstein).
        \item Dépôt Calcaire (Muschelkalk).
        \item Érosion / Régression $\rightarrow$ Surface D2 (Lacune : partie supérieure Muschelkalk + éventuel Ladinien manquant).
        \item Transgression / Dépôt Argiles/Gypses (Keuper).
        \item Tectonique extensive : Faille normale F (recoupe tout, donc postérieur au Keuper).
    \end{enumerate}
    \item \textbf{Événement le plus récent : La Faille F.} Elle recoupe l'Unité 6 (Keuper) qui est la plus jeune formation sédimentaire. Principe de recoupement : ce qui coupe est plus jeune que ce qui est coupé.
\end{enumerate}
\end{corrige}

\begin{aretenir}[Synthèse : La discordance, marqueur de cycle géologique]
\begin{itemize}
    \item Une \textbf{discordance} est une surface d'érosion séparant deux séries de pendages différents (angulaire) ou identiques (parallèle/disconformité).
    \item Sa formation suit le cycle : \textbf{Dépôt $\rightarrow$ Déformation (optionnel) $\rightarrow$ Émersion/Érosion $\rightarrow$ Transgression $\rightarrow$ Nouveau Dépôt}.
    \item Sur une coupe : \textbf{Troncature} des couches inférieures + \textbf{Surface irrégulière/zigzag} + \textbf{Parallélisme} des couches supérieures à cette surface + \textbf{Conglomérat basal} fréquent.
    \item Elle définit une \textbf{lacune stratigraphique} (temps non enregistré) et sépare l'histoire en \textbf{deux cycles tectono-sédimentaires} distincts.
    \item Au Cameroun : Discordance albienne (Douala, transgression post-rift) et Discordance socle/couverture Crétacé (Garoua, Logone).
\end{itemize}
\end{aretenir}

\coursec{Les fossiles stratigraphiques}

Dans la section précédente, nous avons vu que les \cle{discordances stratigraphiques} sont des surfaces témoins d'une interruption de la sédimentation : elles permettent de repérer, dans une coupe, les grandes étapes de l'histoire d'une région (sédimentation, déformation, érosion, nouvelle sédimentation). Mais les discordances, comme les principes de superposition et de recoupement, ne permettent d'ordonner les événements qu'à l'intérieur d'une même région. Comment savoir si une couche de calcaire observée près de Foumban est du même âge qu'une couche de grès observée dans l'Extrême-Nord, ou même en Europe ? Les principes de superposition et de continuité sont impuissants à relier des affleurements éloignés, séparés par des milliers de kilomètres. C'est ici qu'intervient un outil remarquable : les \cle{fossiles stratigraphiques}.

\courssub{Qu'est-ce qu'un fossile ? Rappels}

\begin{definition}[Fossile]
Un \cle{fossile} est un reste (coquille, os, dent, feuille, tronc\ldots) ou une trace d'activité (empreinte de pas, galerie, coprolithe, moule\ldots) d'un organisme ayant vécu autrefois, conservé dans une roche, le plus souvent une \cle{roche sédimentaire}.
\end{definition}

La fossilisation est un événement rare : pour qu'un organisme se fossilise, il faut généralement qu'il possède des parties dures (coquille, squelette) et qu'il soit enfoui rapidement après sa mort dans un sédiment (vase, sable, cendre volcanique), à l'abri de l'oxygène et des décomposeurs. Au fil du temps, le sédiment se transforme en roche (diagénèse) et l'organisme, ou son empreinte, est conservé. C'est pourquoi l'immense majorité des fossiles se trouve dans les roches sédimentaires, et presque jamais dans les roches magmatiques ou métamorphiques, dont la formation détruit toute matière organique.

\begin{savaistu}[Les fossiles, archives de la vie]
Le mot « fossile » vient du latin \emph{fossilis}, « qu'on a déterré ». Au XVIII\textsuperscript{e} siècle, le naturaliste Georges Cuvier montra que les fossiles sont les restes d'espèces aujourd'hui disparues, et non de simples « jeux de la nature ». Depuis, les paléontologues ont décrit des centaines de milliers d'espèces fossiles, qui retracent plus de 3 milliards d'années d'histoire de la vie sur Terre.
\end{savaistu}

\courssub{Du fossile au fossile stratigraphique : l'intuition}

Imaginons la situation suivante, très concrète. Dans une carrière de la région de l'Ouest, un élève découvre dans une couche de calcaire une coquille enroulée en spirale : une \cle{ammonite}. Son professeur lui affirme : « Cette couche date du Secondaire, elle a environ 150 millions d'années. » Comment peut-il être aussi affirmatif, sans aucun appareil de mesure ?

L'idée directrice est simple : les espèces vivantes \textbf{apparaissent, évoluent puis s'éteignent}. Chaque espèce n'a donc existé que pendant un intervalle de temps limité de l'histoire de la Terre. Si l'on connaît cet intervalle, la présence de l'espèce dans une couche permet de situer l'âge de cette couche dans cet intervalle. Le fossile devient alors un véritable « marqueur temporel ».

\begin{definition}[Fossile stratigraphique]
Un \cle{fossile stratigraphique} (ou \cle{fossile index}, ou fossile caractéristique) est un fossile qui permet de \textbf{dater relativement} la couche qui le contient, c'est-à-dire de l'attribuer à une ère, une période ou un étage de l'échelle des temps géologiques.
\end{definition}

\courssub{Les trois critères d'un bon fossile stratigraphique}

Tous les fossiles ne sont pas de bons marqueurs temporels. Pour être utile à la datation relative, un fossile stratigraphique doit réunir \textbf{trois critères} :

\begin{definition}[Critères d'un bon fossile stratigraphique]
Un bon fossile stratigraphique présente :
\begin{enumerate}
\item une \cle{large répartition géographique} : l'espèce a vécu sur de vastes étendues (idéalement sur tous les continents), ce qui permet de corréler des couches très éloignées les unes des autres ;
\item une \cle{courte durée d'existence} : l'espèce est apparue et s'est éteinte rapidement (évolution rapide), ce qui définit un intervalle de temps étroit et donc une datation précise ;
\item une \cle{grande abondance} : l'espèce a laissé de nombreux fossiles, faciles à trouver dans les couches, ce qui rend son utilisation pratique.
\end{enumerate}
\end{definition}

On peut résumer ces critères par une formule à retenir : un bon fossile stratigraphique est \textbf{répandu, éphémère et abondant}.

\begin{attention}[Tous les fossiles ne datent pas !]
Un fossile à \textbf{longue durée de vie} n'est \textbf{pas} un bon fossile stratigraphique. Par exemple, certains coquillages (huîtres, moules, lingules) ont des représentants quasi identiques depuis des centaines de millions d'années jusqu'à aujourd'hui : trouver une telle coquille dans une couche ne renseigne presque pas sur son âge. De même, un fossile très rare ou limité à une petite région ne permet pas de corréler des couches éloignées. Ne confonds jamais « fossile » et « fossile stratigraphique » : tout fossile stratigraphique est un fossile, mais la réciproque est fausse.
\end{attention}

\courssub{Le principe d'identité paléontologique}

L'utilisation des fossiles stratigraphiques repose sur un principe fondamental, établi au XIX\textsuperscript{e} siècle notamment par l'ingénieur anglais William Smith, surnommé « le père de la stratigraphie ».

\begin{propriete}[Principe d'identité paléontologique]
Deux couches sédimentaires contenant les \textbf{mêmes fossiles stratigraphiques} sont considérées comme étant du \textbf{même âge}, même si elles sont très éloignées géographiquement et même si leur nature (lithologie) est différente.
\end{propriete}

Ce principe est \textbf{admis} au niveau du Probatoire (sa démonstration n'est pas exigible), mais il repose sur un fondement solide : \textbf{l'évolution des espèces est irréversible}. Une espèce éteinte ne réapparaît jamais ; une association d'espèces donnée n'a donc existé qu'une seule fois dans l'histoire de la Terre. Retrouver la même association fossilifère dans deux couches, c'est donc retrouver le même « instant » de l'histoire de la vie.

\begin{experience}[Démarche : corréler deux coupes éloignées]
\textbf{Situation :} on étudie deux affleurements distants de plusieurs centaines de kilomètres. La coupe A présente des couches de calcaire ; la coupe B, des couches de grès et d'argile. Peut-on dire si certaines couches sont contemporaines ?

\textbf{Observation :} dans la couche A2 (calcaire) et dans la couche B3 (grès), on découvre le même fossile : une ammonite de l'espèce \emph{Perisphinctes}.

\textbf{Hypothèse :} si le principe d'identité paléontologique s'applique, A2 et B3 sont du même âge.

\textbf{Vérification :} on recherche d'autres fossiles stratigraphiques. A2 et B3 contiennent aussi le même foraminifère, tandis que les couches voisines (A1, A3, B2, B4) contiennent des associations fossilifères différentes, cohérentes avec l'ordre de superposition.

\textbf{Interprétation :} la double coïncidence fossilifère confirme l'hypothèse.

\textbf{Conclusion :} A2 et B3 sont contemporaines, bien que de lithologies différentes : elles se sont déposées à la même époque, mais dans des milieux de dépôt différents (mer calme et profonde pour le calcaire, milieu côtier agité pour le grès).
\end{experience}

\begin{popfigure}
\begin{tikzpicture}[scale=0.92]
% ===== Coupe A =====
\begin{scope}
\fill[popgoldL] (0,0) rectangle (4,1);
\fill[popmuted!40] (0,1) rectangle (4,2.2);
\fill[popcream] (0,2.2) rectangle (4,3.4);
\draw[thick,popdark] (0,0) rectangle (4,3.4);
\draw[popdark] (0,1)--(4,1) (0,2.2)--(4,2.2);
\node[font=\small] at (2,0.5) {A1 : argile};
\node[font=\small] at (2,1.6) {A2 : calcaire};
\node[font=\small] at (2,2.8) {A3 : calcaire};
\node[font=\small\bfseries,popblue] at (2,-0.5) {Coupe A (Ouest)};
% ammonite stylisée dans A2
\begin{scope}[shift={(1.1,1.6)},scale=0.28]
\draw[thick,popink] (0,0) circle (1);
\draw[thick,popink] (0,0) circle (0.62);
\draw[thick,popink] (0,0) circle (0.3);
\draw[thick,popink] (1,0) arc (0:-80:1) -- (0.35,-0.85);
\end{scope}
\node[font=\scriptsize,popink] at (2.6,1.35) {ammonite \emph{Perisphinctes}};
\end{scope}
% ===== Coupe B =====
\begin{scope}[shift={(8.5,0)}]
\fill[popgreenL] (0,0) rectangle (4,0.9);
\fill[poporangeL] (0,0.9) rectangle (4,1.7);
\fill[popmuted!40] (0,1.7) rectangle (4,2.9);
\fill[popgoldL] (0,2.9) rectangle (4,3.4);
\draw[thick,popdark] (0,0) rectangle (4,3.4);
\draw[popdark] (0,0.9)--(4,0.9) (0,1.7)--(4,1.7) (0,2.9)--(4,2.9);
\node[font=\small] at (2,0.45) {B1 : argile};
\node[font=\small] at (2,1.3) {B2 : sable};
\node[font=\small] at (2,2.3) {B3 : gr\`es};
\node[font=\small] at (2,3.15) {B4 : argile};
\node[font=\small\bfseries,popblue] at (2,-0.5) {Coupe B (Extr\^eme-Nord)};
% ammonite stylisée dans B3
\begin{scope}[shift={(1.1,2.3)},scale=0.28]
\draw[thick,popink] (0,0) circle (1);
\draw[thick,popink] (0,0) circle (0.62);
\draw[thick,popink] (0,0) circle (0.3);
\draw[thick,popink] (1,0) arc (0:-80:1) -- (0.35,-0.85);
\end{scope}
\node[font=\scriptsize,popink] at (2.7,2.05) {ammonite \emph{Perisphinctes}};
\end{scope}
% flèche de corrélation
\draw[very thick,popink,dashed,-{Stealth}] (4.1,1.6) .. controls (5.6,2.9) and (7.4,2.9) .. (8.4,2.3);
\node[font=\small\itshape,popink] at (6.25,3.15) {m\^eme \^age};
\end{tikzpicture}
\legende{Corrélation de deux coupes éloignées grâce à un fossile stratigraphique. Les couches A2 (calcaire) et B3 (grès) ont des lithologies différentes mais contiennent la même ammonite : d'après le principe d'identité paléontologique, elles sont du même âge. Les milieux de dépôt étaient simplement différents d'une région à l'autre.}
\end{popfigure}

\begin{attention}[Ne pas corréler sur la seule lithologie]
Deux couches de même nature (par exemple deux couches de sable) ne sont \textbf{pas nécessairement} du même âge : on trouve du sable à presque toutes les époques ! Inversement, deux couches de natures différentes peuvent être contemporaines. Seul le contenu fossilifère (ou, en l'absence de fossiles, la continuité physique des couches) permet une corrélation fiable.
\end{attention}

\courssub{Les principaux fossiles stratigraphiques et l'échelle des temps géologiques}

Certains groupes d'organismes fossiles sont devenus les « vedettes » de la datation relative, car ils remplissent parfaitement les trois critères :

\begin{itemize}
\item les \cle{trilobites}, arthropodes marins disparus, abondants et très diversifiés au \textbf{Primaire} (Paléozoïque) ; ils s'éteignent à la fin de cette ère ;
\item les \cle{graptolites}, organismes coloniaux flottants, excellents marqueurs du début du \textbf{Primaire} (Ordovicien-Silurien) ;
\item les \cle{ammonites}, mollusques céphalopodes à coquille enroulée, très abondantes et à évolution rapide au \textbf{Secondaire} (Mésozoïque) ; elles disparaissent il y a 66 millions d'années, en même temps que les dinosaures ;
\item les \cle{bélemnites}, céphalopodes à rostre interne, caractéristiques du \textbf{Secondaire} ;
\item les \cle{rudistes}, bivalves sessiles constructeurs de récifs, marqueurs du Crétacé (fin du Secondaire) ;
\item les \cle{foraminifères}, protozoaires marins à coquille calcaire microscopique, utiles du Primaire au \textbf{Quaternaire}, notamment les \cle{nummulites} au \textbf{Tertiaire} (Paléogène) ;
\item les \cle{pollens et spores}, très abondants, utilisés surtout pour le \textbf{Quaternaire} et le Tertiaire supérieur.
\end{itemize}

\begin{popfigure}
\begin{tikzpicture}[scale=1]
% Tableau TikZ : fossiles stratigraphiques et ères
\fill[popblueL] (0,4.8) rectangle (12,5.6);
\node[font=\small\bfseries,popdark] at (2,5.2) {\`Ere};
\node[font=\small\bfseries,popdark] at (5.6,5.2) {Dur\'ee (Ma)};
\node[font=\small\bfseries,popdark] at (9.2,5.2) {Fossiles stratigraphiques typiques};
\draw[popdark] (0,4.8) rectangle (12,5.6);
% Ligne Primaire
\fill[popcream] (0,3.6) rectangle (12,4.8);
\node[font=\small\bfseries,popdark] at (2,4.2) {Primaire (Pal\'eozo\"ique)};
\node[font=\small] at (5.6,4.2) {540 \`a 250};
\node[font=\small,align=center] at (9.2,4.2) {trilobites, graptolites,\\ brachiopodes};
\draw[popdark] (0,3.6) rectangle (12,4.8);
% Ligne Secondaire
\fill[poporangeL] (0,2.4) rectangle (12,3.6);
\node[font=\small\bfseries,popdark] at (2,3) {Secondaire (M\'esozo\"ique)};
\node[font=\small] at (5.6,3) {250 \`a 66};
\node[font=\small,align=center] at (9.2,3) {ammonites, b\'elemnites,\\ rudistes};
\draw[popdark] (0,2.4) rectangle (12,3.6);
% Ligne Tertiaire
\fill[popgreenL] (0,1.2) rectangle (12,2.4);
\node[font=\small\bfseries,popdark] at (2,1.8) {Tertiaire (C\'enozo\"ique)}; % Programme MINESEC utilise encore Tertiaire/Quaternaire
\node[font=\small] at (5.6,1.8) {66 \`a 1,6};
\node[font=\small,align=center] at (9.2,1.8) {nummulites (foraminif\`eres),\\ mammifères};
\draw[popdark] (0,1.2) rectangle (12,2.4);
% Ligne Quaternaire
\fill[poppinkL] (0,0) rectangle (12,1.2);
\node[font=\small\bfseries,popdark] at (2,0.6) {Quaternaire};
\node[font=\small] at (5.6,0.6) {1,6 \`a 0};
\node[font=\small,align=center] at (9.2,0.6) {pollens, spores,\\ grands mammifères};
\draw[popdark] (0,0) rectangle (12,1.2);
% séparations verticales
\draw[popdark] (3.4,0)--(3.4,5.6) (7.2,0)--(7.2,5.6);
\end{tikzpicture}
\legende{Principaux fossiles stratigraphiques associés aux grandes ères de l'échelle des temps géologiques (Ma = millions d'années). La présence de l'un de ces fossiles dans une couche permet d'attribuer cette couche à l'ère correspondante.}
\end{popfigure}

\begin{exemplebox}[Les trilobites, marqueurs du Primaire]
Un géologue découvre des trilobites dans une couche de schiste. Or les trilobites n'ont vécu qu'au Primaire (entre environ 540 et 250 millions d'années) et se sont totalement éteints à la fin de cette ère (crise Permien-Trias). La couche de schiste peut donc être attribuée sans hésiter à l'\textbf{ère Primaire}. Si l'espèce de trilobite est connue avec précision, on peut même affiner la datation à une période (Cambrien, Ordovicien\ldots).
\end{exemplebox}

\begin{savaistu}[Les ammonites, horloges du Secondaire]
Les ammonites ont vécu dans les mers du monde entier pendant plus de 300 millions d'années, avant de disparaître il y a 66 millions d'années, lors de la crise Crétacé-Paléogène qui a aussi emporté les dinosaures. Leur évolution était si rapide que chaque espèce n'a souvent vécu qu'un à deux millions d'années : ce sont de véritables « horloges » qui permettent de dater et de corréler des couches sur \textbf{tous les continents}. Une ammonite trouvée dans une falaise d'Angleterre peut ainsi certifier qu'une couche d'Amérique du Sud ou d'Afrique est exactement du même âge. On définit ainsi des \emph{biozones} à ammonites, unités de base de l'échelle stratigraphique du Mésozoïque.
\end{savaistu}

\courssub{Application : attribuer une formation à une ère ou une période}

En pratique, le raisonnement du géologue (et celui qu'on attend de toi au Probatoire) suit toujours la même logique en trois étapes.

\begin{methode}[Dater une couche grâce à son contenu fossilifère]
\begin{enumerate}
\item \textbf{Identifier} les fossiles présents dans la couche (nom du groupe, voire de l'espèce).
\item \textbf{Confronter} ces fossiles à leurs intervalles d'existence connus dans l'échelle des temps géologiques : retenir l'intervalle \emph{commun} à tous les fossiles trouvés (le plus restrictif, intersection des intervalles).
\item \textbf{Conclure} en attribuant la couche à l'ère, la période ou l'étage correspondant, en justifiant par le principe d'identité paléontologique.
\end{enumerate}
\end{methode}

\begin{exoresolu}[Datation d'une formation de la falaise de Limbe]
\textbf{Énoncé.} Dans une falaise littorale de la région du Sud-Ouest, on relève de bas en haut trois couches sédimentaires : C1 (grès sans fossile), C2 (calcaire contenant des ammonites et des bélemnites), C3 (sable contenant des nummulites, foraminifères caractéristiques du Tertiaire). Attribuer chaque couche fossilifère à une ère et préciser la chronologie relative des dépôts.

\tcblower
\textbf{Solution détaillée.}
\begin{itemize}
\item \textbf{Couche C2} : les ammonites et les bélemnites sont des fossiles stratigraphiques caractéristiques de l'ère \textbf{Secondaire} (Mésozoïque, 250 à 66 Ma), où ils étaient abondants, répandus et à évolution rapide. La couche C2 est donc d'âge Secondaire.
\item \textbf{Couche C3} : les nummulites sont des foraminifères caractéristiques de l'ère \textbf{Tertiaire} (Paléogène, Cénozoïque). La couche C3 est donc d'âge Tertiaire.
\item \textbf{Couche C1} : dépourvue de fossiles, elle ne peut pas être datée directement. Mais d'après le \textbf{principe de superposition}, C1 est sous C2, donc C1 est \emph{antérieure} au Secondaire : elle est d'âge Primaire ou plus ancienne.
\end{itemize}
\textbf{Chronologie relative} (du plus ancien au plus récent) : dépôt de C1 (avant le Secondaire) $\rightarrow$ dépôt de C2 (Secondaire) $\rightarrow$ dépôt de C3 (Tertiaire). Cette succession est cohérente : les couches deviennent plus récentes vers le haut, et le contenu fossilifère confirme l'ordre des ères.
\end{exoresolu}

\begin{aretenir}[L'essentiel sur les fossiles stratigraphiques]
\begin{itemize}
\item Un \cle{fossile} est un reste ou une trace d'organisme conservé dans une roche sédimentaire.
\item Un \cle{fossile stratigraphique} permet de dater relativement la couche qui le contient ; il doit être \textbf{répandu géographiquement, à courte durée d'existence et abondant}.
\item \textbf{Principe d'identité paléontologique} (admis) : deux couches contenant les mêmes fossiles stratigraphiques sont du même âge, quelles que soient leur distance et leur lithologie. Il repose sur l'évolution irréversible des espèces.
\item Fossiles repères : \textbf{trilobites et graptolites} (Primaire), \textbf{ammonites, bélemnites et rudistes} (Secondaire), \textbf{nummulites} (Tertiaire), \textbf{pollens et spores} (Quaternaire).
\item Les fossiles permettent la \textbf{corrélation} de couches éloignées et l'attribution d'une formation à une ère ou une période de l'échelle des temps géologiques.
\end{itemize}
\end{aretenir}

\begin{exercice}[Corrélation de deux régions]
On étudie deux coupes distantes de 800 km. Coupe X : X1 (schiste à trilobites), X2 (calcaire à ammonites). Coupe Y : Y1 (argile sans fossile), Y2 (grès à ammonites), Y3 (sable à nummulites).
\begin{enumerate}
\item Attribuer chaque couche fossilifère à une ère en justifiant.
\item Quelles couches des deux coupes sont contemporaines ? Justifier à l'aide d'un principe.
\item Peut-on dater la couche Y1 ? Que peut-on dire de son âge ?
\end{enumerate}
\end{exercice}

\begin{corrige}[Exercice : corrélation de deux régions]
\begin{enumerate}
\item X1 contient des trilobites, fossiles stratigraphiques du \textbf{Primaire} : X1 est d'âge Primaire. X2 et Y2 contiennent des ammonites, fossiles stratigraphiques du \textbf{Secondaire} : elles sont d'âge Secondaire. Y3 contient des nummulites, foraminifères caractéristiques du \textbf{Tertiaire} : Y3 est d'âge Tertiaire.
\item X2 et Y2 sont \textbf{contemporaines} : elles contiennent les mêmes fossiles stratigraphiques (ammonites) ; d'après le \textbf{principe d'identité paléontologique}, elles sont du même âge malgré leur éloignement et leurs lithologies différentes (calcaire contre grès).
\item Y1 ne contient aucun fossile : on ne peut pas la dater directement. Cependant, par le \textbf{principe de superposition}, Y1 est sous Y2 (Secondaire) : Y1 est donc \emph{antérieure} au Secondaire, c'est-à-dire d'âge Primaire ou plus ancien.
\end{enumerate}
\end{corrige}

Ainsi, les fossiles stratigraphiques complètent les principes de la stratigraphie : là où la superposition ordonne les couches d'une même région, les fossiles relient entre elles les régions du monde entier. Il ne reste plus qu'à assembler tous ces outils — superposition, recoupement, discordances, fossiles — pour reconstituer la chronologie complète des événements géologiques d'une région : ce sera l'objet de la dernière section de cette leçon.

\coursec{Reconstitution de la chronologie des événements géologiques}

Après avoir identifié les couches sédimentaires, compris les principes fondamentaux de la stratigraphie (superposition, continuité, horizontalité originelle, recoupement), reconnu les discordances et appris à utiliser les fossiles stratigraphiques comme des « horloges » géologiques, nous possédons maintenant tous les outils pour mener l'enquête ultime : \textbf{reconstituer l'histoire complète d'une région}. C'est l'aboutissement du raisonnement géologique : transformer une coupe statique en un film dynamique, ordonner les événements du plus ancien au plus récent, et justifier rigoureusement chaque étape. C'est exactement ce que l'on attend de vous au Probatoire, et c'est la méthode qu'utilisent les géologues de la SNH (Société Nationale des Hydrocarbures) pour reconstruire l'histoire du bassin de Douala et guider l'exploration pétrolière au Cameroun.

\courssub{Démarche générale : les 4 étapes de l'enquête géologique}

Face à une coupe géologique (schéma, photo aérienne ou affleurement réel), ne vous précipitez pas. La rigueur scientifique impose une démarche en quatre temps, que vous devez appliquer systématiquement.

\begin{methode}[Les 4 étapes pour établir la chronologie des événements géologiques d'une région]
\begin{enumerate}
    \item \textbf{Observer et légender la coupe.} Repérer toutes les couches (lithologie : calcaire, grès, schiste, etc.), leur attitude (horizontales, inclinées, verticales, plissées), les structures tectoniques (plis, failles), les discordances (surfaces d'érosion) et les éventuelles intrusions magmatiques. Légender chaque élément avec un code clair (ex. : C1, C2, P1 pour un pli, F1 pour une faille, D1 pour une discordance).
    \item \textbf{Identifier les événements élémentaires.} Chaque couche correspond à un \textbf{dépôt} (sédimentation). Chaque pli correspond à un \textbf{plissement} (déformation compressive). Chaque faille correspond à une \textbf{fracturation} (mise en place d'une faille, normale ou inverse). Chaque discordance signale un épisode d'\textbf{émersion + érosion} (souvent suivi d'une \textbf{transgression} marine pour les couches sus-jacentes). Chaque dyke ou filon-couche signale une \textbf{intrusion magmatique}.
    \item \textbf{Appliquer les principes de datation relative.} Pour chaque couple d'événements, déterminer lequel est antérieur à l'autre en invoquant explicitement le principe utilisé :
    \begin{itemize}
        \item \textbf{Principe de superposition :} dans une série non renversée, la couche du bas est plus ancienne que celle du haut.
        \item \textbf{Principe de recoupement :} une structure (faille, filon, discordance) qui en recoupe une autre est plus jeune que la structure recoupée.
        \item \textbf{Principe d'identité paléontologique :} deux couches contenant le même fossile stratigraphique sont contemporaines (isochrones).
    \end{itemize}
    \item \textbf{Ordonner les événements du plus ancien au plus récent.} Rédiger la chronologie sous forme de liste numérotée ou de frise verticale, en justifiant \emph{chaque} relation d'ordre par le principe invoqué. Vérifier qu'aucun événement n'a été oublié (surtout l'érosion !).
\end{enumerate}
\end{methode}

\begin{attention}[Piège fatal : l'épisode d'érosion « invisible »]
C'est l'erreur n°1 au Probatoire. Devant une discordance angulaire, beaucoup d'élèves écrivent : « 1. Dépôt des couches inférieures. 2. Plissement. 3. Dépôt des couches supérieures. » \textbf{FAUX.} Entre le plissement (qui déforme les couches) et le dépôt des couches horizontales sus-jacentes, il s'est passé \textbf{deux} événements majeurs : \textbf{l'émersion du secteur (soulèvement) suivie d'une érosion intense} qui a tronqué le pli (créant la surface plane de la discordance), puis \textbf{une transgression marine} (envahissement par la mer) permettant le nouveau dépôt. \textbf{Oublier l'érosion, c'est perdre la moitié de l'histoire géologique et des points faciles.}
\end{attention}

\courssub{Cas 1 : Série sédimentaire simple non perturbée (application de la superposition)}

Commençons par le cas le plus simple, socle de tout raisonnement ultérieur.

\begin{exemplebox}[Coupe du Mont Cameroun : série volcanico-sédimentaire récente]
Observons une coupe simplifiée du flanc Sud-Ouest du Mont Cameroun (volcan actif), montrant une alternance de coulées de lave basaltique et de couches de cendres (tufs) non déformées, reposant sur le socle métamorphique ancien.
\end{exemplebox}

\begin{popfigure}
\begin{tikzpicture}[scale=0.9]
% Socle
\fill[poppurple!30] (0,0) rectangle (10,1.5);
\node[align=center] at (5,0.75) {\textbf{Socle métamorphique}\\(gneiss, micaschistes)\\(Précambrien)};
% Couche 1 : Lave ancienne
\fill[poporange!50] (0,1.5) rectangle (10,2.5);
\node at (5,2) {\textbf{L1 : coulée basaltique ancienne}};
% Couche 2 : Cendres
\fill[popgold!50] (0,2.5) rectangle (10,3.3);
\node at (5,2.9) {\textbf{C1 : niveau de cendres (tuf)}};
% Couche 3 : Lave intermédiaire
\fill[poporange!70] (0,3.3) rectangle (10,4.3);
\node at (5,3.8) {\textbf{L2 : coulée basaltique intermédiaire}};
% Couche 4 : Cendres récentes
\fill[popgold!70] (0,4.3) rectangle (10,5.1);
\node at (5,4.7) {\textbf{C2 : niveau de cendres récentes}};
% Couche 5 : Lave très récente (1999/2000)
\fill[popink!50] (0,5.1) rectangle (10,6);
\node at (5,5.55) {\textbf{L3 : coulée de 1999-2000 (historique)}};
% Surface actuelle
\draw[thick, popgreen!80!black, decorate, decoration={snake, amplitude=1pt, segment length=4pt}] (0,6) -- (10,6);
\node[popgreen!80!black] at (5,6.4) {Surface topographique actuelle (sol, végétation)};
% Légende gauche
\draw[<->, thick] (-0.5,0) -- (-0.5,6);
\node[rotate=90] at (-1.2,3) {Profondeur / altitude relative};
% Flèches d'âge
\draw[->, thick, popdark] (10.5, 5.5) -- (10.5, 0.5);
\node[rotate=-90, align=left, popdark] at (11.3, 3) {Plus récent $\rightarrow$ Plus ancien};
\end{tikzpicture}
\legende{Coupe schématique simplifiée du flanc Sud-Ouest du Mont Cameroun. Série non déformée, horizontale : application directe du principe de superposition.}
\end{popfigure}

\textbf{Chronologie établie (du plus ancien au plus récent) :}
\begin{enumerate}
    \item \textbf{Mise en place du socle métamorphique} (Précambrien). \textit{Justification : il constitue la base de la coupe et supporte toutes les formations sus-jacentes (superposition).}
    \item \textbf{Dépôt de la coulée L1.} \textit{Justification : elle repose directement sur le socle (superposition).}
    \item \textbf{Dépôt du niveau de cendres C1.} \textit{Justification : il recouvre L1 (superposition).}
    \item \textbf{Dépôt de la coulée L2.} \textit{Justification : elle recouvre C1 (superposition).}
    \item \textbf{Dépôt du niveau de cendres C2.} \textit{Justification : il recouvre L2 (superposition).}
    \item \textbf{Dépôt de la coulée L3 (1999-2000).} \textit{Justification : elle est sommitale, historique, et recouvre tout (superposition).}
\end{enumerate}
Ici, un seul principe suffit : la \textbf{superposition}. Aucune déformation, aucune faille, aucune discordance n'intervient. L'histoire est une simple pile de « crêpes » géologiques.

\courssub{Cas 2 : Série plissée puis recoupée par une faille (le piège du recoupement)}

La complexité arrive avec la tectonique. Un pli se forme \emph{après} le dépôt des couches qu'il déforme. Une faille se forme \emph{après} les roches qu'elle décale. Mais quelle est la relation d'âge entre le pli et la faille ? C'est là que le \textbf{principe de recoupement} devient l'outil décisif.

\begin{experience}[Analyse d'une coupe plissée faillée : « Le pli est-il plus vieux que la faille ? »]
\textbf{Document :} coupe géologique schématique (type examen Probatoire) montrant un anticlinal et un synclinal affectant des couches C1 à C4, le tout recoupé par une faille normale F1 qui décale les couches.

\begin{popfigure}
\begin{tikzpicture}[scale=0.85]
% Couches plissées (bloc gauche)
\draw[thick] (0,0) .. controls (1,0.2) and (3,0.2) .. (4,0) .. controls (5,-0.2) and (7,-0.2) .. (8,0) .. controls (9,0.2) and (11,0.2) .. (12,0);
\draw[thick] (0,1) .. controls (1,1.2) and (3,1.2) .. (4,1) .. controls (5,0.8) and (7,0.8) .. (8,1) .. controls (9,1.2) and (11,1.2) .. (12,1);
\node at (6,0.5) {\textbf{C1}};
\draw[thick] (0,2) .. controls (1,2.2) and (3,2.2) .. (4,2) .. controls (5,1.8) and (7,1.8) .. (8,2) .. controls (9,2.2) and (11,2.2) .. (12,2);
\node at (6,1.5) {\textbf{C2}};
\draw[thick] (0,3) .. controls (1,3.2) and (3,3.2) .. (4,3) .. controls (5,2.8) and (7,2.8) .. (8,3) .. controls (9,3.2) and (11,3.2) .. (12,3);
\node at (6,2.5) {\textbf{C3}};
\draw[thick] (0,4) .. controls (1,4.2) and (3,4.2) .. (4,4) .. controls (5,3.8) and (7,3.8) .. (8,4) .. controls (9,4.2) and (11,4.2) .. (12,4);
\node at (6,3.5) {\textbf{C4}};

% Faille F1
\draw[very thick, popink] (11.8, 4.8) -- (12.6, -0.6);
\node[popink, right] at (12.4, 4.6) {\textbf{F1}};

% Bloc droit (compartiment abaissé) : couches décalées vers le bas
\draw[thick, popblue] (12.9, 3.4) -- (14.5, 3.4);
\draw[thick, popblue] (12.9, 2.4) -- (14.5, 2.4);
\draw[thick, popblue] (12.9, 1.4) -- (14.5, 1.4);
\draw[thick, popblue] (12.9, 0.4) -- (14.5, 0.4);
\draw[thick, popblue] (12.9, -0.6) -- (14.5, -0.6);
\node[popblue] at (13.7, 2.9) {C4};
\node[popblue] at (13.7, 1.9) {C3};
\node[popblue] at (13.7, 0.9) {C2};
\node[popblue] at (13.7, -0.1) {C1};

% Annotations
\node[align=left] at (2, 4.6) {Anticlinal\\(voûte)};
\node[align=left] at (10, 4.6) {Synclinal\\(creux)};
\node[align=center, popink] at (13.7, 4.2) {Compartiment\\abaissé};
\draw[<->, thick] (-0.5,0) -- (-0.5,4);
\node[rotate=90] at (-1.2,2) {Altitude relative};
\end{tikzpicture}
\legende{Coupe géologique : série plissée (C1 à C4) recoupée par une faille normale F1. Observer le décalage (rejet) des couches de part et d'autre de F1 : le compartiment de droite est abaissé.}
\end{popfigure}

\textbf{Raisonnement guidé :}
\begin{enumerate}
    \item \textbf{Observation :} les couches C1, C2, C3, C4 sont plissées (ondulées). La faille F1 est un plan qui coupe \emph{à travers} les couches plissées et les décale visiblement (rejet).
    \item \textbf{Hypothèses :} soit le pli est antérieur à la faille, soit la faille est antérieure au pli.
    \item \textbf{Test (principe de recoupement) :} la faille F1 \textbf{recoupe} les couches C1-C4 \emph{et} la structure plissée (l'anticlinal et le synclinal sont tronqués et décalés par F1).
    \item \textbf{Interprétation :} la structure qui recoupe est la plus jeune. Donc \textbf{F1 est postérieure au plissement}.
    \item \textbf{Conclusion :} le plissement a eu lieu \emph{avant} la mise en place de la faille.
\end{enumerate}

\textbf{Chronologie complète (6 événements) :}
\begin{enumerate}
    \item \textbf{Dépôt de C1} (le plus ancien, base de la série).
    \item \textbf{Dépôt de C2} (superposition sur C1).
    \item \textbf{Dépôt de C3} (superposition sur C2).
    \item \textbf{Dépôt de C4} (superposition sur C3, sommet de la série).
    \item \textbf{Plissement (phase compressive)} : formation de l'anticlinal et du synclinal. \textit{Justification : le pli déforme les couches C1 à C4, déposées initialement horizontales (principe d'horizontalité originelle).}
    \item \textbf{Mise en place de la faille F1 (phase extensive)}. \textit{Justification : la faille F1 recoupe et décale les couches \textbf{et} la structure plissée (principe de recoupement).}
\end{enumerate}
\end{experience}

\begin{attention}[Erreur fréquente : confondre « recouper » et « être déformé par »]
Une faille qui \textbf{recoupe} un pli est \textbf{plus jeune} que le pli.
Une faille qui est \textbf{elle-même plissée} (courbée avec les couches) est \textbf{plus ancienne} que le pli (elle a été déformée par lui).
\textbf{Règle d'or :} regardez la géométrie de la faille sur la coupe. Droite et tronquant le pli $\rightarrow$ postérieure. Courbée en suivant le pli $\rightarrow$ antérieure. Au Probatoire, les coupes montrent presque toujours une faille droite recoupant un pli.
\end{attention}

\courssub{Cas 3 : Série complexe avec discordance angulaire et faille (synthèse totale)}

C'est le cas « roi » du Probatoire : une coupe combinant plissement, érosion (discordance), nouveau dépôt, et faille finale. C'est l'histoire type des bassins sédimentaires camerounais (bassin de Douala, bassin du Logone-Birni).

\begin{exoresolu}[Chronologie complète d'une coupe complexe type examen (9 événements)]
\textbf{Énoncé :} à partir de la coupe ci-dessous, établir la chronologie relative des événements géologiques ayant affecté cette région. Justifier chaque étape.

\begin{popfigure}
\begin{tikzpicture}[scale=0.8]
% --- BLOC INFÉRIEUR (plissé, érodé) ---
\draw[thick, poppurple] (0,0) .. controls (1,0.3) and (3,0.3) .. (4,0) .. controls (5,-0.3) and (7,-0.3) .. (8,0) .. controls (9,0.3) and (11,0.3) .. (12,0);
\draw[thick, poppurple] (0,0.8) .. controls (1,1.1) and (3,1.1) .. (4,0.8) .. controls (5,0.5) and (7,0.5) .. (8,0.8) .. controls (9,1.1) and (11,1.1) .. (12,0.8);
\node[poppurple] at (6,0.4) {\textbf{C1 (schistes)}};
\draw[thick, popblue] (0,1.6) .. controls (1,1.9) and (3,1.9) .. (4,1.6) .. controls (5,1.3) and (7,1.3) .. (8,1.6) .. controls (9,1.9) and (11,1.9) .. (12,1.6);
\node[popblue] at (6,1.2) {\textbf{C2 (grès)}};
\draw[thick, popgreen!60!black] (0,2.4) .. controls (1,2.7) and (3,2.7) .. (4,2.4) .. controls (5,2.1) and (7,2.1) .. (8,2.4) .. controls (9,2.7) and (11,2.7) .. (12,2.4);
\node[popgreen!60!black] at (6,2.0) {\textbf{C3 (calcaires)}};

% DISCORDANCE D1
\draw[very thick, poporange, decorate, decoration={zigzag, amplitude=2pt, segment length=6pt}] (0,2.4) -- (12,2.4);
\node[poporange, above] at (6,2.45) {\textbf{D1 : discordance angulaire (surface d'érosion)}};

% --- BLOC SUPÉRIEUR (horizontal) ---
\fill[popgold!50] (0,2.4) rectangle (12,3.2);
\node at (6,2.8) {\textbf{C4 (grès)}};
\fill[poppink!50] (0,3.2) rectangle (12,4.0);
\node at (6,3.6) {\textbf{C5 (argiles marines)}};
\fill[poporange!30] (0,4.0) rectangle (12,4.8);
\node at (6,4.4) {\textbf{C6 (calcaires)}};

% FAILLE F1 (recoupe tout)
\draw[very thick, popink] (12.6, 5.2) -- (13.4, -0.6);
\node[popink, right] at (13.2, 5.0) {\textbf{F1 (faille normale)}};

% Compartiment droit (abaissé)
\fill[popgold!50] (13.8, 2.0) rectangle (15.4, 2.8);
\fill[poppink!50] (13.8, 2.8) rectangle (15.4, 3.6);
\fill[poporange!30] (13.8, 3.6) rectangle (15.4, 4.4);
\node at (14.6, 2.4) {C4};
\node at (14.6, 3.2) {C5};
\node at (14.6, 4.0) {C6};
\node[popink, align=center] at (14.6, 5.0) {Compartiment\\abaissé};

% Échelle
\draw[<->, thick] (-0.5,0) -- (-0.5,4.8);
\node[rotate=90] at (-1.2,2.4) {Altitude relative};
% Flèche temps
\draw[->, thick, popdark] (16, 4.5) -- (16, -0.2);
\node[rotate=-90, align=left, popdark] at (16.8, 2.1) {Plus récent $\rightarrow$ Plus ancien};
\end{tikzpicture}
\legende{Coupe type examen : série plissée (C1-C3) tronquée par une discordance angulaire D1, surmontée d'une couverture horizontale (C4-C6), le tout recoupé par une faille normale F1 qui décale l'ensemble.}
\end{popfigure}

\tcblower
\textbf{Solution détaillée :}

\textbf{Étape 1 : Observation et légendage}
\begin{itemize}
    \item \textbf{Soubassement (sous D1) :} 3 couches plissées : C1 (schistes, base), C2 (grès), C3 (calcaires, sommet). Structure : anticlinal et synclinal.
    \item \textbf{Surface D1 :} ligne irrégulière (zigzag) horizontale tronquant net les plis de C1-C3. \textbf{Discordance angulaire.}
    \item \textbf{Couverture (sur D1) :} 3 couches horizontales, non déformées : C4 (grès), C5 (argiles), C6 (calcaires).
    \item \textbf{Structure transversale :} faille normale F1 recoupant \emph{à la fois} le soubassement plissé, la discordance D1 et la couverture horizontale. Rejet visible (compartiment de droite abaissé).
\end{itemize}

\textbf{Étape 2 : Identification des événements élémentaires}

E1 : dépôt C1 ; E2 : dépôt C2 ; E3 : dépôt C3 ; E4 : plissement (déformation de C1-C3) ; E5 : émersion + érosion (création de D1) ; E6 : transgression + dépôt C4 ; E7 : dépôt C5 ; E8 : dépôt C6 ; E9 : mise en place de la faille F1.

\textbf{Étape 3 : Ordonnancement par les principes (du plus ancien au plus récent)}
\begin{enumerate}
    \item \textbf{Dépôt de C1 (schistes).} \textit{Justification : couche la plus basse du soubassement (principe de superposition).}
    \item \textbf{Dépôt de C2 (grès).} \textit{Justification : repose sur C1 (superposition).}
    \item \textbf{Dépôt de C3 (calcaires).} \textit{Justification : repose sur C2, sommet de la série plissée (superposition).}
    \item \textbf{Plissement (phase compressive : formation de l'anticlinal et du synclinal).} \textit{Justification : les couches C1, C2, C3, déposées horizontalement (principe d'horizontalité originelle), sont maintenant ondulées ; le pli affecte \textbf{toutes} les couches du soubassement.}
    \item \textbf{Émersion, érosion et planation (formation de la discordance D1).} \textit{Justification : la surface D1 est horizontale et \textbf{tronque} les couches plissées (C1, C2, C3 affleurent en bandeaux redressés). C'est une surface d'érosion (principe de recoupement : D1 recoupe les plis). \textbf{Ne pas oublier cet événement !}}
    \item \textbf{Transgression marine et dépôt de C4 (grès).} \textit{Justification : C4 repose horizontalement sur la surface d'érosion D1 (superposition sur D1 ; horizontalité originelle de la nouvelle série).}
    \item \textbf{Dépôt de C5 (argiles marines).} \textit{Justification : repose sur C4 (superposition).}
    \item \textbf{Dépôt de C6 (calcaires).} \textit{Justification : repose sur C5, sommet de la couverture (superposition).}
    \item \textbf{Mise en place de la faille normale F1.} \textit{Justification : F1 \textbf{recoupe} le soubassement plissé, \textbf{recoupe} la discordance D1 et \textbf{recoupe} la couverture horizontale C4-C6 (décalage visible). Principe de recoupement : F1 est plus jeune que \textbf{tout} ce qu'elle recoupe.}
\end{enumerate}

\textbf{Étape 4 : Présentation synthétique (frise verticale) — format Probatoire}

\begin{popfigure}
\begin{tikzpicture}[scale=0.75]
% Frise verticale
\draw[very thick, popdark] (0,0) -- (0,9.2);
\foreach \y/\label in {
    0.5/9. Mise en place de la faille F1 (extensive),
    1.5/8. Dépôt C6 (calcaires),
    2.5/7. Dépôt C5 (argiles marines),
    3.5/6. Transgression + dépôt C4 (grès),
    4.5/5. \textbf{Émersion + érosion (discordance D1)},
    5.5/4. Plissement (compressif : anticlinal/synclinal),
    6.5/3. Dépôt C3 (calcaires),
    7.5/2. Dépôt C2 (grès),
    8.5/1. Dépôt C1 (schistes) -- base
} {
    \draw[thick] (-0.1,\y) -- (0.1,\y);
    \node[align=left, anchor=west] at (0.3,\y) {\label};
}
% Flèche temps
\draw[->, thick, popblue] (8.5, 0.2) -- (8.5, 8.8);
\node[popblue, rotate=-90] at (9, 4.5) {Temps géologique $\rightarrow$};
% Barre de séparation discordance
\draw[poporange, very thick, decorate, decoration={zigzag, amplitude=1.5pt, segment length=4pt}] (-0.5, 5) -- (8.5, 5);
\node[poporange, align=center] at (4, 5.35) {\textbf{Discordance angulaire D1} (hiatus temporel majeur)};
\end{tikzpicture}
\legende{Frise chronologique verticale (du plus ancien en bas au plus récent en haut). La discordance D1 (zigzag orange) marque une rupture majeure : hiatus (temps non enregistré) correspondant à l'érosion du pli. La faille F1 est l'événement final.}
\end{popfigure}
\end{exoresolu}

\begin{aretenir}[Rédaction type « Probatoire » pour une chronologie]
Au Probatoire, on attend une liste numérotée \textbf{du plus ancien au plus récent}, chaque ligne contenant :
\textbf{Numéro. Événement (nature + unités concernées). — Justification : principe invoqué + observation précise sur la coupe.}
Exemple : « 5. Émersion et érosion (formation de la discordance D1). — Justification : la surface D1 recoupe les plis affectant C1-C3 (principe de recoupement) et est plane et horizontale (surface d'aplanissement par l'érosion). »
\end{aretenir}

\courssub{Exploitation des données paléontologiques : caler la chronologie dans l'échelle des temps}

Les principes de stratigraphie donnent un ordre \emph{relatif} (A avant B). Les fossiles stratigraphiques (vus dans la section précédente) permettent d'attribuer des \emph{âges relatifs précis} (étages géologiques : Aptien, Albien, Cénomanien, etc.) et de corréler des couches éloignées.

\begin{exemplebox}[Corrélation par fossiles stratigraphiques : le bassin de Douala]
Deux forages réalisés dans le bassin de Douala montrent des séries comparables mais décalées en profondeur. Les fossiles stratigraphiques permettent de repérer les niveaux de même âge.
\end{exemplebox}

\begin{popfigure}
\begin{tikzpicture}[scale=0.8]
% Puits 1 (Logbaba)
\draw[thick] (0,0) -- (0,6);
\draw[thick] (1,0) -- (1,6);
\fill[poppurple!30] (0,0) rectangle (1,1); \node[font=\small] at (0.5,0.5) {Socle};
\fill[popblue!30] (0,1) rectangle (1,2.5); \node[font=\small, align=center] at (0.5,1.75) {C1 : argiles\\à \textit{Ostrea}};
\fill[popgreen!30] (0,2.5) rectangle (1,4); \node[font=\small, align=center] at (0.5,3.25) {C2 : calcaires\\à orbitolines};
\fill[poporange!30] (0,4) rectangle (1,5.5); \node[font=\small, align=center] at (0.5,4.75) {C3 : grès à\\\textit{Trigonia}};
\fill[popgold!30] (0,5.5) rectangle (1,6); \node[font=\small] at (0.5,5.75) {Alluvions};
\node at (0.5,6.3) {\textbf{Puits Logbaba}};

% Puits 2 (Makepe)
\draw[thick] (3,0.5) -- (3,6.5);
\draw[thick] (4,0.5) -- (4,6.5);
\fill[poppurple!30] (3,0.5) rectangle (4,1.5); \node[font=\small] at (3.5,1) {Socle};
\fill[popblue!30] (3,1.5) rectangle (4,3); \node[font=\small, align=center] at (3.5,2.25) {C1' : argiles\\à \textit{Ostrea}};
\fill[popgreen!30] (3,3) rectangle (4,4.5); \node[font=\small, align=center] at (3.5,3.75) {C2' : calcaires\\à orbitolines};
\fill[poporange!30] (3,4.5) rectangle (4,6); \node[font=\small, align=center] at (3.5,5.25) {C3' : grès à\\\textit{Trigonia}};
\fill[popgold!30] (3,6) rectangle (4,6.5); \node[font=\small] at (3.5,6.25) {Alluvions};
\node at (3.5,6.8) {\textbf{Puits Makepe}};

% Lignes de corrélation (isochrones)
\draw[dashed, thick, popink] (0,1) -- (4,1.5);
\draw[dashed, thick, popink] (0,2.5) -- (4,3);
\draw[dashed, thick, popink] (0,4) -- (4,4.5);
\draw[dashed, thick, popink] (0,5.5) -- (4,6);

\node[popink, align=center] at (2, 0.4) {Corrélations par identité paléontologique};
\node[popink, font=\small] at (2, 1.75) {\textit{Ostrea} sp. (Albien)};
\node[popink, font=\small] at (2, 3.25) {Orbitolines (Aptien)};
\node[popink, font=\small] at (2, 4.75) {\textit{Trigonia} (Cénomanien)};
\end{tikzpicture}
\legende{Corrélation de deux forages du bassin de Douala (Logbaba et Makepe) grâce aux fossiles stratigraphiques (\textit{Ostrea}, orbitolines, \textit{Trigonia}). Les lignes pointillées matérialisent les niveaux isochrones (de même âge).}
\end{popfigure}

Grâce aux fossiles stratigraphiques (orbitolines à l'Aptien, \textit{Trigonia} au Cénomanien), les géologues peuvent :
\begin{itemize}
    \item confirmer que C1/C1', C2/C2', C3/C3' sont des \textbf{niveaux de même âge} (isochrones) malgré l'éloignement des forages (principe d'identité paléontologique) ;
    \item reconstituer la \textbf{transgression crétacée} (la mer envahit progressivement le continent) ;
    \item identifier les \textbf{roches mères} (argiles riches en matière organique) et les \textbf{roches réservoirs} (grès poreux) utiles à l'exploration pétrolière.
\end{itemize}

\begin{savaistu}[Le bassin de Douala : une archive géologique exploitée]
Le bassin sédimentaire de Douala (littoral camerounais) est un bassin de marge passive ouvert lors de la séparation de l'Afrique et de l'Amérique du Sud au Crétacé. Sa chronologie, établie \textbf{exactement par la méthode de cette leçon} (superposition + discordances + fossiles + failles), révèle :
\begin{enumerate}
    \item \textbf{Phase de rift (Crétacé inférieur) :} dépôts continentaux (grès, conglomérats) dans des fossés effondrés bordés de failles normales synsédimentaires.
    \item \textbf{Transgression marine (Aptien-Albien) :} dépôt d'argiles noires riches en matière organique (roches mères pétrolifères) et de calcaires à orbitolines.
    \item \textbf{Phase tectonique compressive (Santonien, $\approx$ 85 Ma) :} déformation, soulèvement, érosion $\rightarrow$ \textbf{discordance majeure} (équivalent de notre D1).
    \item \textbf{Phase post-rift (Cénomanien à Actuel) :} dépôts de couverture (grès, calcaires, argiles) subhorizontaux, peu déformés.
\end{enumerate}
C'est en maîtrisant cette chronologie que le Cameroun a découvert et exploité le gaz de \textbf{Logbaba} et développe le projet gazier de \textbf{Makepe}. La géologie stratigraphique n'est pas qu'académique : elle alimente l'économie nationale !
\end{savaistu}

\courssub{Synthèse : la « check-list » pour réussir l'exercice de chronologie}

\begin{aretenir}[Check-list « zéro faute » pour le Probatoire]
Avant de rendre votre copie, vérifiez que vous avez :
\begin{enumerate}
    \item \textbf{Légendé} la coupe (couches C1, C2... ; pli P1 ; faille F1 ; discordance D1 ; intrusion I1).
    \item \textbf{Listé tous} les événements : dépôts (autant que de couches distinctes), déformations (plis, failles), surfaces d'érosion (discordances), intrusions éventuelles.
    \item \textbf{Ordonné} du plus ancien (base) au plus récent (sommet / structure transversale finale).
    \item \textbf{Justifié chaque relation d'ordre} par un principe cité nommément : « superposition », « recoupement », « horizontalité originelle », « identité paléontologique ».
    \item \textbf{N'avez pas oublié l'érosion} entre un pli tronqué et une couverture horizontale (discordance = émersion + érosion + hiatus).
    \item \textbf{Présenté} le résultat sous forme de liste numérotée \textbf{et/ou} de frise verticale.
\end{enumerate}
\end{aretenir}

\begin{exercice}[Entraînement : coupe de la région de Ngaoundéré (Adamaoua)]
On donne la coupe schématique simplifiée suivante :
\begin{itemize}
    \item socle granitique (Précambrien) ;
    \item série volcano-sédimentaire paléozoïque (schistes, grès, calcaires) plissée en anticlinal ;
    \item discordance angulaire tronquant l'anticlinal ;
    \item coulée de basalte (trapp) horizontale recouvrant la discordance ;
    \item faille normale recoupant le basalte et le socle ;
    \item alluvions récentes déposées dans le fossé d'effondrement.
\end{itemize}
\textbf{Travail demandé :}
\begin{enumerate}
    \item Légender la coupe (inventez les codes : S0, C1, C2, C3, P1, D1, B1, F1, A1).
    \item Établir la chronologie relative complète (9 événements) en justifiant chaque étape.
    \item Représenter le résultat sous forme de frise verticale.
\end{enumerate}
\emph{Indication : le basalte (trapp) est une coulée de lave. Une faille qui recoupe une coulée de lave est postérieure au volcanisme.}
\end{exercice}

\begin{corrige}[Exercice : région de Ngaoundéré]
\textbf{1. Légende proposée :} S0 : socle granitique ; C1 : schistes ; C2 : grès ; C3 : calcaires (série paléozoïque) ; P1 : anticlinal ; D1 : discordance angulaire ; B1 : basalte (trapp) ; F1 : faille normale ; A1 : alluvions récentes.

\textbf{2. Chronologie (du plus ancien au plus récent) :}
\begin{enumerate}
    \item \textbf{Mise en place du socle granitique S0} (Précambrien). — \textit{Justification : base de la coupe, recoupé par toutes les structures (recoupement).}
    \item \textbf{Dépôt de C1 (schistes).} — \textit{Justification : superposition sur S0.}
    \item \textbf{Dépôt de C2 (grès).} — \textit{Justification : superposition sur C1.}
    \item \textbf{Dépôt de C3 (calcaires).} — \textit{Justification : superposition sur C2.}
    \item \textbf{Plissement P1 (anticlinal).} — \textit{Justification : déforme C1, C2, C3 déposées horizontalement (horizontalité originelle).}
    \item \textbf{Émersion + érosion (formation de la discordance D1).} — \textit{Justification : D1 recoupe et tronque le pli P1 et les couches C1-C3 (recoupement) ; surface plane = aplanissement par l'érosion.}
    \item \textbf{Éruption volcanique : épanchement du basalte B1 (trapp).} — \textit{Justification : B1 repose horizontalement sur D1 (superposition + horizontalité originelle) et recouvre tout le soubassement.}
    \item \textbf{Mise en place de la faille normale F1.} — \textit{Justification : F1 recoupe B1, D1, P1 et S0 (recoupement) ; elle crée un fossé d'effondrement.}
    \item \textbf{Dépôt des alluvions A1.} — \textit{Justification : elles comblent le fossé créé par F1 (superposition dans le fossé) ; événement le plus récent.}
\end{enumerate}

\textbf{3. Frise verticale :} à tracer sur le modèle de la frise de l'exercice résolu : 9 graduations, le zigzag orange sur D1 (entre les événements 5 et 7), la flèche du temps dirigée vers le haut (du plus ancien en bas au plus récent en haut).
\end{corrige}

\begin{attention}[Dernier conseil de « Prof »]
Face à une coupe complexe le jour de l'examen : \textbf{respirez}. Prenez deux minutes pour \textbf{coder} chaque couche au crayon (hachures différentes). Repérez \textbf{visuellement} qui coupe qui. Écrivez les événements au brouillon \textbf{en vrac} d'abord, puis ordonnez-les. Ne rédigez la version finale qu'une fois l'ordre figé. La rigueur de la \textbf{justification} vaut souvent plus de points que la liste elle-même. Bon courage, géologues en herbe ! Le sous-sol camerounais n'attend que vous pour livrer ses secrets.
\end{attention}

\newpage

\coursec{Activité d'intégration}

\courssub{Objectifs de l'activité}

À l'issue de cette activité d'intégration, tu seras capable de :

\begin{itemize}
\item lire et légender une \cle{coupe géologique} réelle en identifiant les couches, les structures (plis, failles, discordances) et les fossiles qu'elle contient ;
\item appliquer de manière combinée les trois principes de la \cle{stratigraphie} (superposition, continuité, recoupement) pour ordonner des événements géologiques ;
\item utiliser des \cle{fossiles stratigraphiques} pour dater relativement des formations et les rattacher à une ère géologique ;
\item rédiger une \cle{chronologie relative} complète et justifiée, du plus ancien au plus récent, puis la traduire sous forme de frise chronologique.
\end{itemize}

Cette activité mobilise l'ensemble des savoir-faire de la leçon : c'est exactement le type de raisonnement exigé au Probatoire dans les exercices de géologie.

\courssub{Situation}

\textbf{Géologues d'un jour : l'histoire cachée de la carrière de Nkolfoulou}

Au cours d'une sortie pédagogique organisée par le club Sciences de ton lycée, ta classe se rend à \cle{Nkolfoulou}, localité située à une vingtaine de kilomètres de Yaoundé, sur l'axe Yaoundé--Mbalmayo. Là, l'exploitation d'une carrière de matériaux de construction a mis à nu une falaise de près de \SI{25}{m} de hauteur. Le géologue qui accompagne la visite s'arrête devant la paroi et déclare : « Cette falaise est un livre ouvert. Chaque couche, chaque cassure, chaque fossile est une phrase de l'histoire de cette région. À vous de la lire. »

Sur la paroi, vous observez, de bas en haut :

\begin{itemize}
\item \textbf{Ensemble inférieur (couches 1, 2 et 3)} : trois couches sédimentaires \cle{plissées} (déformées en ondulations). La couche 1 (la plus profonde) est un grès ; la couche 2 est un calcaire renfermant des \cle{trilobites} ; la couche 3 est un schiste. Ces trois couches sont \cle{tronquées} en leur sommet par une surface plane d'\cle{érosion} qui les coupe obliquement.
\item \textbf{Ensemble supérieur (couches 4 et 5)} : au-dessus de la surface d'érosion reposent deux couches parfaitement \cle{horizontales}. La couche 4 est un calcaire fossilifère contenant des \cle{ammonites} ; la couche 5 (sommitale) est une argile sableuse sans fossile visible.
\item \textbf{Accident tectonique} : une \cle{faille} recoupe les couches 1, 2, 3 et 4, avec un rejet visible (les couches sont décalées d'environ \SI{2}{m} de part et d'autre), mais elle \textbf{s'arrête sous la couche 5}, qu'elle n'affecte pas.
\end{itemize}

Ton professeur te confie la mission suivante : \emph{« En géologue rigoureux, établis la chronologie relative complète des événements qui ont façonné cette région, justifie chaque étape par un principe de stratigraphie, et attribue chaque ensemble de couches à son ère géologique. »}

\courssub{Consignes}

Travail en groupes de quatre élèves, puis mise en commun au tableau.

\begin{enumerate}
\item \textbf{Légender la coupe.} Reproduis schématiquement la coupe de la carrière et légende-la : couches 1 à 5, surface d'érosion, faille, emplacements des fossiles.
\item \textbf{Identifier les événements.} Dresse la liste de tous les événements géologiques que traduit cette coupe (dépôts, plissement, érosion, sédimentation, fracturation).
\item \textbf{Appliquer le principe de superposition.} À l'intérieur de chaque ensemble (inférieur puis supérieur), ordonne les couches de la plus ancienne à la plus récente. Justifie.
\item \textbf{Appliquer le principe de recoupement.} Le plissement est-il antérieur ou postérieur au dépôt des couches 1 à 3 ? La faille est-elle antérieure ou postérieure au dépôt de la couche 4 ? Et au dépôt de la couche 5 ? Justifie chaque réponse.
\item \textbf{Appliquer le principe d'identité paléontologique.} Sachant que les trilobites caractérisent l'\cle{ère Primaire} (Paléozoïque) et que les ammonites caractérisent l'\cle{ère Secondaire} (Mésozoïque), attribue chaque ensemble de couches à son ère. Que représente alors la surface d'érosion sur le plan chronologique ?
\item \textbf{Rédiger la chronologie complète.} Rédige, du plus ancien au plus récent, la succession de tous les événements géologiques de la région de Nkolfoulou, en justifiant chaque position par le principe utilisé.
\item \textbf{Construire la frise chronologique.} Traduis ta chronologie sous forme d'une frise verticale graduée en temps relatif (plus ancien en bas, plus récent en haut), en y plaçant les ères géologiques correspondantes.
\end{enumerate}

\courssub{Ressources à mobiliser}

\begin{aretenir}[Les outils du géologue pour dater en relatif]
\begin{itemize}
\item \textbf{Principe de superposition} : dans une série de couches non renversées, une couche est plus ancienne que celle qui la recouvre et plus récente que celle qui la supporte.
\item \textbf{Principe de continuité} : une couche a même âge sur toute son extension ; deux affleurements d'une même couche, séparés par une faille ou une vallée, sont contemporains.
\item \textbf{Principe de recoupement} : une structure géologique (faille, pli, filon, surface d'érosion) est plus récente que les couches qu'elle recoupe ou déforme, et plus ancienne que les couches qu'elle n'affecte pas.
\item \textbf{Principe d'identité paléontologique} : deux couches contenant les mêmes fossiles stratigraphiques ont le même âge. Un bon fossile stratigraphique est abondant, largement répandu et d'existence courte (ex. : trilobites $\rightarrow$ Primaire ; ammonites $\rightarrow$ Secondaire).
\item \textbf{Discordance} : surface d'érosion séparant des couches plissées anciennes de couches horizontales plus récentes ; elle traduit un hiatus (lacune) dans l'enregistrement sédimentaire.
\end{itemize}
\end{aretenir}

\begin{attention}[Pièges fréquents]
\begin{itemize}
\item Une couche \textbf{plissée} s'est d'abord déposée \emph{horizontalement} (principe d'horizontalité initiale) : le plissement est donc un événement \emph{postérieur} au dépôt.
\item La faille est \emph{plus récente} que toutes les couches qu'elle décale, mais \emph{plus ancienne} que la première couche intacte qui la surmonte.
\item Ne confonds pas datation \textbf{relative} (ordre des événements) et datation \textbf{absolue} (âge en millions d'années) : ici, seule la chronologie relative est demandée.
\item Attention au vocabulaire : on dit qu'une couche est \emph{antérieure} ou \emph{postérieure} à un événement, \emph{plus ancienne} ou \emph{plus récente} qu'une autre couche ; évite les formulations vagues comme « la couche est avant ».
\end{itemize}
\end{attention}

\courssub{Démarche et corrigé commenté}

\begin{exoresolu}[Chronologie de la carrière de Nkolfoulou — corrigé détaillé]
\textbf{Énoncé.} À partir des observations faites sur la paroi de la carrière de Nkolfoulou (trois couches plissées à trilobites tronquées par une surface d'érosion, surmontées de deux couches horizontales à ammonites, l'ensemble étant recoupé par une faille n'atteignant pas la couche sommitale), établir la chronologie relative complète et justifiée des événements géologiques de la région, puis attribuer chaque série à son ère.

\tcblower

\textbf{Étape 1 : lecture et légendage de la coupe.}

On identifie cinq couches numérotées de bas en haut (1 : grès ; 2 : calcaire à trilobites ; 3 : schiste ; 4 : calcaire à ammonites ; 5 : argile sableuse), une surface d'érosion entre les couches 3 et 4, et une faille recoupant les couches 1 à 4.

\begin{popfigure}
\begin{center}
\begin{tikzpicture}[scale=0.9]
% couches plissées
\draw[fill=popgoldL] (0,0) -- (6,0) -- (6,1.1) .. controls (4.5,1.6) and (3,0.6) .. (1.5,1.1) .. controls (0.75,1.35) and (0.4,1.2) .. (0,1.1) -- cycle;
\draw[fill=popgreenL] (0,1.1) .. controls (0.4,1.2) and (0.75,1.35) .. (1.5,1.1) .. controls (3,0.6) and (4.5,1.6) .. (6,1.1) -- (6,2.0) .. controls (4.5,2.5) and (3,1.5) .. (1.5,2.0) .. controls (0.75,2.25) and (0.4,2.1) .. (0,2.0) -- cycle;
\draw[fill=popblueL] (0,2.0) .. controls (0.4,2.1) and (0.75,2.25) .. (1.5,2.0) .. controls (3,1.5) and (4.5,2.5) .. (6,2.0) -- (6,2.9) .. controls (4.5,3.4) and (3,2.4) .. (1.5,2.9) .. controls (0.75,3.15) and (0.4,3.0) .. (0,2.9) -- cycle;
% surface d'érosion
\draw[popink,very thick] (0,3.35) -- (6,3.35);
% couches horizontales
\draw[fill=poporangeL] (0,3.35) rectangle (6,4.35);
\draw[fill=popcream] (0,4.35) rectangle (6,5.35);
% faille
\draw[popdark,very thick,dashed] (3.6,0) -- (3.3,4.35);
% légendes
\node[right,font=\small] at (6.15,0.55) {1 : grès};
\node[right,font=\small] at (6.15,1.55) {2 : calcaire à trilobites};
\node[right,font=\small] at (6.15,2.45) {3 : schiste};
\node[right,font=\small] at (6.15,3.85) {4 : calcaire à ammonites};
\node[right,font=\small] at (6.15,4.85) {5 : argile sableuse};
\node[above,font=\small\itshape] at (3,3.4) {surface d'érosion (discordance)};
\node[below,font=\small\itshape] at (3.45,-0.1) {faille};
\end{tikzpicture}
\end{center}
\legende{Coupe schématique de la carrière de Nkolfoulou : couches plissées 1 à 3 tronquées par une surface d'érosion, surmontées des couches horizontales 4 et 5 ; la faille recoupe les couches 1 à 4 mais s'arrête sous la couche 5.}
\end{popfigure}

\textbf{Étape 2 : ordre interne des deux ensembles (principe de superposition).}

Dans l'ensemble inférieur comme dans l'ensemble supérieur, aucun indice de renversement n'est observé. D'après le \cle{principe de superposition} :
\[ \text{couche 1} < \text{couche 2} < \text{couche 3} \qquad \text{et} \qquad \text{couche 4} < \text{couche 5} \]
(le signe $<$ signifie ici « plus ancien que »).

\textbf{Étape 3 : datation relative du plissement et de l'érosion (principe de recoupement).}

Les couches 1 à 3 se sont déposées horizontalement, puis ont été \cle{plissées} : le plissement est donc \textbf{postérieur au dépôt de la couche 3}. La surface d'érosion recoupe (tronque) les plis : elle est donc \textbf{postérieure au plissement}. Enfin, les couches 4 et 5, horizontales et non plissées, reposent sur cette surface : elles sont \textbf{postérieures à l'érosion}. L'ensemble constitue une \cle{discordance angulaire}.

\textbf{Étape 4 : datation relative de la faille (principe de recoupement).}

La faille décale les couches 1, 2, 3 \emph{et} 4 : elle est donc \textbf{postérieure au dépôt de la couche 4}. Mais elle n'affecte pas la couche 5 : elle est donc \textbf{antérieure au dépôt de la couche 5}. La faille se situe donc entre le dépôt de la couche 4 et celui de la couche 5.

\textbf{Étape 5 : attribution aux ères (principe d'identité paléontologique).}

Les \cle{trilobites} de la couche 2 sont des fossiles stratigraphiques de l'\textbf{ère Primaire} (Paléozoïque) : les couches 1 à 3 sont primaires. Les \cle{ammonites} de la couche 4 caractérisent l'\textbf{ère Secondaire} (Mésozoïque) : les couches 4 et 5 sont secondaires. La surface d'érosion représente donc un \cle{hiatus} correspondant à la limite Primaire--Secondaire : une longue période pendant laquelle aucune sédimentation n'a été enregistrée dans la région.

\textbf{Étape 6 : chronologie complète (du plus ancien au plus récent).}

\begin{enumerate}
\item Dépôt de la couche 1 (grès) — Primaire ;
\item Dépôt de la couche 2 (calcaire à trilobites) — Primaire ;
\item Dépôt de la couche 3 (schiste) — Primaire ;
\item \textbf{Plissement} des couches 1 à 3 (déformation tectonique) ;
\item \textbf{Érosion} des couches plissées (mise en place de la surface de discordance) — long hiatus ;
\item Dépôt de la couche 4 (calcaire à ammonites) — Secondaire ;
\item \textbf{Faille} recoupant les couches 1 à 4 ;
\item Dépôt de la couche 5 (argile sableuse) — Secondaire, événement le plus récent.
\end{enumerate}

\textbf{Étape 7 : frise chronologique.}

\begin{popfigure}
\begin{center}
\begin{tikzpicture}[scale=0.95]
% axe du temps
\draw[->,very thick,popink] (0,0) -- (0,8.6) node[above,font=\small\bfseries] {Temps (plus récent en haut)};
% ères
\draw[fill=popgreenL,draw=popgreen] (-3.4,0) rectangle (-0.3,3.6);
\draw[fill=poporangeL,draw=poporange] (-3.4,3.6) rectangle (-0.3,8.4);
\node[rotate=90,font=\small\bfseries] at (-1.85,1.8) {Ère Primaire (Paléozoïque)};
\node[rotate=90,font=\small\bfseries] at (-1.85,6.0) {Ère Secondaire (Mésozoïque)};
% événements
\foreach \y/\txt in {0.4/{Dépôt couche 1 (grès)}, 1.2/{Dépôt couche 2 (calcaire à trilobites)}, 2.0/{Dépôt couche 3 (schiste)}, 2.8/{Plissement des couches 1 à 3}, 3.6/{Érosion : surface de discordance (hiatus)}, 4.6/{Dépôt couche 4 (calcaire à ammonites)}, 5.6/{Faille (recoupe 1 à 4)}, 6.8/{Dépôt couche 5 (argile sableuse)}}{
  \fill[popink] (0,\y) circle (1.5pt);
  \node[right,font=\small] at (0.25,\y) {\txt};
}
% hiatus
\draw[<->,popdark,thick] (0.9,3.65) -- (0.9,4.55) node[midway,right,font=\small\itshape] {hiatus};
\end{tikzpicture}
\end{center}
\legende{Frise chronologique relative des événements géologiques de la région de Nkolfoulou : les événements se lisent de bas en haut, du plus ancien au plus récent ; la surface d'érosion marque un hiatus entre l'ère Primaire et l'ère Secondaire.}
\end{popfigure}

\textbf{Conclusion.} Chaque position dans la chronologie est justifiée : la superposition pour l'ordre des dépôts, le recoupement pour le plissement, l'érosion et la faille, l'identité paléontologique pour l'attribution aux ères. La région de Nkolfoulou a donc connu deux grandes phases de sédimentation (Primaire puis Secondaire), séparées par une phase de déformation et d'érosion, et une fracturation tardive intercalée entre les dépôts des couches 4 et 5.
\end{exoresolu}

\courssub{Bilan de l'activité}

Cette activité t'a permis de mettre en évidence que :

\begin{itemize}
\item une coupe géologique se lit comme un \cle{document historique} : chaque couche, structure ou fossile est la trace d'un événement daté relativement aux autres ;
\item les trois principes de stratigraphie sont \textbf{complémentaires} : la superposition ordonne les dépôts, le recoupement situe les déformations (plis, failles) et les interruptions (érosions), l'identité paléontologique rattache les formations aux grandes ères de l'histoire de la Terre ;
\item une \cle{discordance} traduit toujours une succession d'événements : sédimentation, plissement, érosion, puis nouvelle sédimentation, avec un hiatus chronologique souvent considérable ;
\item la datation relative ne donne pas d'âge en années, mais un \textbf{ordre rigoureux et justifié} des événements : c'est la première étape de toute reconstitution de l'histoire géologique d'une région, avant toute datation absolue.
\end{itemize}

Tu disposes désormais de la démarche complète du géologue : observer, identifier, ordonner, justifier. C'est cette méthode qui te sera demandée au Probatoire face à toute coupe géologique inédite.

\newpage

\coursec{Méthodes et exercices résolus}

\courssub{Méthodes et exercices résolus}

\begin{methode}[Appliquer les principes de stratigraphie à une coupe géologique]
Pour analyser une coupe géologique et ordonner les couches sédimentaires, procède toujours dans le même ordre :
\begin{enumerate}
\item \textbf{Repérer les couches sédimentaires} : identifie chaque strate par sa lettre ou son numéro et note sa position (en bas, au milieu, en haut de la coupe).
\item \textbf{Appliquer le principe de superposition} : dans une série sédimentaire non renversée, une couche est plus ancienne que celle qui la recouvre et plus récente que celle qui la supporte. La couche la plus profonde est donc la plus ancienne.
\item \textbf{Vérifier l'absence de renversement} : si la série est plissée, cherche des indices de polarité (discordances, fossiles, sédimentation granoclassée) avant de conclure. En cas de doute, précise « si la série n'est pas renversée ».
\item \textbf{Appliquer le principe de continuité} : une même couche a le même âge sur toute son extension ; deux couches séparées par une faille mais de même nature et de même niveau sont contemporaines.
\item \textbf{Appliquer le principe de recoupement} : toute structure qui recoupe une couche (faille, dyke, filon, intrusion magmatique) est \emph{postérieure} à cette couche ; tout ce qui est recoupé est plus ancien.
\item \textbf{Conclure par une phrase} : écris la succession chronologique du plus ancien au plus récent, en justifiant chaque étape par le nom du principe utilisé.
\end{enumerate}
\textbf{Réflexe} : toujours citer le principe par son nom (superposition, continuité, recoupement) — c'est ce qui rapporte les points de justification.
\end{methode}

\begin{popfigure}
\begin{center}
\begin{tikzpicture}[scale=1.0]
% Couche A (argiles)
\fill[popgoldL] (0,0) rectangle (6,1);
% Couche B (grès)
\fill[poporangeL] (0,1) rectangle (6,2);
% Couche C (calcaires)
\fill[popblueL] (0,2.6) rectangle (6,3.6);
% Filon basaltique F recoupant A et B, s'arrêtant sous C
\fill[popdark] (2.7,0) -- (3.3,0) -- (3.5,2.6) -- (2.9,2.6) -- cycle;
% Surface d'érosion entre B et C
\draw[popink,thick] (0,2.6) -- (6,2.6);
% Lignes de séparation
\draw[popmuted] (0,1) -- (6,1);
\draw[popmuted] (0,2) -- (6,2);
% Légendes pointées
\draw[popmuted,thin] (5.6,0.5) -- (6.4,0.5) node[right,font=\small]{A : argiles};
\draw[popmuted,thin] (5.6,1.5) -- (6.4,1.5) node[right,font=\small]{B : grès};
\draw[popmuted,thin] (5.6,3.1) -- (6.4,3.1) node[right,font=\small]{C : calcaires};
\draw[popmuted,thin] (3.2,1.3) -- (1.2,-0.6) node[below,font=\small]{F : filon de basalte};
\draw[popmuted,thin] (0.8,2.6) -- (0.4,4.1) node[above,font=\small]{surface d'érosion};
\end{tikzpicture}
\end{center}
\legende{Coupe schématique : le filon F recoupe les couches A et B (il leur est donc postérieur, principe de recoupement) mais s'arrête sous la couche C, qui le recouvre (C est donc postérieure à F).}
\end{popfigure}

\begin{exoresolu}[Datation relative d'une série sédimentaire simple]
On réalise une coupe dans une carrière de la région de l'Adamaoua. Du bas vers le haut, on observe : des argiles (couche A), des grès (couche B), des calcaires (couche C). Un filon de basalte F recoupe les couches A et B mais s'arrête sous la couche C.
\begin{enumerate}
\item Classer les couches A, B, C de la plus ancienne à la plus récente.
\item Situer la mise en place du filon F dans cette chronologie.
\end{enumerate}
\tcblower
\textbf{1. Classement des couches.}\\
Les couches A, B et C sont des couches sédimentaires superposées et non perturbées (aucun indice de renversement n'est signalé). D'après le \cle{principe de superposition}, dans une série sédimentaire non renversée, toute couche est plus récente que celle qui la supporte et plus ancienne que celle qui la recouvre. Or A supporte B, et B supporte C. Donc :
\[ A \;(\text{la plus ancienne}) \;\longrightarrow\; B \;\longrightarrow\; C \;(\text{la plus récente}) \]

\textbf{2. Position du filon F.}\\
Le filon F \emph{recoupe} les couches A et B : d'après le \cle{principe de recoupement}, une structure qui recoupe une couche lui est postérieure. Donc F est plus récent que A et que B.\\
En revanche, F s'arrête sous la couche C : il ne la recoupe pas, et C le recouvre. La sédimentation de C est donc postérieure à la mise en place du filon (le filon existait déjà, érodé, lorsque C s'est déposée). Donc F est plus ancien que C.\\
\textbf{Conclusion} : la chronologie relative des événements, du plus ancien au plus récent, est :
\[ \text{dépôt de A} \rightarrow \text{dépôt de B} \rightarrow \text{mise en place du filon F} \rightarrow \text{dépôt de C} \]
\end{exoresolu}

\begin{methode}[Établir la chronologie relative complète des événements d'une région]
Une coupe géologique enregistre non seulement des dépôts, mais aussi des \textbf{événements} : plissements, failles, érosions, intrusions. Pour reconstituer l'histoire complète :
\begin{enumerate}
\item \textbf{Faire l'inventaire} de tous les éléments de la coupe : couches sédimentaires, plis, failles, surfaces d'érosion (discordances), filons ou massifs intrusifs.
\item \textbf{Dater les dépôts} entre eux par le principe de superposition.
\item \textbf{Dater chaque événement tectonique} (pli, faille) par le principe de recoupement : un pli ou une faille est postérieur à la plus récente des couches qu'il affecte, et antérieur à la plus ancienne des couches qui le recouvrent sans être affectées. C'est la méthode du \cle{cadrage} : l'événement est « coincé » entre deux repères.
\item \textbf{Dater les discordances} : la surface d'érosion est postérieure aux couches tronquées et antérieure aux couches qui la surmontent.
\item \textbf{Ordonner tous les événements} dans un tableau ou une frise chronologique, du plus ancien (en bas) au plus récent (en haut).
\item \textbf{Rédiger la conclusion} sous forme d'une « histoire géologique » rédigée en phrases.
\end{enumerate}
\textbf{Réflexe} : pour chaque événement, pose-toi deux questions : « Quelle est la couche la plus récente qu'il affecte ? » et « Quelle est la couche la plus ancienne qui n'est pas affectée ? ».
\end{methode}

\begin{exoresolu}[Histoire géologique d'une région : pli, faille et intrusion]
Une coupe montre, du bas vers le haut : des couches 1, 2 et 3 plissées ; une faille F qui décale les couches 1, 2 et 3 ; une surface d'érosion E qui tronque le pli et la faille ; des couches horizontales 4 et 5 déposées au-dessus de E ; un dyke de granité D qui recoupe les couches 1 à 4 mais pas la couche 5.\\
Reconstituer la chronologie complète des événements géologiques de cette région.
\tcblower
\textbf{Étape 1 : inventaire.} Éléments présents : couches sédimentaires 1, 2, 3, 4, 5 ; un plissement P ; une faille F ; une surface d'érosion E ; un dyke D.

\textbf{Étape 2 : datation des dépôts.} Par le \cle{principe de superposition} (séries non renversées) : 1 est plus ancienne que 2, plus ancienne que 3 ; de même 4 est plus ancienne que 5. Les couches 4 et 5 reposent sur la surface d'érosion E, donc elles sont postérieures à tout ce que E tronque.

\textbf{Étape 3 : cadrage du plissement P.} Le pli affecte les couches 1, 2 et 3 : il est donc \emph{postérieur au dépôt de la couche 3}. La surface E tronque le pli : le plissement est \emph{antérieur à l'érosion E}. Donc : dépôt de 3 $<$ P $<$ E.

\textbf{Étape 4 : cadrage de la faille F.} La faille décale les couches 1, 2, 3 : elle est postérieure au dépôt de 3. La surface E tronque aussi la faille (elle ne décale pas les couches 4 et 5) : F est antérieure à E. On ne peut pas, avec ces seules données, dire si F est antérieure ou postérieure au plissement P ; on note simplement : dépôt de 3 $<$ F $<$ E.

\textbf{Étape 5 : cadrage du dyke D.} Par le \cle{principe de recoupement} : D recoupe les couches 1 à 4, donc D est postérieur au dépôt de 4 ; D ne recoupe pas la couche 5, donc D est antérieur au dépôt de 5.

\textbf{Étape 6 : chronologie finale} (du plus ancien au plus récent) :
\begin{center}
\begin{tabularx}{\linewidth}{|c|X|}\hline
\rowcolor{popblueL} \textbf{Rang} & \textbf{Événement (justification)} \\ \hline
1 & Dépôt des couches 1, 2, 3 (superposition) \\ \hline
2 & Plissement P et faille F (postérieurs à 3, antérieurs à E — recoupement) \\ \hline
3 & Érosion E formant la discordance (tronque P et F) \\ \hline
4 & Dépôt de la couche 4 (recouvre E) \\ \hline
5 & Mise en place du dyke D (recoupe 4, pas 5) \\ \hline
6 & Dépôt de la couche 5 (recouvre D) \\ \hline
\end{tabularx}
\end{center}
\textbf{Conclusion} : après le dépôt des couches 1 à 3, la région a subi une tectonique (plissement et fracturation), puis une longue phase d'érosion, avant une nouvelle transgression marine (dépôt de 4), une intrusion magmatique (dyke D) et enfin le dépôt de la couche 5.
\end{exoresolu}

\begin{methode}[Exploiter les fossiles stratigraphiques pour dater une formation]
Pour dater relativement une couche grâce à ses fossiles :
\begin{enumerate}
\item \textbf{Identifier les fossiles} contenus dans la couche (nom du genre ou de l'espèce : ammonite, trilobite, foraminifère...).
\item \textbf{Vérifier qu'il s'agit de bons fossiles stratigraphiques} : un bon fossile stratigraphique présente quatre critères — large répartition géographique, grande abondance, durée d'existence courte (espèce à évolution rapide), facilement identifiable.
\item \textbf{Attribuer l'âge} : la couche a l'âge de la période d'existence de l'espèce fossile qu'elle contient. Si plusieurs fossiles sont présents, l'âge de la couche correspond à l'\emph{intervalle commun} d'existence de toutes les espèces (intersection des fourchettes).
\item \textbf{Corréler} : deux couches éloignées contenant la même association de fossiles stratigraphiques sont de même âge (\cle{principe d'identité paléontologique}), même si leur nature (faciès) est différente.
\end{enumerate}
\textbf{À retenir} : un fossile stratigraphique date la couche qui le contient, jamais les couches voisines directement — c'est ensuite la superposition qui permet d'encadrer les autres couches.
\end{methode}

\begin{exoresolu}[Datation d'une couche par une association de fossiles]
Dans une falaise du littoral camerounais, une couche de calcaire contient trois espèces fossiles dont les durées d'existence connues sont : espèce X : de $-160$ à $-150$ Ma ; espèce Y : de $-155$ à $-145$ Ma ; espèce Z : de $-170$ à $-152$ Ma.
\begin{enumerate}
\item Préciser l'intervalle d'âge possible de cette couche.
\item Une couche de grès située dans une autre région contient les espèces X et Z. Que peut-on en conclure ?
\end{enumerate}
\tcblower
\textbf{1. Intervalle d'âge de la couche calcaire.}\\
Les trois espèces ont vécu ensemble dans la couche : celle-ci s'est donc déposée pendant l'intervalle de temps où les trois espèces \emph{coexistaient}, c'est-à-dire l'\textbf{intersection} des trois fourchettes :
\begin{itemize}
\item borne la plus récente des débuts d'existence : $\max(-160, -155, -170) = -155$ Ma ;
\item borne la plus ancienne des fins d'existence : $\min(-150, -145, -152) = -152$ Ma.
\end{itemize}
La couche calcaire a donc un âge compris entre $-155$ Ma et $-152$ Ma. C'est l'application du \cle{principe d'identité paléontologique} restreint à une association d'espèces.

\textbf{2. Corrélation avec la couche de grès.}\\
La couche de grès contient X et Z : son âge est compris entre $-160$ Ma (début de X) et $-152$ Ma (fin de Z). Les intervalles d'âge des deux couches se recouvrent (zone commune : $-155$ à $-152$ Ma) : les deux couches \textbf{peuvent être contemporaines} (au moins partiellement), bien que leurs natures soient différentes (calcaire et grès = faciès différents, donc milieux de dépôt différents à la même époque).\\
\textbf{Conclusion} : la couche calcaire date de $-155$ à $-152$ Ma ; la couche de grès lui est au moins partiellement synchrone, ce qui illustre qu'un même âge peut correspondre à des faciès sédimentaires différents.
\end{exoresolu}

\begin{methode}[Reconnaître et exploiter une discordance stratigraphique]
\begin{enumerate}
\item \textbf{Repérer la discordance} sur la coupe : des couches inclinées ou plissées sont tronquées par une surface plane (ou ondulée) d'érosion, sur laquelle reposent des couches horizontales.
\item \textbf{Identifier le type} : discordance angulaire (couches inférieures inclinées/plissées), discordance d'érosion ou paraconcordance (couches parallèles mais séparées par une surface d'érosion), discordance de non-dépôt (lacune de sédimentation sans érosion visible).
\item \textbf{Dater la discordance par cadrage} : elle est postérieure à la plus récente des couches tronquées et antérieure à la plus ancienne des couches sus-jacentes.
\item \textbf{Interpréter l'histoire} : une discordance angulaire traduit la succession : dépôt des couches inférieures $\rightarrow$ tectonique (plissement) $\rightarrow$ érosion (mise à l'air libre) $\rightarrow$ nouvelle transgression et dépôt des couches supérieures. Elle marque une \cle{lacune} dans l'enregistrement sédimentaire.
\end{enumerate}
\textbf{Réflexe} : la présence éventuelle d'un conglomérat de base au pied des couches sus-jacentes confirme la surface d'érosion.
\end{methode}

\begin{popfigure}
\begin{center}
\begin{tikzpicture}[scale=1.0]
% Couches redressées S (schistes) : bandes verticales tronquées
\fill[poppurpleL] (0,0) rectangle (1,2.4);
\fill[poppurple!50] (1,0) rectangle (2,2.4);
\fill[poppurpleL] (2,0) rectangle (3,2.4);
\fill[poppurple!50] (3,0) rectangle (4,2.4);
\fill[poppurpleL] (4,0) rectangle (5,2.4);
\fill[poppurple!50] (5,0) rectangle (6,2.4);
% Surface d'érosion
\draw[popink,very thick] (0,2.4) -- (6,2.4);
% Conglomérat de base G
\fill[popgoldL] (0,2.4) rectangle (6,2.9);
\foreach \x in {0.4,1.1,1.8,2.6,3.3,4.1,4.8,5.5} {\fill[popdark] (\x,2.65) circle (0.07);}
% Sables horizontaux H
\fill[poporangeL] (0,2.9) rectangle (6,3.9);
\draw[popmuted] (0,2.9) -- (6,2.9);
% Légendes pointées
\draw[popmuted,thin] (0.5,1.2) -- (-0.4,1.2) node[left,font=\small]{S : schistes redressés (trilobites, Primaire)};
\draw[popmuted,thin] (5.5,2.4) -- (6.4,2.1) node[right,font=\small]{surface d'érosion};
\draw[popmuted,thin] (5.5,2.65) -- (6.4,2.75) node[right,font=\small]{G : conglomérat de base};
\draw[popmuted,thin] (5.5,3.4) -- (6.4,3.5) node[right,font=\small]{H : sables horizontaux (ammonites, Secondaire)};
\end{tikzpicture}
\end{center}
\legende{Discordance angulaire : les couches redressées S sont tronquées par une surface d'érosion surmontée d'un conglomérat de base G puis de couches horizontales H. La discordance enregistre une lacune : plissement puis érosion entre le Primaire et le Secondaire.}
\end{popfigure}

\begin{exoresolu}[Interprétation d'une discordance angulaire]
Sur une coupe, on observe des couches de schistes S fortement redressées (presque verticales), tronquées par une surface plane surmontée d'un conglomérat G puis de couches horizontales de sable H. Les schistes S contiennent des trilobites du Primaire ; les sables H contiennent des ammonites du Secondaire.
\begin{enumerate}
\item Nommer la structure observée entre S et l'ensemble G--H.
\item Décrire les étapes de l'histoire géologique qui ont conduit à cette structure.
\end{enumerate}
\tcblower
\textbf{1. Nature de la structure.}\\
Des couches redressées (S) sont tronquées par une surface d'érosion sur laquelle reposent des couches horizontales (G puis H) : il s'agit d'une \cle{discordance angulaire}, matérialisée par la surface d'érosion. Le conglomérat G est le \emph{conglomérat de base} de la série sus-jacente, témoin du remaniement des terrains sous-jacents.

\textbf{2. Histoire géologique.} Les étapes, de la plus ancienne à la plus récente, sont :
\begin{enumerate}
\item \textbf{Dépôt des schistes S} en milieu sédimentaire, au Primaire (présence de trilobites, bons fossiles stratigraphiques de cette ère) ;
\item \textbf{Tectonique} : une phase de compression plisse et redresse les couches S (elles passent de l'horizontale à la verticale) — le plissement est postérieur au dépôt de S (principe de recoupement : il affecte S) ;
\item \textbf{Érosion} : la région, soulevée, est mise à l'air libre ; l'érosion arase les couches redressées et crée la surface d'érosion — cette érosion est postérieure au plissement et antérieure au dépôt de G ;
\item \textbf{Transgression marine au Secondaire} : la mer recouvre la surface érodée ; les débris remaniés forment le conglomérat de base G, puis se déposent les sables H contenant des ammonites (fossiles stratigraphiques du Secondaire).
\end{enumerate}
\textbf{Conclusion} : la discordance angulaire enregistre une longue lacune (fin du Primaire à début du Secondaire) correspondant à une phase de plissement suivie d'érosion, avant la reprise de la sédimentation au Secondaire.
\end{exoresolu}

\begin{methode}[Exploiter un document du type « coupe + tableau de fossiles » au Probatoire]
Face à un sujet complet (coupe géologique + données paléontologiques), adopte une stratégie de copie efficace :
\begin{enumerate}
\item \textbf{Lire l'énoncé en entier} et souligner les verbes de consigne : « classer », « justifier », « reconstituer », « dater ».
\item \textbf{Annoter la coupe} au brouillon : numéroter les couches de bas en haut, repérer failles, plis, filons, discordances.
\item \textbf{Construire un tableau à deux colonnes} : « Événement » / « Justification (principe utilisé) ». Chaque ligne doit citer un principe nommément.
\item \textbf{Rédiger la frise chronologique finale} du plus ancien au plus récent, en précisant les éventuelles indéterminations (« on ne peut pas départager F et P »).
\item \textbf{Vérifier la cohérence} : aucun événement récent ne doit se trouver sous un événement ancien ; tout élément recoupant doit être postérieur à tout élément recoupé.
\end{enumerate}
\end{methode}

\begin{exoresolu}[Sujet de synthèse type Probatoire]
Une coupe géologique montre du bas vers le haut : couche 1 (calcaire à trilobites), couche 2 (grès), toutes deux plissées ; faille F décalant 1 et 2 ; surface d'érosion E ; couche 3 (calcaire à ammonites) horizontale reposant sur E ; dyke basaltique D recoupant 1, 2, F, E et 3.
\begin{enumerate}
\item Établir la chronologie relative de tous les événements.
\item Que peut-on dire des âges relatifs des couches 1 et 3 ?
\end{enumerate}
\tcblower
\textbf{1. Chronologie relative.}
\begin{itemize}
\item \textbf{Dépôts de 1 puis 2} : principe de superposition, 1 plus ancienne que 2.
\item \textbf{Plissement P} : il affecte 1 et 2, donc postérieur à 2 ; il est tronqué par E, donc antérieur à E.
\item \textbf{Faille F} : elle décale 1 et 2, donc postérieure à 2 ; elle est tronquée par E, donc antérieure à E. (Impossible de départager P et F avec ces données.)
\item \textbf{Érosion E} : postérieure à P et F, antérieure à 3.
\item \textbf{Dépôt de 3} : recouvre E, donc postérieur à E.
\item \textbf{Dyke D} : il recoupe \emph{tous} les éléments, y compris la couche 3 : c'est l'événement \textbf{le plus récent} de la coupe (principe de recoupement).
\end{itemize}
Frise chronologique (ancien $\rightarrow$ récent) :
\[ \text{dépôt 1} \rightarrow \text{dépôt 2} \rightarrow \text{pli P et faille F} \rightarrow \text{érosion E} \rightarrow \text{dépôt 3} \rightarrow \text{dyke D} \]

\textbf{2. Âges relatifs de 1 et 3.}\\
La couche 1 contient des trilobites (fossiles stratigraphiques du Primaire) et la couche 3 des ammonites (fossiles stratigraphiques du Secondaire) : d'après le \cle{principe d'identité paléontologique}, la couche 1 est du Primaire et la couche 3 du Secondaire. La couche 1 est donc \emph{nettement plus ancienne} que la couche 3, et la discordance E correspond à une lacune couvrant la transition Primaire--Secondaire (plissement puis érosion).\\
\textbf{Conclusion} : la région a connu une sédimentation primaire, une phase tectonique avec érosion, puis une reprise de la sédimentation au Secondaire, avant une intrusion magmatique tardive (dyke D).
\end{exoresolu}

\begin{attention}[Les 5 erreurs qui coûtent des points au Probatoire]
\begin{enumerate}
\item \textbf{Inverser le sens de la superposition} : la couche la plus ancienne est en \emph{bas}, la plus récente en \emph{haut} — et seulement si la série n'est pas renversée. Oublier cette condition (« série non renversée ») fait perdre le point de justification.
\item \textbf{Confondre « recoupe » et « est recoupé »} : une faille ou un filon qui \emph{recoupe} une couche est \emph{plus récent} qu'elle. Beaucoup d'élèves écrivent l'inverse. Réflexe : « celui qui coupe est le plus jeune ».
\item \textbf{Dater un événement sans cadrage double} : pour un pli ou une faille, il faut donner les deux bornes (couche la plus récente affectée ET couche la plus ancienne non affectée). Une seule borne ne suffit pas pour situer l'événement.
\item \textbf{Oublier de nommer les principes} : écrire « la couche 2 est au-dessus donc plus récente » sans citer le \emph{principe de superposition} fait perdre le point de justification. Chaque étape doit être rattachée à un principe nommé (superposition, continuité, recoupement, identité paléontologique).
\item \textbf{Mal exploiter une association de fossiles} : avec plusieurs espèces, l'âge de la couche est l'\emph{intersection} des durées d'existence, pas la réunion. Prendre la fourchette la plus large (réunion) est une erreur classique ; prendre aussi un fossile quelconque (longue durée d'existence) comme fossile stratigraphique sans vérifier les critères (répartition large, abondance, durée courte) est sanctionné.
\end{enumerate}
\end{attention}

\newpage

\coursec{Exercices}

\courssub{Je vérifie mes connaissances}

\begin{exercice}[QCM : les principes de la stratigraphie]
Pour chaque question, une seule réponse est correcte. Recopie la lettre de la bonne réponse.

\textbf{Question 1.} Dans une série de couches sédimentaires non renversées, une couche est d'autant plus ancienne qu'elle est :
\begin{enumerate}[label=\alph*)]
\item plus épaisse ;
\item plus profonde ;
\item plus riche en fossiles ;
\item plus proche de la surface.
\end{enumerate}

\textbf{Question 2.} Un filon de basalte recoupe trois couches de calcaire. Le basalte est :
\begin{enumerate}[label=\alph*)]
\item plus ancien que les trois couches ;
\item contemporain des trois couches ;
\item plus récent que les trois couches ;
\item impossible à dater relativement.
\end{enumerate}

\textbf{Question 3.} Le principe de continuité affirme qu'une couche sédimentaire :
\begin{enumerate}[label=\alph*)]
\item a le même âge sur toute son étendue ;
\item contient toujours les mêmes fossiles ;
\item ne peut jamais être plissée ;
\item est toujours horizontale.
\end{enumerate}

\textbf{Question 4.} Un bon fossile stratigraphique est une espèce :
\begin{enumerate}[label=\alph*)]
\item à longue durée d'existence et à répartition restreinte ;
\item à courte durée d'existence et à large répartition géographique ;
\item encore vivante aujourd'hui ;
\item présente dans une seule localité.
\end{enumerate}
\end{exercice}

\begin{exercice}[Vrai ou faux ?]
Indique si chacune des affirmations suivantes est vraie ou fausse, puis justifie ta réponse en une ou deux phrases.
\begin{enumerate}
\item Une faille qui recoupe une couche est plus ancienne que cette couche.
\item Deux couches situées dans deux régions différentes et contenant les mêmes fossiles stratigraphiques peuvent être considérées comme de même âge.
\item Une discordance angulaire témoigne d'une phase de plissement suivie d'une érosion avant le dépôt des couches supérieures.
\item L'huître \emph{Crassostrea gigas}, encore vivante aujourd'hui, est un excellent fossile stratigraphique.
\end{enumerate}
\end{exercice}

\begin{exercice}[Définitions]
Définis en une ou deux phrases chacun des termes suivants :
\begin{enumerate}
\item strate (ou couche sédimentaire) ;
\item fossile stratigraphique ;
\item discordance stratigraphique ;
\item datation relative.
\end{enumerate}
\end{exercice}

\begin{exercice}[Application directe des principes]
Sur une coupe géologique, on observe du bas vers le haut : une couche de grès (couche A), une couche de calcaire (couche B) et une couche de marne (couche C). L'ensemble est recoupé par une faille F, et un filon de granite G recoupe la faille F.
\begin{enumerate}
\item Classe les couches A, B et C de la plus ancienne à la plus récente. Quel principe utilises-tu ?
\item La faille F est-elle antérieure ou postérieure à la couche C ? Justifie.
\item Classe les cinq éléments (A, B, C, F, G) du plus ancien au plus récent.
\end{enumerate}
\end{exercice}

\courssub{Je m'entraîne}

\begin{exercice}[Deux affleurements distants]
Deux affleurements situés de part et d'autre de la Sanaga, distants d'environ \SI{50}{km}, présentent des couches de natures différentes : du calcaire à l'affleurement 1 et du grès fossilifère à l'affleurement 2. Les deux couches contiennent néanmoins les mêmes espèces d'ammonites.
\begin{enumerate}
\item Peut-on considérer que ces deux couches se sont déposées à la même époque ? Justifie ta réponse.
\item Nomme le principe de datation relative mis en jeu.
\item Explique pourquoi la nature différente des roches ne contredit pas ta conclusion.
\end{enumerate}
\end{exercice}

\begin{exercice}[Les étapes d'une discordance angulaire]
Une coupe géologique réalisée lors de travaux routiers montre des couches plissées de schistes surmontées par des couches horizontales de conglomérat puis de calcaire ; la surface de contact entre les deux ensembles est irrégulière.
\begin{enumerate}
\item Comment appelle-t-on la surface de contact observée ?
\item Décris et ordonne les quatre étapes qui ont conduit à la formation de cette discordance (dépôt des couches inférieures, plissement, érosion, transgression marine avec nouveau dépôt).
\item Que représente le conglomérat de base au sommet de la surface irrégulière ?
\end{enumerate}
\end{exercice}

\begin{exercice}[Le bon fossile stratigraphique]
On a découvert dans une couche de marne trois espèces fossiles :
\begin{itemize}
\item espèce 1 : un coquillage ayant vécu pendant 200 millions d'années, présent sur tous les continents ;
\item espèce 2 : une ammonite ayant vécu pendant 2 millions d'années, connue en Afrique, en Europe et en Amérique du Sud ;
\item espèce 3 : un corail connu uniquement dans cette région, ayant vécu pendant 1 million d'années.
\end{itemize}
\begin{enumerate}
\item Rappelle les trois critères d'un bon fossile stratigraphique.
\item Quelle espèce est la meilleure pour dater la couche de marne ? Justifie.
\item Explique pourquoi un coquillage encore vivant aujourd'hui ne peut pas servir à dater une formation géologique.
\end{enumerate}
\end{exercice}

\begin{exercice}[Faille et filon : qui est le plus récent ?]
Sur une carte géologique simplifiée d'une région de l'Adamaoua, on observe quatre couches sédimentaires horizontales numérotées de 1 (la plus profonde) à 4 (la plus superficielle). Une faille F recoupe les couches 1 à 3 mais s'arrête sous la couche 4. Un filon de dolérite D recoupe les couches 1 à 4.
\begin{enumerate}
\item Classe les couches 1 à 4 de la plus ancienne à la plus récente en citant le principe utilisé.
\item Situe la faille F dans le temps par rapport aux couches 1, 3 et 4. Justifie.
\item Situe le filon D par rapport à l'ensemble.
\item Établis la chronologie complète des six éléments (couches 1 à 4, faille F, filon D) du plus ancien au plus récent.
\end{enumerate}
\end{exercice}

\courssub{Je me prépare au Probatoire}

\begin{exercice}[Type Probatoire : la carrière de Nomayos]
Lors de l'exploitation d'une carrière de sable au sud de Yaoundé, une coupe géologique a été dégagée. On y observe, du bas vers le haut :
\begin{itemize}
\item une couche de grès (couche 1) ;
\item une couche de calcaire à ammonites (couche 2) ;
\item une couche de marne (couche 3) ;
\item une couche de sable (couche 4) ;
\item une couche d'argile (couche 5).
\end{itemize}
L'ensemble des cinq couches est horizontal. Une faille F recoupe les couches 1 à 5. Un filon de basalte B recoupe les couches 1 à 4 ainsi que la faille F, mais s'arrête sous la couche 5.

\textbf{Questions :}
\begin{enumerate}
\item Définis le terme « strate » et rappelle le principe de superposition.
\item Classe les couches 1 à 5 de la plus ancienne à la plus récente en justifiant ta réponse.
\item En appliquant le principe de recoupement, montre que la faille F est postérieure à la couche 5.
\item Montre que le filon de basalte B est à la fois postérieur à la faille F et antérieur à la couche 5. Que peux-tu en conclure sur l'ordre de F et B ?
\item Établis la chronologie complète des sept éléments (couches 1 à 5, faille F, filon B) du plus ancien au plus récent, sous forme d'une liste ordonnée.
\item La couche 2 contient des ammonites. Rappelle les trois critères qui font des ammonites de bons fossiles stratigraphiques et explique l'intérêt de leur présence pour comparer cette carrière à un autre site du pays.
\end{enumerate}
\end{exercice}

\begin{exercice}[Type Probatoire : la coupe de la falaise de Limbé]
Une falaise du littoral de Limbé, au pied du Mont Cameroun, présente la succession suivante, du bas vers le haut :
\begin{itemize}
\item des couches de schistes fortement plissées (ensemble A) ;
\item une surface érodée irrégulière portant des galets ;
\item un conglomérat de base (couche B) ;
\item des couches horizontales de calcaire (couche C) puis de marne (couche D) ;
\item une coulée de lave récente (couche E) qui recouvre la marne.
\end{itemize}

\textbf{Questions :}
\begin{enumerate}
\item Nomme la structure représentée par le contact entre l'ensemble A et les couches B--C--D.
\item Décris, dans l'ordre chronologique, les quatre étapes qui ont conduit à la formation de cette structure (dépôt, plissement, érosion, transgression).
\item La coulée de lave E repose sur la couche D sans la recouper. Peut-on appliquer le principe de recoupement pour la dater ? Propose un raisonnement permettant de situer E par rapport à D.
\item Établis la chronologie complète des événements géologiques de cette région, du plus ancien au plus récent.
\item Construis une frise chronologique (flèche du temps orientée vers le présent) sur laquelle tu placeras dans l'ordre : dépôt des schistes, plissement, érosion, dépôt du conglomérat, dépôt du calcaire, dépôt de la marne, émission de la coulée de lave.
\end{enumerate}
\end{exercice}

\begin{exercice}[Type Probatoire : synthèse sur la région de Bafoussam]
Une coupe géologique dressée dans la région de Bafoussam montre :
\begin{itemize}
\item à la base, des couches plissées de grès et de schiste (couches 1 et 2) contenant des fossiles de trilobites ;
\item une surface d'érosion irrégulière ;
\item au-dessus, trois couches horizontales : un conglomérat (couche 3), un calcaire à ammonites (couche 4) et une marne (couche 5) ;
\item une faille F qui recoupe les couches 1 à 4 mais s'arrête sous la couche 5 ;
\item un filon de granite G qui recoupe uniquement les couches 1 et 2 et est tronqué par la surface d'érosion.
\end{itemize}

\textbf{Questions :}
\begin{enumerate}
\item Rappelle les trois principes fondamentaux de la stratigraphie en les énonçant précisément.
\item Montre, en justifiant chaque étape, que le filon G est antérieur à la phase d'érosion, alors que la faille F lui est postérieure.
\item Situe la faille F dans le temps par rapport aux couches 4 et 5.
\item Les couches 1 et 2 contiennent des trilobites et la couche 4 des ammonites. Que peut-on dire de l'âge relatif de ces couches ? Quel principe paléontologique utilises-tu ?
\item Établis la chronologie complète de tous les événements géologiques de cette région (dépôts, plissement, intrusion du granite, érosion, faille) du plus ancien au plus récent.
\item Rédige un court paragraphe (5 à 8 lignes) racontant l'histoire géologique de cette région en langage courant, comme si tu l'expliquais à un élève de Seconde.
\end{enumerate}
\end{exercice}

\newpage

\coursec{Corrigés des exercices}

\begin{corrige}[Exercice 1 — QCM : les principes de la stratigraphie]

\textbf{Question 1 : réponse b).} D'après le \cle{principe de superposition}, dans une série non renversée, chaque couche est plus ancienne que celle qui la recouvre : plus une couche est profonde, plus elle est ancienne. L'épaisseur (a) dépend des conditions de dépôt, pas de l'âge ; la richesse en fossiles (c) n'a aucun lien avec l'ancienneté ; la proximité de la surface (d) caractérise au contraire les couches récentes.

\textbf{Question 2 : réponse c).} D'après le \cle{principe de recoupement}, une structure (filon, faille) qui recoupe des couches est postérieure à toutes les couches qu'elle traverse : le basalte est donc plus récent que les trois couches de calcaire.

\textbf{Question 3 : réponse a).} Le \cle{principe de continuité} (ou d'identité paléontologique appliqué à la strate) affirme qu'une même couche a le même âge sur toute son étendue, même si sa nature ou son épaisseur varient latéralement. Les réponses b), c) et d) sont fausses : les fossiles peuvent varier avec le milieu, et une couche peut être plissée ou basculée après son dépôt.

\textbf{Question 4 : réponse b).} Un bon \cle{fossile stratigraphique} est une espèce à \cle{courte durée d'existence} (précision de la datation), à \cle{large répartition géographique} (corrélation entre régions éloignées) et \cle{abondante} (facile à trouver). Les réponses a), c) et d) décrivent des espèces inutilisables pour dater finement ou corréler.
\end{corrige}

\begin{corrige}[Exercice 2 — Vrai ou faux ?]

\begin{enumerate}
\item \textbf{Faux.} D'après le principe de recoupement, une faille qui recoupe une couche est \emph{postérieure} (plus récente) à cette couche : la couche devait exister avant de pouvoir être fracturée.
\item \textbf{Vrai.} C'est le \cle{principe d'identité paléontologique} : deux couches contenant les mêmes fossiles stratigraphiques sont considérées comme de même âge, même si elles sont très éloignées et de natures différentes.
\item \textbf{Vrai.} Une \cle{discordance angulaire} se forme en quatre étapes : dépôt des couches inférieures, plissement (tectonique), érosion qui arase l'ensemble, puis transgression marine avec dépôt de nouvelles couches horizontales. L'angle entre les deux séries témoigne bien du plissement antérieur.
\item \textbf{Faux.} \emph{Crassostrea gigas} vit encore aujourd'hui : sa durée d'existence est très longue (plusieurs millions d'années). Or un bon fossile stratigraphique doit avoir une courte durée d'existence ; cette huître ne permet donc aucune datation précise.
\end{enumerate}
\end{corrige}

\begin{corrige}[Exercice 3 — Définitions]

\begin{enumerate}
\item \textbf{Strate (couche sédimentaire)} : ensemble de sédiments déposés de façon continue dans des conditions sensiblement identiques, limité par deux surfaces de sédimentation appelées \cle{joints de strates}.
\item \textbf{Fossile stratigraphique} : reste ou trace d'une espèce ayant vécu pendant une courte période, répandue sur une vaste zone géographique et abondante, permettant de dater relativement la couche qui le contient.
\item \textbf{Discordance stratigraphique} : surface d'érosion séparant deux séries de couches dont les orientations ou les conditions de dépôt diffèrent ; elle traduit une interruption de la sédimentation (hiatus) liée à des événements tectoniques et/ou érosifs.
\item \textbf{Datation relative} : méthode consistant à classer les événements géologiques les uns par rapport aux autres (antériorité, postériorité, contemporanéité) sans leur attribuer d'âge chiffré en années.
\end{enumerate}
\end{corrige}

\begin{corrige}[Exercice 4 — Application directe des principes]

\begin{enumerate}
\item Ordre de la plus ancienne à la plus récente : \textbf{A (grès) < B (calcaire) < C (marne)}. On applique le \cle{principe de superposition} : dans une série non renversée, toute couche est plus récente que celle qui la supporte.
\item La faille F recoupe la couche C : d'après le \cle{principe de recoupement}, F est \textbf{postérieure} à C (la couche devait être en place pour être fracturée).
\item Le filon G recoupe la faille F, donc G est postérieur à F. Chronologie complète, du plus ancien au plus récent :
\[ \text{A} \;<\; \text{B} \;<\; \text{C} \;<\; \text{F} \;<\; \text{G} \]
\end{enumerate}
\end{corrige}

\begin{corrige}[Exercice 5 — Deux affleurements distants]

\begin{enumerate}
\item \textbf{Oui}, on peut considérer que ces deux couches se sont déposées à la même époque : elles contiennent les mêmes espèces d'ammonites, qui sont d'excellents fossiles stratigraphiques.
\item C'est le \cle{principe d'identité paléontologique} : deux couches contenant les mêmes fossiles stratigraphiques sont de même âge, quelle que soit la distance qui les sépare.
\item La nature d'une roche sédimentaire dépend du \emph{milieu de dépôt}, pas de l'époque : à une même époque, la mer pouvait déposer du calcaire en zone claire et peu profonde (affleurement 1) et du grès plus près d'un apport détritique (affleurement 2). On parle de \cle{faciès} différents mais contemporains. La différence lithologique ne contredit donc pas la contemporanéité établie par les fossiles.
\end{enumerate}
\end{corrige}

\begin{corrige}[Exercice 6 — Les étapes d'une discordance angulaire]

\begin{enumerate}
\item La surface de contact irrégulière entre les schistes plissés et les couches horizontales sus-jacentes est une \cle{surface de discordance angulaire} (ou surface d'érosion).
\item Les quatre étapes, dans l'ordre chronologique :
\begin{enumerate}
\item \textbf{Dépôt} des couches de schistes à l'horizontale, dans un milieu sédimentaire ;
\item \textbf{Plissement} de ces couches sous l'effet de contraintes tectoniques (orogenèse) ;
\item \textbf{Érosion} de l'ensemble plissé, souvent mis à l'air libre, qui arase les plis et crée la surface irrégulière ;
\item \textbf{Transgression marine} : la mer envahit la surface érodée et de nouvelles couches horizontales (conglomérat puis calcaire) se déposent en discordance.
\end{enumerate}
\item Le conglomérat de base est le \textbf{premier dépôt de la transgression} : il rassemble les galets arrachés à la surface érodée et remaniés par la mer. Il marque la reprise de la sédimentation et scelle la discordance.
\end{enumerate}
\end{corrige}

\begin{corrige}[Exercice 7 — Le bon fossile stratigraphique]

\begin{enumerate}
\item Les trois critères d'un bon fossile stratigraphique :
\begin{itemize}
\item \textbf{courte durée d'existence} de l'espèce (datation précise) ;
\item \textbf{large répartition géographique} (corrélations possibles entre régions éloignées) ;
\item \textbf{grande abondance} (facile à trouver dans les couches).
\end{itemize}
\item C'est l'\textbf{espèce 2} (l'ammonite) : courte durée d'existence (2 millions d'années) et large répartition (trois continents). L'espèce 1 a vécu 200 millions d'années : elle ne permet qu'une datation très imprécise. L'espèce 3, bien que de courte durée, n'est connue que localement : elle ne permet aucune corrélation avec d'autres régions.
\item Un coquillage encore vivant aujourd'hui a une durée d'existence qui s'étend jusqu'au présent, donc très longue : sa présence dans une couche indique seulement que la couche s'est déposée \emph{à un moment quelconque} de cette longue période. Il ne permet ni datation précise ni corrélation fine entre couches.
\end{enumerate}
\end{corrige}

\begin{corrige}[Exercice 8 — Faille et filon : qui est le plus récent ?]

\begin{enumerate}
\item Ordre de la plus ancienne à la plus récente : \textbf{1 < 2 < 3 < 4}, par application du \cle{principe de superposition} (série horizontale non renversée).
\item La faille F recoupe les couches 1 à 3 : elle est donc \textbf{postérieure à la couche 3} (principe de recoupement). Mais elle s'arrête \emph{sous} la couche 4, qui la scelle : elle est donc \textbf{antérieure à la couche 4}. Conclusion : F s'est produite entre le dépôt de 3 et le dépôt de 4.
\item Le filon D recoupe les couches 1 à 4 : il est \textbf{postérieur à toutes les couches}, donc aussi postérieur à la faille F. C'est l'élément le plus récent de la coupe.
\item Chronologie complète, du plus ancien au plus récent :
\[ 1 \;<\; 2 \;<\; 3 \;<\; \text{F} \;<\; 4 \;<\; \text{D} \]
\end{enumerate}

\begin{attention}[Point de vigilance]
L'argument décisif est double : la faille \emph{recoupe} 1 à 3 (postériorité) mais est \emph{scellée} par 4 (antériorité). Une structure qui s'arrête sous une couche est antérieure à cette couche : c'est le raisonnement inverse du recoupement, à savoir formuler explicitement au Probatoire.
\end{attention}
\end{corrige}

\begin{corrige}[Exercice 9 — Type Probatoire : la carrière de Nomayos]

\begin{enumerate}
\item Une \textbf{strate} est un ensemble de sédiments déposés de façon continue dans des conditions sensiblement identiques, limité par deux joints de strates. \textbf{Principe de superposition} : dans une série de couches sédimentaires non renversées, toute couche est plus récente que celle qui la supporte et plus ancienne que celle qui la recouvre.
\item Les couches étant horizontales et non renversées, le principe de superposition donne, de la plus ancienne à la plus récente :
\[ 1 \;<\; 2 \;<\; 3 \;<\; 4 \;<\; 5 \]
\item La faille F recoupe les couches 1 à 5. D'après le principe de recoupement, une structure qui recoupe une couche lui est postérieure ; F recoupant la couche 5, F est \textbf{postérieure à la couche 5}.
\item Le filon B recoupe la faille F : il est donc postérieur au \textbf{premier jeu} de la faille F. Mais B s'arrête sous la couche 5, qui le recouvre sans être recoupée : B est donc \textbf{antérieur à la couche 5}. Or F recoupe la couche 5 (question 3) : cela prouve un \textbf{rejeu de la faille F postérieur au dépôt de la couche 5}. La chronologie relative est donc : premier jeu de F $<$ B $<$ 5 $<$ rejeu de F. Au niveau Première, on retient la chronologie simplifiée où l'événement « faille F » est placé à sa position finale (rejeu), soit après 5.
\item Chronologie complète, du plus ancien au plus récent (version simplifiée attendue) :
\[ 1 \;<\; 2 \;<\; 3 \;<\; 4 \;<\; \text{B} \;<\; 5 \;<\; \text{F} \]
\item Les trois critères : courte durée d'existence, large répartition géographique, grande abondance. Les ammonites de la couche 2 permettent, par le \cle{principe d'identité paléontologique}, d'affirmer que toute couche d'un autre site camerounais contenant les mêmes espèces d'ammonites est de même âge que la couche 2 : on peut ainsi \textbf{corréler} la carrière de Nomayos avec d'autres coupes du pays malgré l'éloignement et les différences de faciès.
\end{enumerate}

\textbf{Barème indicatif (sur 10) :} définitions 1 pt ; classement des couches justifié 1,5 pt ; postériorité de F 1,5 pt ; double argument sur B et conclusion F > B (avec mention du rejeu) 2 pts ; chronologie complète 2 pts ; critères et intérêt des ammonites 2 pts.
\end{corrige}

\begin{attention}[Exercice 9 — Points de vigilance]
\begin{itemize}
\item La question 4 est un piège classique : B recoupe F (donc B > F \emph{premier jeu}) mais B < 5 (scellement) et F > 5 (recoupement de 5). La contradiction apparente se lève par le \textbf{rejeu de la faille} : F a bougé une première fois (recoupé par B), puis une seconde fois après le dépôt de 5 (recoupant 5). La rédaction attendue doit expliciter ce rejeu.
\item Toujours citer le nom du principe utilisé : « principe de superposition », « principe de recoupement » — une justification sans nom de principe perd des points.
\item Pour la chronologie, ne pas oublier d'insérer les événements (intrusion B, faille F) entre les numéros de couches.
\end{itemize}
\end{attention}

\begin{corrige}[Exercice 10 — Type Probatoire : la coupe de la falaise de Limbé]

\begin{enumerate}
\item Le contact entre l'ensemble plissé A et les couches horizontales B--C--D, matérialisé par une surface érodée, est une \cle{discordance angulaire}.
\item Les quatre étapes chronologiques :
\begin{enumerate}
\item \textbf{dépôt} des couches de schistes à l'horizontale ;
\item \textbf{plissement} de ces couches par des contraintes tectoniques ;
\item \textbf{érosion} de l'édifice plissé (mise à l'air libre, arasement) créant la surface irrégulière à galets ;
\item \textbf{transgression marine} : dépôt en discordance du conglomérat de base (B), puis du calcaire (C) et de la marne (D).
\end{enumerate}
\item \textbf{Non}, le principe de recoupement ne s'applique pas directement car E ne recoupe pas D : E \emph{repose} sur D. On raisonne alors par \textbf{superposition} : la coulée E recouvre la marne D, elle est donc postérieure à D. (On peut aussi invoquer le principe de recoupement au sens large : une coulée qui s'étale sur une surface est postérieure à tout ce qui constitue cette surface.)
\item Chronologie complète, du plus ancien au plus récent :
\[ \text{dépôt de A} < \text{plissement de A} < \text{érosion} < \text{B} < \text{C} < \text{D} < \text{E} \]
\item Frise chronologique (le temps s'écoule vers la droite, vers le présent) :

\begin{popfigure}
\begin{tikzpicture}[scale=1]
\draw[-{Stealth},very thick,popink] (0,0) -- (15.2,0);
\node[below] at (0.2,-0.05) {\small passé};
\node[below] at (14.9,-0.05) {\small présent};
\foreach \x/\t in {0.8/{dépôt des schistes}, 3.1/{plissement}, 5.2/{érosion}, 7.3/{dépôt du conglomérat}, 9.4/{dépôt du calcaire}, 11.5/{dépôt de la marne}, 13.8/{coulée de lave}}{
  \fill[poporange] (\x,0) circle (2.2pt);
  \node[above,align=center,font=\scriptsize,text width=1.9cm] at (\x,0.15) {\t};
}
\end{tikzpicture}
\legende{Frise chronologique des événements géologiques de la falaise de Limbé, du plus ancien (à gauche) au plus récent (à droite).}
\end{popfigure}
\end{enumerate}

\textbf{Barème indicatif (sur 10) :} nom de la structure 1 pt ; quatre étapes ordonnées 3 pts ; raisonnement sur E 2 pts ; chronologie complète 2 pts ; frise correctement orientée et ordonnée 2 pts.
\end{corrige}

\begin{attention}[Exercice 10 — Points de vigilance]
\begin{itemize}
\item Ne pas confondre « recouvrir » et « recouper » : la coulée E \emph{recouvre} D, on utilise alors la superposition, pas le recoupement. Le jury attend cette distinction.
\item Sur la frise, la flèche du temps doit être orientée vers le présent et les sept événements dans l'ordre exact ; le plissement et l'érosion sont des \emph{événements}, pas des couches.
\end{itemize}
\end{attention}

\begin{corrige}[Exercice 11 — Type Probatoire : synthèse sur la région de Bafoussam]

\begin{enumerate}
\item Les trois principes fondamentaux de la stratigraphie :
\begin{itemize}
\item \textbf{principe de superposition} : dans une série non renversée, une couche est plus récente que celle qui la supporte et plus ancienne que celle qui la recouvre ;
\item \textbf{principe de continuité} : une même couche a le même âge sur toute son étendue, même si son faciès varie latéralement ;
\item \textbf{principe de recoupement} : une structure géologique (faille, filon) qui recoupe des couches est postérieure à toutes les couches recoupées.
\end{itemize}
\item \textbf{Filon G antérieur à l'érosion} : G recoupe les couches 1 et 2 mais est \emph{tronqué} par la surface d'érosion : l'érosion a arasé le sommet du filon, donc l'intrusion de G est antérieure à la phase d'érosion. \textbf{Faille F postérieure à l'érosion} : F recoupe les couches 3 et 4, qui sont \emph{postérieures} à la surface d'érosion (elles reposent au-dessus) ; F est donc postérieure au dépôt de ces couches, et par conséquent postérieure à l'érosion.
\item F recoupe la couche 4 : F est \textbf{postérieure à 4} (principe de recoupement). F s'arrête sous la couche 5, qui la scelle : F est \textbf{antérieure à 5}. Donc F s'est produite entre le dépôt de 4 et le dépôt de 5.
\item Les couches 1 et 2 (à trilobites) sont situées \emph{sous} la discordance, la couche 4 (à ammonites) \emph{au-dessus} : par superposition, 1 et 2 sont nettement plus anciennes que 4. De plus, par le \cle{principe d'identité paléontologique}, les trilobites caractérisent des terrains anciens (Paléozoïque) et les ammonites des terrains plus récents (Mésozoïque) : les faunes différentes confirment que ces couches appartiennent à des époques très éloignées, séparées par un long hiatus (la discordance).
\item Chronologie complète, du plus ancien au plus récent :
\[ \text{dépôt de 1} < \text{dépôt de 2} < \text{plissement} < \text{intrusion de G} < \text{érosion} < \text{dépôt de 3} < \text{dépôt de 4} < \text{faille F} < \text{dépôt de 5} \]
(Le plissement précède nécessairement l'érosion qui arase les plis ; G recoupe 1 et 2 plissées et est tronqué par l'érosion, donc il se place entre le plissement et l'érosion.)
\item \textbf{Histoire géologique de la région (rédaction attendue, 5 à 8 lignes) :} Il y a très longtemps, des sédiments se sont déposés en mer et ont formé des couches de grès et de schiste où vivaient des trilobites. Ces couches ont ensuite été plissées par d'énormes forces tectoniques, puis traversées par du granit en fusion qui s'est solidifié en filon. L'ensemble a été soulevé et arasé par l'érosion pendant une très longue période. La mer est ensuite revenue : elle a déposé un conglomérat, puis un calcaire peuplé d'ammonites. Une faille a alors fracturé toutes ces couches, avant qu'une dernière couche de marne ne recouvre le tout et scelle la faille.
\end{enumerate}

\textbf{Barème indicatif (sur 12) :} trois principes énoncés 2 pts ; double démonstration G/F par rapport à l'érosion 2,5 pts ; position de F entre 4 et 5 1,5 pt ; argument paléontologique 2 pts ; chronologie complète 2,5 pts ; paragraphe rédigé 1,5 pt.
\end{corrige}

\begin{attention}[Exercice 11 — Points de vigilance]
\begin{itemize}
\item Le mot « \emph{tronqué} » est la clé de la question 2 : un filon tronqué par une surface d'érosion est forcément antérieur à cette érosion. Ne pas se contenter de « G recoupe 1 et 2 ».
\item Pour la faille, exiger le double argument : recoupement de 4 (postériorité) ET scellement par 5 (antériorité).
\item Dans la chronologie, ne pas oublier les \emph{événements} (plissement, érosion, intrusion) entre les dépôts de couches : une liste limitée aux couches est incomplète.
\item Le paragraphe final doit être rédigé en phrases, en langage courant, sans vocabulaire technique non expliqué : c'est un savoir-être de communication scientifique évalué au Probatoire.
\end{itemize}
\end{attention}

\newpage

\coursec{Fiche bilan}

\begin{popfigure}
\centering
\begin{center}\tcbox[colback=popblue,colframe=popblue,coltext=white,arc=9pt,boxsep=4pt,left=6pt,right=6pt]{\bfseries\small Chronologie des évènements géologiques}\end{center}
\vspace{-2pt}
\begin{tcbraster}[raster columns=3,raster equal height=rows,raster column skip=6pt,raster row skip=6pt,enhanced,blankest]
\begin{tcolorbox}[enhanced,colback=white,colframe=poporange,boxrule=0.9pt,arc=3pt,title={Stratigraphie},fonttitle=\bfseries\scriptsize,coltitle=white,colbacktitle=poporange,left=4pt,right=4pt,top=2pt,bottom=2pt,boxsep=1pt,toptitle=1.5pt,bottomtitle=1.5pt,halign=left]
\begin{itemize}[nosep,leftmargin=*,itemsep=1pt,font=\scriptsize]
\item couches sédimentaires = strates
\item strate = témoin du passé
\end{itemize}
\end{tcolorbox}
\begin{tcolorbox}[enhanced,colback=white,colframe=popgreen,boxrule=0.9pt,arc=3pt,title={Principes},fonttitle=\bfseries\scriptsize,coltitle=white,colbacktitle=popgreen,left=4pt,right=4pt,top=2pt,bottom=2pt,boxsep=1pt,toptitle=1.5pt,bottomtitle=1.5pt,halign=left]
\begin{itemize}[nosep,leftmargin=*,itemsep=1pt,font=\scriptsize]
\item superposition
\item continuité
\item recoupement
\end{itemize}
\end{tcolorbox}
\begin{tcolorbox}[enhanced,colback=white,colframe=poppurple,boxrule=0.9pt,arc=3pt,title={Discordances},fonttitle=\bfseries\scriptsize,coltitle=white,colbacktitle=poppurple,left=4pt,right=4pt,top=2pt,bottom=2pt,boxsep=1pt,toptitle=1.5pt,bottomtitle=1.5pt,halign=left]
\begin{itemize}[nosep,leftmargin=*,itemsep=1pt,font=\scriptsize]
\item angulaire
\item progressive
\item d'érosion
\end{itemize}
\end{tcolorbox}
\begin{tcolorbox}[enhanced,colback=white,colframe=poppink,boxrule=0.9pt,arc=3pt,title={Fossiles stratigraphiques},fonttitle=\bfseries\scriptsize,coltitle=white,colbacktitle=poppink,left=4pt,right=4pt,top=2pt,bottom=2pt,boxsep=1pt,toptitle=1.5pt,bottomtitle=1.5pt,halign=left]
\begin{itemize}[nosep,leftmargin=*,itemsep=1pt,font=\scriptsize]
\item large répartition
\item courte durée
\end{itemize}
\end{tcolorbox}
\begin{tcolorbox}[enhanced,colback=white,colframe=popgold,boxrule=0.9pt,arc=3pt,title={Évènements},fonttitle=\bfseries\scriptsize,coltitle=white,colbacktitle=popgold,left=4pt,right=4pt,top=2pt,bottom=2pt,boxsep=1pt,toptitle=1.5pt,bottomtitle=1.5pt,halign=left]
\begin{itemize}[nosep,leftmargin=*,itemsep=1pt,font=\scriptsize]
\item plis, failles
\item intrusions
\end{itemize}
\end{tcolorbox}
\begin{tcolorbox}[enhanced,colback=white,colframe=popblue!70,boxrule=0.9pt,arc=3pt,title={Méthode},fonttitle=\bfseries\scriptsize,coltitle=white,colbacktitle=popblue!70,left=4pt,right=4pt,top=2pt,bottom=2pt,boxsep=1pt,toptitle=1.5pt,bottomtitle=1.5pt,halign=left]
\begin{itemize}[nosep,leftmargin=*,itemsep=1pt,font=\scriptsize]
\item du plus ancien au plus récent
\end{itemize}
\end{tcolorbox}
\end{tcbraster}

\legende{Carte mentale de la leçon : les principes de la stratigraphie, les discordances et les fossiles stratigraphiques permettent d'établir la chronologie relative des évènements géologiques d'une région.}
\end{popfigure}

\begin{bilan}
\textbf{Définitions clés}
\begin{itemize}
  \item \cle{Stratigraphie} : science qui étudie l'organisation, la succession et la datation relative des couches (strates) sédimentaires.
  \item \cle{Strate} : couche de sédiments déposée lors d'un épisode de sédimentation ; son âge est celui de son dépôt.
  \item \cle{Discordance} : surface d'érosion séparant deux séries de couches d'âges différents ; elle témoigne d'une interruption de la sédimentation (émersion, tectonique, érosion).
  \item \cle{Fossile stratigraphique} : fossile à large répartition géographique et à courte durée d'existence, permettant de dater la couche qui le contient.
  \item \cle{Datation relative} : classement des évènements géologiques du plus ancien au plus récent, sans âge chiffré en années.
\end{itemize}

\textbf{Les trois principes de la stratigraphie (à connaître par c\oe ur)}
\begin{center}
\begin{tabularx}{\linewidth}{|l|X|}\hline
\rowcolor{popblueL} \textbf{Principe} & \textbf{Énoncé et conséquence} \\ \hline
Superposition & Dans une série non renversée, une couche est plus récente que celle qu'elle recouvre et plus ancienne que celle qui la surmonte. \\ \hline
Continuité & Une couche a le même âge sur toute son étendue ; des couches de même nature et de même niveau sont contemporaines. \\ \hline
Recoupement & Une structure (faille, intrusion, pli affectant des couches) est plus récente que les couches qu'elle recoupe ou déforme. \\ \hline
\end{tabularx}
\end{center}

\textbf{Les discordances stratigraphiques}
\begin{itemize}
  \item \cle{Discordance angulaire} : les couches supérieures horizontales reposent sur des couches inférieures plissées ou inclinées et érodées. Séquence : dépôt $\rightarrow$ plissement $\rightarrow$ érosion $\rightarrow$ nouvelle sédimentation.
  \item \cle{Discordance d'érosion} : surface d'érosion plane entre deux séries parallèles (lacune sédimentaire).
  \item \cle{Discordance progressive} : les strates supérieures s'amenuisent vers la bordure du bassin (transgression marine).
\end{itemize}

\textbf{Conditions d'un bon fossile stratigraphique}
\begin{itemize}
  \item large répartition géographique (espèce cosmopolite) ;
  \item courte durée d'existence (évolution rapide) ;
  \item abondance et facilité d'identification.
  \item Exemples : ammonites, trilobites, foraminifères.
\end{itemize}

\textbf{Méthode express : établir une chronologie relative}
\begin{enumerate}
  \item Repérer les couches et les numéroter de bas en haut (superposition).
  \item Repérer les structures qui recoupent : faille, dyke, pli $\rightarrow$ postérieures aux couches affectées.
  \item Repérer les discordances : elles marquent des étapes d'érosion entre deux phases de dépôt.
  \item Utiliser les fossiles stratigraphiques pour attribuer un âge relatif (étage) aux couches.
  \item Rédiger la chronologie du plus ancien au plus récent, en justifiant chaque étape par un principe.
\end{enumerate}
\end{bilan}

\begin{aretenir}[Checklist avant le Probatoire]
\begin{itemize}
  \item[$\square$] Je sais définir la stratigraphie, une strate et une discordance.
  \item[$\square$] Je connais par c\oe ur les trois principes : superposition, continuité, recoupement.
  \item[$\square$] Je sais que dans une série non renversée, la couche la plus ancienne est en bas.
  \item[$\square$] Je sais qu'une faille ou une intrusion est plus récente que les couches qu'elle recoupe.
  \item[$\square$] Je sais reconnaître une discordance angulaire sur une coupe géologique.
  \item[$\square$] Je sais décrire les étapes conduisant à une discordance angulaire (dépôt, plissement, érosion, nouveau dépôt).
  \item[$\square$] Je connais les trois critères d'un bon fossile stratigraphique.
  \item[$\square$] Je sais citer des exemples de fossiles stratigraphiques (ammonites, trilobites).
  \item[$\square$] Je sais distinguer datation relative et datation absolue.
  \item[$\square$] Je sais rédiger une chronologie du plus ancien au plus récent en justifiant chaque étape.
  \item[$\square$] Je vérifie toujours si la série est renversée avant d'appliquer le principe de superposition.
  \item[$\square$] Je légende correctement les coupes géologiques (couches, failles, surfaces de discordance).
\end{itemize}
\end{aretenir}

\findecours

\end{document}
